\documentclass[11pt,a4paper]{article}
\usepackage[margin=26mm,headheight=14pt]{geometry}
\usepackage[T1]{fontenc}
\usepackage{lmodern}
\usepackage{amsmath,amssymb,amsthm,mathtools}
\usepackage{mathrsfs}
\usepackage{microtype,longtable,booktabs,array,xcolor,graphicx}
\usepackage{enumitem}
\usepackage[colorlinks=true,linkcolor=blue!45!black,citecolor=blue!45!black,urlcolor=blue!45!black]{hyperref}
\usepackage{xurl}
\usepackage{bookmark}
\hypersetup{pdftitle={RESEARCH SA10: Stratified spaces and stratified topoi},pdfauthor={GPT-6 Astra (OpenAI, ChatGPT Pro mode)}}
\usepackage{fancyhdr}
\pagestyle{fancy}\fancyhf{}\fancyhead[L]{\small RESEARCH\_SA10}\fancyhead[R]{\small Checked research manuscript}\fancyfoot[C]{\thepage}
\setlength{\parskip}{4pt}
\setlength{\parindent}{0pt}
\setlist{itemsep=2pt,topsep=4pt}
\providecommand{\tightlist}{\setlength{\itemsep}{0pt}\setlength{\parskip}{0pt}}
\allowdisplaybreaks[2]
\newcommand{\Sh}{\operatorname{Sh}}
\newcommand{\Topos}{\mathbf{Topos}}
\newcommand{\Pro}{\operatorname{Pro}}
\newcommand{\Hom}{\operatorname{Hom}}
\newcommand{\colim}{\operatorname*{colim}}
\newcommand{\hocolim}{\operatorname*{hocolim}}
\newcommand{\holim}{\operatorname*{holim}}
\newcommand{\Exit}{\operatorname{Exit}}
\newcommand{\Map}{\operatorname{Map}}
\newtheorem{result}{Statement}[section]
\newtheorem{theorem}[result]{Theorem}
\newtheorem{proposition}[result]{Proposition}
\newtheorem{lemma}[result]{Lemma}
\newtheorem{corollary}[result]{Corollary}
\theoremstyle{definition}
\newtheorem{definition}[result]{Definition}
\newtheorem{remark}[result]{Remark}
\title{Grothendieck's stratified spaces and stratified topoi\\
\large Cotes 150--151: corrected gluing, homotopy, and control data}
\author{AI-produced by GPT-6 Astra\\\small OpenAI, ChatGPT Pro mode}
\date{19 November 2026}
\begin{document}
\maketitle
\begin{abstract}
We settle the stated questions about cotes 150--151 by distinguishing
ordinary topoi, pro-topoi of neighborhoods, and infinity-topoi.
The successor diagram and even the full manuscript pair/triple diagram
fail. A synchronized \v{C}ech construction gives a precise replacement,
with an exact deletion criterion for honest open diagrams. We derive
the corresponding higher shape and exit-path statements at their
proper truncation levels, give the complete conditional-meet
classification for stable-graph strata, and separate weak recursive
cone data from controlled manifold resolutions. Explicit counterexamples
settle the proposed finite-set presentations and the ambiguous TOP,
PL, and smooth margin assertions. Every result is stated with its
hypotheses, proof, source locators, and the scope of any finite certificate.
No priority is claimed.
\end{abstract}
\tableofcontents
\clearpage
\section*{Reading conventions and proof status}
\addcontentsline{toc}{section}{Reading conventions and proof status}
All scan locators refer to PDF page numbers; the archive pencil page is
one less. The attached second-pass TeX files are readings to be checked,
not authoritative mathematical premises. Each assertion below carries a
stable statement ID and a status. The accompanying results register,
check-event table, and progress ledger record the scope of verification.
Known results are used only at the cited locators in opened sources.
The symbols \emph{Proved}, \emph{Refuted}, and \emph{Proved under stated
hypotheses} refer to precisely the written statement, not to a broader
unwritten version of an item. We write in the first person plural to
describe the AI research process, and attribute none of this work to a
human author.
\part{The gluing of tubes}
\section{A.1. The corrected gluing statement}
\subsection{Conventions and the distinction between three categories}
We write $\Sh(Z)$ for the Grothendieck topos of sheaves of sets on a
topological space $Z$. Geometric morphisms point in the direction of maps of
spaces; their inverse-image functors point in the opposite direction.
Throughout this part, a colimit of topoi means a bicategorical colimit,
with invertible comparison 2-cells, in $\Topos$.
For a closed subset $Y\subset Z$, the notation
\[
 \mathcal N_Z(Y)=(\Sh(N))_{N\supset Y,\ N\text{ open in }Z}
\]
denotes a \emph{pro-topos}, with transition morphisms induced by inclusions
of smaller neighborhoods. It is not an object of $\Topos$. We use
$2\Pro(\Topos)$ with mapping categories
\[
 \Hom_{2\Pro}( (E_a)_a,(T_b)_b)
 \simeq \mathop{\mathrm{bilim}}_b
       \mathop{\mathrm{bicolim}}_a\Hom_{\Topos}(E_a,T_b).
\tag{A.1.1}\label{eq:prohom}
\]
The second colimit is filtered. Constant pro-objects are denoted by $cE$.
Finally, the infinity-topos of sheaves of spaces is denoted by
$\Sh_\infty(Z)$. Its shape is discussed in Part B; it must not be identified
with a functor preserving arbitrary bicategorical colimits of ordinary topoi.

\begin{result}[A1.1; Proved: the neighborhood-limit collapse]\label{A1.1}
Suppose that $Z$ is a regular Hausdorff space and $Y\subset Z$ is closed.
The inverse limit in $\Topos$ of the open subtopoi $\Sh(N)\subset\Sh(Z)$,
for $N$ ranging over all neighborhoods of $Y$, is $\Sh(Y)$. The analogous
limit of $\Sh(N\setminus Y)$ is the empty topos. Thus taking this limit
before gluing discards precisely the punctured-neighborhood information.
\end{result}
\begin{proof}
A geometric morphism $f:E\to\Sh(Z)$ factors through $\Sh(N)$ exactly when
$f^*N=1_E$, where $N$ is regarded as a subterminal object. Regularity gives
\[
 Z\setminus Y=\bigcup_{N\supset Y}\operatorname{int}(Z\setminus N).
\]
Indeed, for $z\notin Y$ choose an open neighborhood $V$ of $z$ whose closure
avoids $Y$, and put $N=Z\setminus\overline V$. Then $N$ contains $Y$ and
$V\subset\operatorname{int}(Z\setminus N)$.
If $f$ factors through every $N$, then
$f^*\operatorname{int}(Z\setminus N)=0_E$ for every $N$. Since $f^*$
preserves unions, $f^*(Z\setminus Y)=0_E$. This is exactly the condition for
factorization through the complementary closed subtopos $\Sh(Y)$.
Factorizations and their 2-morphisms are unique in the required fully
faithful sense, so this is the limit universal property. Apply the same
argument over $Z\setminus Y$ to the opens $N\setminus Y$; their excluded
interiors cover the whole space. The resulting factorization is through
the empty subtopos. This rechecks the substantive assertion of R8, under
the stated separation hypothesis; paracompactness alone without separation
is not what this proof uses.
\end{proof}

We fix the manuscript's axioms a)--d): the closed subsets $(X_i)_{i\in I}$
are locally finite, $i\leq j$ implies $X_i\subset X_j$,
$X=\bigcup_iX_i$, and
\[
 X_i\cap X_j=\bigcup_{k\leq i,j}X_k.
\tag{A.1.2}\label{eq:abcd}
\]
Here $I$ is a poset. For base change it is essential to use the order-based
closed subset
\[
 R_i^{\mathrm{ord}}=\bigcup_{k:\,i\nleq k}X_k,
 \qquad S_i=X_i\setminus\bigcup_{k<i}X_k.
\tag{A.1.3}\label{eq:Riord}
\]
When all $S_i$ are nonempty, this agrees with the support-based $R_i$ of
151/057. When empty strata or coincident closed subsets occur, it need not
agree; we do not silently replace one definition by the other.

\begin{result}[A1.2; Proved: open-tube descent, including infinitely many strata]\label{A1.2}
For every tuple $\alpha=(N_i)_{i\in I}$ of open neighborhoods $N_i\supset X_i$
in $X$, put
\[
 W_i^\alpha=N_i\setminus R_i^{\mathrm{ord}},\qquad
 C_q^\alpha=\Sh\left(\coprod_{(i_0,\ldots,i_q)\in I^{q+1}}
              W_{i_0}^\alpha\cap\cdots\cap W_{i_q}^\alpha\right).
\tag{A.1.4}\label{eq:tubecech}
\]
The $W_i^\alpha$ form an open cover of $X$. Their augmented Cech diagram
has the following properties.
\begin{enumerate}
\item For every fixed $\alpha$, its truncation in degrees $0,1,2$ has
      $\Sh(X)$ as its bicolimit in $\Topos$.
\item Regard each $C_q=(C_q^\alpha)_\alpha$ as a \emph{single synchronized
      pro-object}. Then
      \[
        \mathop{\mathrm{bicolim}}_{[q]\in\Delta_{\leq2}^{\mathrm{op}}}C_q
        \simeq c\Sh(X)
        \quad\text{in }2\Pro(\Topos).
      \tag{A.1.5}\label{eq:procech}
      \]
\item If $I$ is finite, one can equivalently use the individual nonempty
      intersections $W_A=(\Sh(\bigcap_{i\in A}W_i^\alpha))_\alpha$ for
      $1\leq |A|\leq3$, with all their face inclusions. For infinite $I$,
      the synchronized pro-coproducts in (\ref{eq:tubecech}) are part of the
      assertion; no interchange of an infinite product and a filtered
      colimit is asserted.
\end{enumerate}
No Thom--Mather, Whitney, or homotopy-link assumption is needed. Such
assumptions are needed for geometric models and homotopy replacements of
the intersections, not for this descent theorem.
\end{result}
\begin{proof}
For $x\in X$, the finite set $I_x=\{i:x\in X_i\}$ is nonempty and downward
directed by (\ref{eq:abcd}). It has a least element $i(x)$. If $x\in X_k$,
then $i(x)\leq k$; consequently $x\notin R_{i(x)}^{\mathrm{ord}}$ and
$x\in W_{i(x)}^\alpha$. Local finiteness makes every union in
(\ref{eq:Riord}) closed, so all these sets are open.

We give the full descent argument. For an arbitrary open cover $(W_i)$,
a sheaf $F$ on $X$ restricts to sheaves $F_i$ on $W_i$ and isomorphisms
$g_{ij}:F_i|_{W_i\cap W_j}\to F_j|_{W_i\cap W_j}$, with
$g_{ii}=1$ and $g_{jk}g_{ij}=g_{ik}$ on triple intersections. Conversely,
given these data, define, for each open $O\subset X$,
\[
 F(O)=\left\{(s_i)_i\in\prod_i F_i(O\cap W_i):
             g_{ij}(s_i|)=s_j|\text{ for all }i,j\right\}.
\]
The sheaf axioms for the $F_i$ prove both existence and uniqueness of
gluing in this formula. On $W_i$, the cocycle condition identifies this
sheaf with $F_i$. The identical construction for morphisms proves an
equivalence of categories, not merely a bijection on isomorphism classes.

To obtain the topos universal property, let $T$ be any topos and suppose
geometric morphisms $f_i:\Sh(W_i)\to T$ are identified on overlaps with
the stated cocycle. Apply sheaf descent to $f_i^*(A)$ for every $A\in T$,
and to their arrows, to obtain $f^*:T\to\Sh(X)$. This functor preserves
finite limits and all colimits because restriction to the open cover does,
and because those restrictions are jointly conservative. More explicitly,
a right adjoint exists: for $F\in\Sh(X)$, take in $T$ the equalizer of
the two arrows from $\prod_i f_{i*}(F|_{W_i})$ to
$\prod_{i,j}f_{ij*}(F|_{W_i\cap W_j})$. The adjunction follows from the
displayed formula for sheaf gluing. Natural transformations between the
$f_i^*$ glue in the same way. This proves the bicolimit universal property.
Degrees $0,1,2$ encode respectively objects, overlap isomorphisms, and
their cocycle; no further degree is used.

The poset of tuples $\alpha$, ordered by componentwise shrinking, is
cofiltered, even when $I$ is infinite. The category
$\Delta_{\leq2}^{\mathrm{op}}$ has finitely many objects and morphisms.
Filtered bicolimits of categories commute with finite bilimits: every
object, isomorphism, and equation in a finite cone is represented at one
common stage, and equality of two such finite systems is witnessed after
a common refinement. This proves essential surjectivity and full
faithfulness of the comparison. Applying this fact to
(\ref{eq:prohom}), and commuting limits with limits, gives, for any
pro-topos $(T_b)_b$,
\begin{align*}
 \mathop{\mathrm{bilim}}_{q\leq2}\Hom_{2\Pro}(C_q,(T_b)_b)
 &\simeq\mathop{\mathrm{bilim}}_b
     \mathop{\mathrm{bicolim}}_\alpha
     \mathop{\mathrm{bilim}}_{q\leq2}\Hom(C_q^\alpha,T_b)\\
 &\simeq\mathop{\mathrm{bilim}}_b\Hom(\Sh(X),T_b).
\end{align*}
The last equivalences are compatible with all refinements because their
augmentations are the identity on $X$. This is
(\ref{eq:procech}). For finite $I$, descent data on the finite intersection
poset of subsets of cardinality at most three are the same data already
written out, proving the final assertion.
\end{proof}

The empty case has $I=\varnothing$ and $X=\varnothing$ and gives the
initial topos; a one-stratum cover gives the identity diagram. Empty
indexed strata are allowed throughout. Local finiteness cannot simply
be dropped: take the real line with one closed singleton $X_x=\{x\}$
for every real $x$, indexed discretely. The cover, monotonicity and
intersection axioms still hold, but the resulting $W_x=\{x\}$ are not
open. The closed-piece colimit has the discrete topology and is not
the usual real line. Thus both the openness step here and the closed
quotient step of A4.1 have a concrete failure without that hypothesis.

The theorem supplies a precise correction of 151/069, but its intersections
are additional stones. They are not, in general, any of the
$V_{ij},V_{ij}^*,V_{i,k}^j$ of 151/057--061. When $X$ is not already one of
the closed $X_i$, the global tubes used in (\ref{eq:tubecech}) amount to
adjoining an ambient top index with closed subset $X$; its open stratum may
be empty. The following counterexample proves that this extra intersection
data cannot be replaced by simply allowing all pairs and triples in the
manuscript's formulas.

\subsection{The successor and all-pair counterexamples}
\begin{result}[A1.3; Refuted: 151/069, and Lead L1's all-pair repair]\label{A1.3}
The successor-pair diagram of 151/069 does not have $\Sh(X)$ as its
colimit. Even the diagram of \emph{all} manuscript pair and triple stones,
with \emph{all} their canonical inclusions over $X$, does not suffice.
Both assertions fail for the Whitney and conically smooth stratification
\[
 X=\mathbb C,\qquad X_0=\{0\},\quad X_1=\mathbb R,\quad X_2=\mathbb C.
\]
Disconnected open strata are allowed, as in the question. The all-pair
and all-triple failure also has a connected-stratum example in
$\{0\}\subset\mathbb R^2\subset\mathbb R^4$, described below.
\end{result}
\begin{proof}
The five connected cells are $o=\{0\}$, the positive and negative rays
$p,n$, and the upper and lower half-planes $u,d$. Their incidence poset
has relations
\[
 o<p<u,\quad o<n<u,\quad o<p<d,\quad o<n<d.
\]
Every covariant set-valued functor on this poset yields a sheaf on the
plane with those cell stalks and generization maps: take its sheaf on the
finite Alexandrov poset and pull it back along the continuous cell map.
The same construction applies to the punctured plane and the real line.
This explicitly realizes all sheaves used below.

For the successor version assign a singleton to $o,p,n,d$ and the set
$\{0,1\}$ to $u$. On $V_{01}$ use singleton stalks and their unique maps.
On $V_{12}$ use the sheaf on the punctured plane whose maps $p\to u$ and
$n\to u$ have values $0$ and $1$, respectively. Its other nonidentity
maps have singleton targets. These data agree on all successor stones.
Any sheaf on the plane restricting to them would have a section
representing the unique stalk element at $o$. After shrinking to a disk,
its restriction to the connected upper half-disk must have a single value.
Restrictions near the positive and negative rays force that value to be
both $0$ and $1$. This is impossible. This independently verifies the
counterexample recorded as SA10-C3-14 and SA10-S3-K3.

For the stronger counterexample retain the same stalk sets, but put value
$1$ on \emph{both} maps $p,n\to u$ in $V_{12}$. On $V_{02}$ put value
$0$ on both maps, and also on $o\to u$. All remaining maps to singleton
stalks are unique. We list the nine distinct stones; writing
$\operatorname{germ}_Y Z$ abbreviates their pro-neighborhood systems:
\[
\begin{array}{c|c@{\qquad}c|c}
S_0&\{0\}&S_1&\mathbb R^*\\
S_2&\mathbb C\setminus\mathbb R&V_{01}&\operatorname{germ}_0\mathbb R\\
V_{01}^*&(\operatorname{germ}_0\mathbb R)\setminus\{0\}
&V_{12}&(\operatorname{germ}_{\mathbb R}\mathbb C)\setminus\{0\}\\
V_{12}^*&(\operatorname{germ}_{\mathbb R}\mathbb C)\setminus\mathbb R
&V_{02}&\operatorname{germ}_0\mathbb C\\
V_{02}^*&(\operatorname{germ}_0\mathbb C)\setminus\mathbb R&&
\end{array}
\]
Besides the standard three arrows for each pair, the only new generating
inclusions are $V_{01}\to V_{02}$ and
$V_{02}^*\to V_{12}^*$. Composites include
$V_{01}^*\to V_{12}$ and $V_{01}^*\to V_{02}$. To check completeness of
this list, a germ containing $0$ cannot map to a stone deleting $0$;
real and nonreal supports cannot contain one another; and a germ along
the whole real line cannot be contained in an arbitrarily small disk at
$0$. The remaining possibilities are exactly those listed. In particular,
$V_{12}\not\subset V_{02}$. The inclusion
$V_{02}^*\subset V_{12}^*$ does hold, because every neighborhood of the
whole real line contains a disk about $0$.

All listed restrictions agree identically on stalks and on every map they
see. For the unique strict triple $0<1<2$, formula (11) of 151/060 gives
$S_{0,2}=S_{1,2}=\mathbb R$, hence
$V_{0,2}^{0}=V_{0,2}^{1}=V_{02}^*$ and $V_{0,2}^{2}=V_{02}$.
Triples with repeated indices reproduce pair stones. They impose no new
condition. But the actual overlap of $V_{02}$ and $V_{12}$ contains the
punctured disk germ, where the two $p\to u$ maps differ. If a global sheaf
existed, its identifications on $S_1,S_2$ and on these overlaps would make
the two maps identical. The discrepancy survives every shrinking, so no
global sheaf exists.

We checked the primary scan 151/060 as well as its transcription. In
formula (10), the overwritten first index of $R$ is potentially read as
$i$ or $j$; the following sentence specifies $R_{i,k}$. Neither reading
repairs this counterexample. For the only interior index $j=1$, the
deletion of $S_{1,2}=\mathbb R$ absorbs both $R_{0,2}=\varnothing$ and
$R_{1,2}=\{0\}$. At the upper endpoint, reading $R_{j,k}$ instead
merely replaces the repeated $V_{02}$ by $V_{02}^*$, both already in
the pair diagram. The doubtful letter $S$ versus $\Sigma$ denotes the
set explicitly defined in (11). The later doubtful bound in the
pseudo-partition is not used: (11) and its adjacent $j=k$ note specify
the strict upper bound. If that bound were changed to $l\leq k$, its
extra $X_k$ would make the corresponding triple stone empty; this
weaker diagram still cannot detect the missing overlap.

This sheaf obstruction tests a genuine topos colimit universal property.
The object-classifying topos $\mathbf{Set}^{\mathrm{FinSet}}$ has
$\Hom(E,\mathbf{Set}^{\mathrm{FinSet}})\simeq E$: a left-exact functor
$\mathrm{FinSet}^{\mathrm{op}}\to E$ is $K\mapsto A^K$ for a unique
object $A$, and the presheaf-topos universal property gives the claimed
equivalence. Thus the compatible sheaves above are a compatible cocone of
geometric morphisms to this classifier, with no extension from $\Sh(X)$.
The argument applies unchanged in the pro-bicategory, using representatives
and their germs. For a version with every open stratum connected, replace
the chain by $\{0\}\subset\mathbb R^2\subset\mathbb R^4$. On the closed
plane use the singleton sheaf, and on its complement use the constant
two-point sheaf. The two constant choices $0,1$ define two global
closed--open recollement sheaves $F^0,F^1$ by B1.1. Put $F^0$ on $V_{02}$
and $F^1$ on $V_{12}$. Their restrictions to all named strata and to
$V_{01}$ agree, whereas their maps from the middle stratum to the top
one disagree on every punctured-origin overlap. The list of possible
germ inclusions and the equality $V_{02}^{0}=V_{02}^{1}=V_{02}^*$ are
unchanged. The strata are connected smooth manifolds; the nested linear
subspaces satisfy Whitney's conditions and have their product conical
control data. Thus connectedness does not repair the all-pair claim.
The program \texttt{check\_A\_finite.py} verifies the
finite functor equations and all 14 nonidentity inclusions in the
all-pair example. The analytic realization just given verifies the
topological hypotheses, which the finite program alone could not verify.
\end{proof}

\begin{result}[A1.4; Proved: precise universal minimality]\label{A1.4}
For ordinary-topos open-cover descent, vertices, pair overlaps, and triple
cocycles are the smallest universal Cech arities. In the subset-indexed
presentation of A1.2(3), each vertex, pair, or triple intersection object
can be indispensable for an appropriate finite open cover;
intersections of four or more members are unnecessary. There is no
claim that every slot of every particular cover is indispensable.
Ordered repetitions and degeneracies in the full Cech diagram are not
being counted as independently necessary pieces.
In binary gluing, the three-piece span
$\Sh(N)\leftarrow\Sh(N\cap U)\to\Sh(U)$ is universally indispensable;
an extra soul $\Sh(Y)\to\Sh(N)$ with no other incidence is redundant.
\end{result}
\begin{proof}
The sufficiency is the explicit descent proof of A1.2. To see necessity
of a vertex slot, let a point belong only to that member of a cover;
deleting it removes all information at the point. For a pair slot, let
only those two cover members contain a point. Without their intersection,
one may assign stalk sets of cardinalities one and two at the two copies,
which cannot come from a single sheaf. For a triple slot, let exactly the
three members contain a point and assign a two-element set to each.
Take two pair identifications to be the identity and the third to be the
transposition. Every inverse rule holds, but the triple cocycle fails.
Other cover members may be taken disjoint from that point. These are
finite counterexamples; the transposition calculation is in the program.
The binary examples are the same missing-overlap example and the missing
vertex examples, realized, if desired, by a tube and its complement in
the real line. An object with its sole outgoing arrow into an already
present object contributes no additional bicolimit data: its cocone map
is forced by composition, with a contractible choice of comparison.
That proves the assertion about the soul and also explains why absolute
minimality of the manuscript's named list cannot be demanded.
\end{proof}

\subsection{Comparison with unzipping and SGA 4}
\begin{result}[A1.5; Proved: exact criterion for deleting open pieces]\label{A1.5}
Let $X$ be a $T_1$ space, $C$ a small category, and
$D:C\to\operatorname{Open}(X)$ a diagram of open subsets and inclusions.
For $x\in X$, let $C_x$ be the full subcategory on the objects $c$ with
$x\in D(c)$. The canonical cocone has $\Sh(X)$ as its bicolimit if and
only if the groupoid obtained by inverting every arrow of $C_x$ is
equivalent to the terminal groupoid, for every $x$. Equivalently, each
$|N(C_x)|$ is nonempty, connected, and simply connected.
This gives an exact test for any proposed deletion of objects or arrows
from a specified honest-open gluing diagram. An inclusion-minimal
subdiagram is precisely one satisfying this test for which every further
permitted deletion fails the test at some point.
\end{result}
\begin{proof}
A sheaf descent datum on $D$ is a family of sheaves $F_c$ with
isomorphisms $F_d|_{D(c)}\simeq F_c$ for every arrow $c\to d$, satisfying
the functorial identities. At $x$ its stalks therefore define a functor
from the groupoidification of $C_x$ to sets. If that groupoid is terminal,
all these stalks identify uniquely, with no possible monodromy.
Take the disjoint union of the etale spaces of the $F_c$, and quotient
by the identifications on the open subspaces specified by arrows. The
quotient map is open: the saturation of an open set is the union of its
images under finite strings of these open-subspace homeomorphisms.
At a point $x\in D(c)$, connectedness gives surjectivity from the
$c$-stalk to the quotient fibre, and trivial groupoid monodromy gives
injectivity. Thus each etale space maps homeomorphically onto the
quotient over $D(c)$. The quotient is an etale space over $X$, since
the $D(c)$ cover $X$. It gives a sheaf restricting to the specified
datum. Morphisms descend by the same quotient construction, proving an
equivalence of categories. Applying the argument to inverse-image
functors and their natural transformations, as in A1.2, proves the
topos bicolimit property.

Conversely, if $C_x$ is empty, two distinct skyscraper sheaves at the
closed point $x$ restrict identically to the entire diagram. If $C_x$
is disconnected, choose stalk sets of different cardinalities on two
components. If it is connected with nontrivial groupoid isotropy, choose
the regular action of one of its nontrivial automorphism groups and
transport it to all its objects. These choices are nonconstant descent
data in sets on $C_x$. They extend to sheaf descent data on $D$: for
$x\in D(c)$ use the skyscraper with the specified stalk, and for
$x\notin D(c)$ use the terminal sheaf. Restrictions to opens missing
$x$ are terminal, so all required identifications and equations hold.
A sheaf on $X$ would induce the constant stalk datum on $C_x$, contrary
to the choice. The object classifier in A1.3 converts this obstruction
to failure of the topos universal property. Finally, the nerve of a
category has fundamental groupoid equal to its groupoidification:
edges give arrows and formal inverses, and its triangles impose exactly
the composition relations. Higher simplices do not change that
presentation. This proves the stated equivalent topological criterion.
\end{proof}

The geometric unzipping theorem of Ayala--Francis--Tanaka
\cite[Theorem 7.3.8, Corollary 7.3.7 and Proposition 8.2.3]{AFT}
provides, under conical smoothness and their finite-dimensional hypotheses,
a resolution with corners, a compact link bundle along a deepest stratum,
and its cone neighborhood. Taking that neighborhood and its open complement
produces exactly a binary cover to which A1.2 applies. Repeating the
geometric construction retains the boundary incidences and their
overlaps; it does not identify them with the manuscript's all-pair list.
The plane example therefore refutes the asserted equivalence between
these two procedures in Lead L1.

SGA 4, Expose IV, Theorem 9.5.4 \cite{SGA4}, reconstructs a topos from an
open subtopos, its closed complement, and the functor $i^*j_*$. Its objects
are triples $(A,B,A\to i^*j_*B)$. The arrow is not required to be invertible.
Our open-cover descent uses invertible overlap identifications; the
noninvertible generization information is already present inside each
mixed tube. A1.3 shows that omitting an actual overlap fails to compare
exactly this information. The two formulations are consistent once the
neighborhood system, its overlap, and the gluing functor are retained.

\section{A.2. Base change of tubes and locality}
\begin{result}[A2.1; Proved: exact neighborhood criterion]\label{A2.1}
Let $f:Z'\to Z$ be continuous and $Y\subset Z$ closed. The canonical
comparison from the neighborhood pro-object of $f^{-1}Y$ in $Z'$ to the
pullback of $\mathcal N_Z(Y)$ is an isomorphism over $Z'$ if the family
$\{f^{-1}N:N\supset Y\text{ open}\}$ is cofinal among neighborhoods of
$f^{-1}Y$. Equivalently, for every closed $F'\subset Z'$ disjoint from
$f^{-1}Y$, there is an open $N\supset Y$ with $f^{-1}N\cap F'=\varnothing$.
This holds for every closed continuous map $f$, in particular for proper
maps between locally compact Hausdorff spaces. No local-immersion
hypothesis is necessary.
\end{result}
\begin{proof}
The two neighborhood systems have the canonical inclusion morphisms over
$Z'$. Cofinality is precisely the assertion that each neighborhood in the
first system contains a pullback neighborhood, giving inverse morphisms
of these pro-subspaces. Replacing an open neighborhood $O'$ by its closed
complement gives the equivalent condition. If $f$ is closed, then
$f(F')$ is closed and disjoint from $Y$, so $N=Z\setminus f(F')$ works.
All comparisons are induced by actual inclusions, hence the same
isomorphisms hold for their pro-topoi and after deleting compatible
closed subsets. Conversely, an inverse in the category of pro-subspaces
over $Z'$ must, for each $O'$, factor some pullback neighborhood into
$O'$; this is cofinality. This is an exact criterion for this specified
over-space comparison, not a claim about abstract isomorphisms of
unlabelled pro-objects.
\end{proof}

\begin{result}[A2.2; Refuted: unrestricted open base change and support-based $R$]\label{A2.2}
Open immersions do not generally satisfy A2.1. Moreover, formula (4) of
151/071 for the support-based $R_{ij}$ fails under arbitrary inverse
images, even for a closed immersion. The order-based definition
(\ref{eq:Riord}) does commute with every inverse image.
\end{result}
\begin{proof}
Take $Z=\mathbb R$, $Y=\{0\}$ and $Z'=\mathbb R\setminus\{0\}$.
The induced center is empty, so its neighborhood pro-object has the empty
set as a cofinal stage. The pulled-back neighborhoods are punctured
interval germs, and none is empty. They cannot factor into the empty
stage over $Z'$, so the comparison is not an isomorphism. Their pro-homotopy
types even have two components instead of zero.

For the second assertion use $X_0=\{0\}\subset X_1=\mathbb R$ and the
closed immersion $f:\{0\}\hookrightarrow\mathbb R$. Originally $R_1=X_0$.
After pullback $X'_0=X'_1=\{0\}$, so no $X'_k$ fails to contain $X'_1$;
the recomputed support-based $R'_1$ is empty. But $f^{-1}R_1=\{0\}$.
The discrepancy is removed if the index condition $i\nleq k$ is held
fixed: inverse image commutes with the union in (\ref{eq:Riord}), and
likewise with the fixed-index unions defining $S_{j,k}$.

For completeness, the order-based identity is always
$R_i^{\mathrm{ord}}\cap X_i=\bigcup_{h<i}X_h$. Indeed, intersect each
$X_k$ with $X_i$ and use (d); a lower index $h\leq i,k$ cannot equal
$i$ when $i\nleq k$. Conversely every $h<i$ occurs in the defining
union. If $S_i$ is nonempty, the argument of A4.4 shows
$X_i\subset X_k$ exactly when $i\leq k$, recovering the manuscript's
support-based identity (4) under this hypothesis. Without it the
coincident-piece example above refutes (4). The separate parenthesis
after (12) of 151/060 is false even with nonempty strata: in the
three-point chain $X_0=\{0\}\subset X_1=\{0,1\}\subset
X_2=\{0,1,2\}$, take $(i,j,k)=(0,1,2)$. Then
$R_{1,2}\cap S_{1,2}=\{0\}$, whereas $\dot X_0=\varnothing$.
The finite identities are checked in the attached program. This
rechecks the substantive R3 qualification and its separately false
parenthesis without relying on the earlier check labels.
\end{proof}

The supplied second-pass transcription of 151/071 explicitly inserts
\emph{mais pas} before the reference to an open immersion. Thus the prompt
and \texttt{study/STUDY.md}, line 221, misstate this sentence. The
mathematical negation above is independent of that reading: it proves the
claim false if the parenthesis were read positively, and confirms the
negative reading if it is retained.

\begin{result}[A2.3; Proved: the corrected locality statement]\label{A2.3}
Descent for a specified open-cover diagram is local on $X$: restrictions
of that cover to an open cover of $X$ verify the same gluing equivalence,
and these local equivalences glue. The argument proposed and retracted
in 151/072 is invalid when it replaces the pullback of a closed-center
pro-neighborhood by the independently recomputed neighborhood system
on an open subset.
\end{result}
\begin{proof}
For honest opens, $(W_i\cap W_j)\cap O=(W_i\cap O)\cap(W_j\cap O)$,
with identical equalities for triples. The explicit sheaf construction
in A1.2 commutes with restricting its argument $O$ to any smaller open.
Therefore it is a local inverse to restriction, and inverses agree on
overlaps by uniqueness of sheaf gluing. Applying this level by level
proves locality for a specified pulled-back synchronized system.
A2.2 shows that recomputation can change even an empty center into a
nonempty punctured pro-object. Thus the missing comparison in 151/072
cannot be inferred from open base change. This refutes the inference,
while the valid local descent theorem remains as strong as stated.
\end{proof}

\section{A.3. Differences of cribles}
\begin{result}[A3.1; Proved and Refuted: exact closed restriction, false open homotopy comparison]\label{A3.1}
Let $I''\subset I'\subset I$ be lower ideals, and put
$Z=X_{I''}$, $X'=X_{I'}$, $U=X'\setminus Z$. The induced strata of $U$
are exactly the $S_i$ with $i\in I'\setminus I''$. For retained $i,j$,
passing from $X$ to the closed union $X'$ preserves the pair stones
defined inside $X_j$ (with fixed order-based deletions, or nonempty strata).
Passing to $U$ need not preserve even their number of connected
components or their homotopy types, including in the conically smooth
Whitney example of A1.3.
\end{result}
\begin{proof}
If $x\in S_i$ and $x\in X_k$ for $k\in I''$, then $i\leq k$ forces
$i\in I''$. This proves the assertion about remaining strata and the
same argument gives the induced lower unions. Since $I'$ contains every
index below each of its elements, the ambient $X_j$ and all its relevant
closed deletions are unchanged on restriction to $X'$. For support-based
deletions under nonempty strata, $X_i\subset X_k$ is equivalent to
$i\leq k$, so this is the same argument.

Now take $I'=\{0,1,2\}$ and $I''=\{0\}$ in A1.3. Originally
\[
 V_{12}=(\operatorname{germ}_{\mathbb R}\mathbb C)\setminus\{0\}.
\]
A cofinal family of neighborhoods of $\mathbb R$ consists of strips
$\{|y|<\varepsilon(x)\}$ with $\varepsilon:\mathbb R\to(0,\infty)$
continuous. Such strips are homeomorphic to $\mathbb R\times(-1,1)$,
and deleting the origin gives homotopy type $S^1$. The inclusions between
shrunk strips preserve the once-around-origin generator and are homotopy
equivalences, so the pro-homotopy type is constant $S^1$.
In $U=\mathbb C^*$ the induced center is $\mathbb R^*$.
There is a cofinal family of two disjoint variable-width strips over the
positive and negative rays, each of width less than $|x|/2$. Each
component contracts onto its ray. The induced tube therefore has the
constant homotopy type of two points. Cofinality of both strip families
follows by choosing a positive continuous minorant of the distance to
the complement, using paracompactness of the real line or its open rays.

Deleting the real center from the original strip gives two contractible
components, upper and lower. Deleting the center from the induced pair
of strips gives four contractible components. Hence the maps on deleted
tubes induce the nonbijective component map $\{1,2,3,4\}\to\{1,2\}$.
Both comparisons of 151/075 fail to be homotopy equivalences. This is a
cofinal-system calculation, not an artifact of a badly chosen representative.
\end{proof}

\begin{result}[A3.2; Proved: comparison of constructible sheaves]\label{A3.2}
For an open union of strata $j:U\hookrightarrow X'$, restriction
$j^*: \Sh_{\mathrm{c}}(X')\to\Sh_{\mathrm{c}}(U)$ is essentially
surjective and has the fully faithful extension-by-empty left adjoint
$j_!$. It is an equivalence exactly when $Z=X'\setminus U$ is empty.
Its full reconstruction is the recollement with $\Sh_{\mathrm{c}}(Z)$
and the nearby-section functor of B.1, with its constructibility conditions.
\end{result}
\begin{proof}
For sheaves of sets, $j_!$ exists for every open immersion. Its stalks
are the original stalks on $U$ and empty on $Z$, so it is constructible
for this stratification. The identity $j^*j_!=1$ proves essential
surjectivity and full faithfulness of $j_!$. If $Z\ne\varnothing$,
let $i:Z\hookrightarrow X'$ and compare the two sheaves
$i_*\underline{\{0\}}$ and $i_*\underline{\{0,1\}}$. They are
constructible: on $Z$ they are locally constant, and on $U$ both have
singleton stalks. They are not isomorphic at a point of $Z$, but their
restrictions to $U$ are isomorphic. Thus restriction cannot be an
equivalence. If $Z=\varnothing$ it is the identity. The recollement
description follows from the explicit inverse in B.1 or SGA 4 IV 9.5.4.
\end{proof}

\section{A.4. Closed diagrams, infima, and stable graphs}
\begin{result}[A4.1; Proved: the reconstruction already valid in R2]\label{A4.1}
Under a)--d), the strata $S_i$ partition $X$, and the colimit of the
closed-subset diagram $i\mapsto X_i$, equivalently the coequalizer of
its flag diagram, is $X$ as a topological space. Empty strata may be
discarded without destroying a)--d).
\end{result}
\begin{proof}
The least element $i(x)$ constructed in A1.2 is precisely the unique
index with $x\in S_{i(x)}$. Any two representatives of $x$ in $X_i$ and
$X_j$ are identified through $X_k$ for some $k\leq i,j$ containing $x$,
by (d). Thus the underlying set coequalizer is $X$. The map
$\coprod_i X_i\to X$ is closed: a closed subset of the coproduct is a
family of closed subsets of the closed $X_i$, and their union is closed
by local finiteness. It is surjective, hence a quotient map, proving the
topological assertion. Finally,
$X_k=\bigcup_{i\leq k,\,S_i\ne\varnothing}X_i$, because $i(x)\leq k$
for every $x\in X_k$. This identity verifies the cover and intersection
axioms after discarding empty strata; the other axioms are inherited.
This rechecks SA10-C3-01, -02, and -03.
The primary scan 151/070 really writes a strict lower bound in its
repeated intersection axiom. Taken literally, this forces $X$ empty
under local finiteness: for any supposed point choose its least
containing index $i(x)$ and apply that axiom to $i=j=i(x)$. Its right
side has no member containing the point. Thus all nonempty applications
here explicitly use the nonstrict axiom of 151/052, not the slip on
151/070. This verifies the mathematical correction R6.
\end{proof}

\begin{result}[A4.2; Refuted: the corollary of 151/052--053]\label{A4.2}
A finite meet-semilattice diagram of closed embeddings need not embed
its members in its colimit, even if its associated ordered space satisfies
(23a)--(23b) of 151/053.
\end{result}
\begin{proof}
Let $I$ be the seven proper subsets of $\{1,2,3\}$, including the empty
set, ordered by inclusion. Its meet is intersection. Assign the singleton
$\{*\}$ to $\varnothing$ and the discrete three-point space
$\{*,0,1\}$ to each of the other six indices. Bottom arrows send $*$ to
$*$. Every proper transition map between the other indices is the identity except
$X_{\{1\}}\to X_{\{1,2\}}$, which swaps $0,1$ and fixes $*$. Every
composable pair of strict arrows starts at the bottom and fixes $*$, so this is a
functor, and all maps are closed embeddings. The disjoint union is a
finite discrete ordered space: comparability is given by the graphs of
the transition maps. Reflexivity, antisymmetry, and transitivity hold,
and its order relation is closed. Thus (23a)--(23b) hold.
The six non-bottom indices form a hexagon. Going once around it transports
$0$ to $1$, since exactly one transition is a transposition. Hence the
nineteen points of the disjoint union have two equivalence classes in the
colimit, one for $*$ and one for $0,1$. No three-point member embeds.
Deleting the $*$ orbit gives the variant with an empty bottom, twelve
points and a singleton colimit. Both complete finite checks are in
\texttt{check\_A\_finite.py}; the first version avoids any convention
about whether an indexed partition may contain an empty part.
\end{proof}

\begin{result}[A4.3; Proved: exact and geometric repairs of the corollary]\label{A4.3}
For a finite diagram of closed embeddings indexed by a finite poset,
all maps to its topological colimit are closed embeddings if and only if
no zigzag of transition identifications identifies two distinct points
in one member. If $I$ is a finite meet-semilattice, the stronger condition
\[
 X_{i\wedge j}\xrightarrow{\sim}X_i\times_{X_k}X_j
 \quad(i,j\leq k)
\tag{A.4.1}\label{eq:cartesianmeet}
\]
is sufficient, and then the intersection of the images of $X_i,X_j$ in
the colimit is exactly the image of $X_{i\wedge j}$.
\end{result}
\begin{proof}
The zigzag condition is exactly injectivity of each set map to the
colimit. Assuming it, any zigzag joining two members can have a repeated
index removed, because a revisit to an index must revisit the same
point. There are only finitely many simple index paths. Transport of a
closed subset along each step is either a closed image under a closed
embedding or its inverse image. Thus its saturation meets each member
in a finite union of closed subsets. The quotient topology shows that
each member's map is closed, and the same argument applies to every
closed subset of it. An injective continuous closed map is an embedding.
The converse follows immediately from injectivity.

Under (\ref{eq:cartesianmeet}), follow a zigzag from a point of $X_i$.
Inductively maintain a common lower index and a point mapping both to
the starting point and to the current representative. An upward step
preserves this datum. For a downward step to $X_j$, use the pullback
square with the old lower index and $j$ inside their common upper
index, replacing the lower index by their meet. The point lifts
uniquely. At the end, representatives in $X_i$ and $X_j$ therefore come
from a lower index, and hence from $i\wedge j$. Taking $i=j$ proves the
zigzag condition. Taking arbitrary $i,j$ gives the claimed intersection
identity, including its subspace topology since all maps are closed
embeddings. This proves the repair without assuming that a poset
meet alone forces a pullback square.
\end{proof}

\begin{result}[A4.4; Proved: the class defined by (24)]\label{A4.4}
Under a)--d) with nonempty strata, the strengthened condition (24) of
151/054 is equivalent to the family of closed subsets $\{X_i\}$ being
closed under binary intersection, with the intersection represented by
one index. The conditional version, requiring infima only for pairs
with a common lower bound, is equivalent to this assertion for nonempty
intersections. The literal all-pairs version additionally needs a
bottom index representing the empty intersection whenever disjoint
closed strata occur.
\end{result}
\begin{proof}
If $X_i\subset X_j$ and $S_i\ne\varnothing$, choose $x\in S_i$.
Then $i=i(x)\leq j$. Thus $i\mapsto X_i$ is an order embedding.
An index representing $X_i\cap X_j$ is consequently precisely the
greatest lower bound. Conversely, if the infimum exists, (d) says that
the union of all lower closed subsets is its closed subset, giving the
intersection equality. Since every retained $X_k$ contains a nonempty
stratum, an intersection is nonempty exactly when a common lower
index exists. A formal empty bottom handles the remaining pairs.
Every finite meet-semilattice is realized in this general topological
class by its upper-Alexandrov space, with $X_i=\downarrow i$; so the
criterion describes exactly a closure-system condition and imposes
no hidden smoothness assertion.
\end{proof}

\begin{result}[A4.5; Proved: stable-graph classification and L4]\label{A4.5}
Let $2g-2+n>0$, and order isomorphism classes of connected stable graphs
of genus $g$ with $n$ individually labelled legs by reverse contraction:
$H\leq G$ if $H$ contracts to $G$. Every pair having a common lower
bound has an infimum exactly for
\[
 g=0\ (n\geq3),\qquad g=1\ (1\leq n\leq3),\qquad (g,n)=(2,0).
\tag{A.4.2}\label{eq:smallgn}
\]
Consequently the conditional interpretation of Lead L4 and the quoted
EXIT\_TOWER A.2 is correct. Literal (24) requires the empty-bottom
qualification of A4.4. This statement concerns graph types, not fixed
contractions or the incidence category.
\end{result}
\begin{proof}
A stable graph has vertex genera $g_v\geq0$, labelled legs, loops and
multiple edges, with $2g_v-2+\operatorname{val}(v)>0$ and
$g=\sum_vg_v+|E|-|V|+1$. For genus zero it is a tree, and its edges
determine stable bipartitions of the labelled legs. Distinct edges give
distinct bipartitions. A common refinement supplies a tree containing
the two sets of bipartitions; contract precisely the edges outside their
union. The resulting tree is independent of the common refinement:
root at one labelled leg, orient each split by the side avoiding that
leg, and reconstruct vertices from the inclusion hierarchy of these
nested or disjoint subsets. This reconstruction is unique. Every other
common refinement contracts to it, so it is the infimum.

For the remaining positive cases we give a complete finite certificate,
with a proof of exhaustive generation. The program
\texttt{stable\_graph\_meets.py} starts with the corolla, then applies
every inverse single-edge contraction: a loop lowering a positive
vertex genus by one, or a split of a vertex into two joined by an edge,
with all partitions of its genus and incident half-edges and legs.
It rejects precisely the unstable splits and identifies graphs by all
vertex permutations preserving the data. Contracting any edge of a
stable graph preserves stability; iterating contracts every graph to
the corolla. Thus induction on edge number proves that the enumeration
is exhaustive. The bound $|E|\leq3g-3+n$ follows by summing the vertex
stability inequalities, so the process terminates. The program computes
all contraction closures, all pairs, and all minimal common refinements.
For $(g,n)=(1,1),(1,2),(1,3),(2,0)$ the graph counts are respectively
$2,5,23,7$; all $1,10,253,21$ unordered pairs were checked, and none
has two minimal common refinements. The graph lists and contraction
closures are attached as finite certificates.

Here are uniform counterexamples in every other case; they also explain
the failures independently of enumeration. If $g\geq3$, let $G$ be the
one-edge separating graph with vertex genera $1,g-1$, and let $H$ be
the one-loop graph with one vertex of genus $g-1$. Put all markings,
if any, on the second vertex of $G$. One two-edge common refinement
has a loop on a genus-zero vertex joined to a genus-$(g-1)$ vertex;
the other has vertex genera $1,g-2$, with the loop on the second
vertex. Put all markings on the second vertex in both. Loop
contraction gives $G$, and bridge contraction gives $H$. Both are
stable and are nonisomorphic, since one has a genus-zero vertex and
the other does not. Their only one-edge contractions are $G,H$, which
are different. They are therefore distinct minimal common refinements.

If $g=2,n\geq1$, take $G$ with two genus-one vertices joined by an
edge, all legs on one vertex, and $H$ with one genus-one vertex,
all legs, and one loop. The refinements consist of a genus-zero loop
vertex joined to a genus-one vertex; in one refinement all legs lie
on the loop vertex, and in the other they lie on the other vertex.
They are stable, nonisomorphic, and minimal by the same two-edge
argument.

Finally let $g=1,n\geq4$. Let $G$ be a two-edge cycle with leg
partition $\{1\}\mid\{2,\ldots,n\}$, and $H$ the analogous cycle
with partition $\{2\}\mid\{1,3,\ldots,n\}$. A common refinement is
the triangle with vertex leg sets $\{1\},\{2\},\{3,\ldots,n\}$.
Another is the square whose opposite vertices carry $\{1\},\{2\}$,
and whose other vertices carry $\{3\},\{4,\ldots,n\}$. Every
vertex is rational and stable. Contracting suitable edges gives both
$G$ and $H$. In any common refinement obtained from the square,
the vertices carrying legs $1$ and $2$ must each remain separate
from all other marked vertices. Every square edge touches one of
these vertices, so no edge can be contracted while retaining both
requirements. The square is minimal. The triangle is also minimal,
since a two-edge common refinement would have to equal both $G$ and
$H$. These two nonisomorphic minimal refinements rule out an infimum.
This covers all stable $(g,n)$ not in (\ref{eq:smallgn}).

For the application to moduli spaces, the local smoothing parameters
of a stable curve correspond to its nodes, and smoothing a selected
set contracts the corresponding edges. The stable-graph stratum
closure relation is therefore exactly this order; see
\cite[Proposition 4.14, Section 6.1, and Section 7.3]{Schmitt}.
Each common refined type occurs, by gluing smooth pointed components
of its prescribed genera. Thus the combinatorial failures are actual
failures of the closure-intersection condition in the stated moduli
stratification, not merely formal graphs.
\end{proof}

\part{The stratified homotopy hypothesis at its first levels}
\IfFileExists{B_levels.tex}{% AI-produced by GPT-6 Astra (OpenAI, ChatGPT Pro mode), 2026-11-19.
% This segment uses only amsmath, amssymb, amsthm, and standard LaTeX commands.
% Bibliographic keys are supplied in B_references.bib.

\section{B.1. What the two-stratum sheaf description actually requires}

We write $\operatorname{Loc}(T)$ for the category of locally constant sheaves
of sets on a topological space $T$. A local system is not assumed to have
finite stalks. We keep this distinction because it affects the answer.

\paragraph{B1.1 (Proved: unconditional closed--open recollement).}
Let $i:Y\hookrightarrow X$ be a closed inclusion of topological spaces and
let $j:U=X\setminus Y\hookrightarrow X$ be its complementary open inclusion.
There is an equivalence
\[
 \operatorname{Sh}(X)\simeq
 \{(L,E,\alpha):L\in\operatorname{Sh}(U),\ E\in\operatorname{Sh}(Y),\
             \alpha:E\longrightarrow i^{-1}j_*L\}.
 \tag{B1.1}
\]
Morphisms on the right are pairs of sheaf morphisms commuting with $\alpha$.
Its restriction to the full subcategory $\mathcal F$ of sheaves locally
constant on $U$ and $Y$ is obtained by imposing
$L\in\operatorname{Loc}(U)$ and $E\in\operatorname{Loc}(Y)$. No normal
triviality, local contractibility, separation, or finiteness assumption is
needed for this statement.

\begin{proof}
Put $T=i^{-1}j_*L$ and form the pullback in $\operatorname{Sh}(X)$
\[
 F=j_*L\mathbin{\times}_{i_*T}i_*E,
 \tag{B1.2}
\]
where the first map is the adjunction unit and the second is $i_*\alpha$.
At $x\in U$ the two sheaves supported on $Y$ have terminal stalk, so
$F_x=L_x$. At $y\in Y$ the pullback on stalks is
$T_y\times_{T_y}E_y=E_y$. The resulting closed--open comparison is
$\alpha$. Conversely, adjunction supplies a map from every sheaf $F$ to
the pullback constructed from $(j^{-1}F,i^{-1}F)$; its map on every stalk
is an isomorphism. Sheaves on a topological space have enough points,
so this is an isomorphism. The pullback universal property proves the
assertion on morphisms. Empty $Y$ or $U$ gives the corresponding identity
equivalence. This is the explicit form of closed--open recollement; compare
\cite[Example A.8.4]{LurieHA}.
\end{proof}

\paragraph{B1.2 (Proved under an exact comparison hypothesis).}
Suppose in addition that $N\subset X$ is an open neighbourhood of $Y$ with
a retraction $r:N\to Y$, and write $N^*=N\setminus Y$ and
$p=r|_{N^*}:N^*\to Y$. For $L\in\operatorname{Loc}(U)$ there is a natural
germ comparison
\[
 \theta_L:p_*(L|_{N^*})\longrightarrow i^{-1}j_*L.
 \tag{B1.3}
\]
Indeed a section over $p^{-1}W$ gives a germ along the open subset
$W\subset Y$ by using the neighbourhood $r^{-1}W$.
If every $\theta_L$ is an isomorphism, then
\[
 \mathcal F\simeq
 \{(L,E,\varphi):L\in\operatorname{Loc}(U),\ E\in\operatorname{Loc}(Y),\
       \varphi:p^{-1}E\longrightarrow L|_{N^*}\}.
 \tag{B1.4}
\]
More generally, without assuming (B1.3) to be an isomorphism, the precise
necessary and sufficient condition for this functor of systems to be an
equivalence over the fixed pair $(L,E)$ is that
\[
 \operatorname{Hom}_Y(E,p_*(L|_{N^*}))\longrightarrow
 \operatorname{Hom}_Y(E,i^{-1}j_*L)
 \tag{B1.5}
\]
be bijective. Requiring the description after restriction to every open
subset of $Y$, and for every locally constant $E$, makes the isomorphism
condition on (B1.3) necessary as well as sufficient.

\begin{proof}
Adjunction identifies the set of maps $\varphi$ with the source of
(B1.5). Under (B1.1), its image is exactly the right-hand gluing map.
Thus bijectivity on each pair of objects is exactly full faithfulness
and essential surjectivity over the two restriction categories. It also
gives the assertion for varying objects, since morphisms are the pairs
commuting with their gluing maps. For the local necessity assertion we
take the constant one-point sheaf on each open $W\subset Y$: maps from
that sheaf are sections over $W$. Bijectivity for all $W$ is precisely
that $\theta_L$ is an isomorphism of sheaves.
\end{proof}

The usual two-stratum conical tube satisfies this criterion. More
explicitly, suppose $Y$ is paracompact and has a basis of contractible
open sets, that
$r:N\to Y$ is locally a bundle of open cones on compact, locally
contractible links, and that $U$ is locally path connected and semilocally
simply connected. We allow disconnected links. In a contractible base
chart $W$, the punctured tube has the form
$W\times K\times(0,\epsilon)$. A local system there is pulled back from
$K$, and its sections are unchanged by shrinking $W$ and $\epsilon$.
Compactness of $K$ makes the constant-radius conical neighbourhoods
cofinal at each point of $Y$. Consequently (B1.3) is an isomorphism on
stalks. Radial contraction also identifies
$\operatorname{Loc}(N)$ with $\operatorname{Loc}(Y)$. Monodromy now gives
the groupoid formulation in 151/035--036, with
$\Pi_1N\simeq\Pi_1Y$. Here monodromy means transport in an actual
locally constant sheaf; the assumptions on the spaces make this an
equivalence with all set-valued functors on their fundamental groupoids.

\paragraph{B1.3 (Refuted: an unconditional uniform pro-tube formula).}
The assertion on 151/035 is false if, without normal or neighbourhood
stability hypotheses, a map on a pro-tube means a morphism represented
on one sufficiently small neighbourhood. This fails even for a compact
metric $X$ with two strata, one a point and the other a disjoint union
of one-dimensional manifolds.

\begin{proof}
Let $X$ be the Hawaiian earring in $\mathbb R^2$: its $n$th circle $C_n$
has centre $(1/n,0)$ and radius $1/n$. Let $Y=\{o\}$, where $o=(0,0)$,
and $U=X\setminus\{o\}$. Each $C_n\setminus\{o\}$ is an open-and-closed
component of $U$, homeomorphic to an open interval. Let $L$ be the
constant two-point local system on $U$, and let $E=\mathbb N$ on $Y$.
For each $n$, choose a sufficiently small ball about $o$ that misses the
point of $C_n$ opposite $o$. On its punctured intersection with $X$,
there is a locally constant function that is $1$ on the upper arm of
$C_n$, $0$ on its lower arm, and $0$ on all other circles. Denote its
germ by $a_n\in(i^{-1}j_*L)_o$. The map $\alpha(n)=a_n$ defines an
object of (B1.1), hence a sheaf $F\in\mathcal F$.

Suppose it came from the manuscript's ordinary tube system. A locally
constant sheaf on the tube, with stalk $\mathbb N$ at $o$, is constant
on some single open neighbourhood $W$ of $o$. The map to $L$ on the
deleted tube would therefore be represented on a smaller such $W$.
Every neighbourhood of $o$ contains an entire $C_n$ for all sufficiently
large $n$. On the connected set $C_n\setminus\{o\}$, a morphism from
the constant $\mathbb N$-sheaf to the constant two-point sheaf sends
each fixed element to a constant value. This contradicts the two
distinct arm-values in $a_n$. Relabelling the stalk $\mathbb N$ does
not change the contradiction. No uniform representative exists.

There are no omitted hypotheses in this counterexample: $X$ is compact
metric; $Y$ is closed; $U$ is a disjoint union of open intervals; and
both restrictions of $F$ are locally constant by (B1.2). What fails is
uniform stabilization of neighbourhood gluing. In particular, the
unconditional statement (B1.1), not (B1.4), is the correct replacement.
The example also explains why one cannot assert that the pro-neighbourhood
of an arbitrary closed subset is homotopy equivalent to that subset.
\end{proof}

There is a second independent proviso in the wording of 151/035:
ordinary fundamental groupoids do not classify all local systems on
arbitrary spaces. For example, on a nondiscrete Cantor space the
fundamental groupoid has only identity arrows, so a representation may
assign a singleton to one non-isolated point and a two-point set to
every other point. No locally constant sheaf has those stalk
cardinalities. Thus the purely groupoid wording also requires a
monodromy hypothesis, even in the one-stratum case.

\paragraph{B1.4 (Proved: the inertia formula, with its precise scope).}
In the connected case let $D=\pi_1(N^*)$, $B=\pi_1(Y)$ and
$G=\pi_1(U)$, with homomorphisms $p_*:D\to B$ and $i_*:D\to G$.
Assume $p_*$ is onto and put $I=\ker p_*$. Then (18) of 151/036 is
equivalent to the $B$-equivariant map
\[
 E\longrightarrow (i_*^*L)^I.
 \tag{B1.6}
\]
Indeed a $D$-equivariant map from a $D$-set on which $I$ acts trivially
has image in the $I$-fixed points, and the induced action on those fixed
points factors through $D/I=B$. Conversely such a $B$-map is
$D$-equivariant. If $p$ is a fibration with connected fibre $K$, its
homotopy exact sequence gives
\[
 \pi_2(Y)\longrightarrow\pi_1(K)\longrightarrow D
      \longrightarrow B\longrightarrow1.
 \tag{B1.7}
\]
Thus $I$ is the quotient of $\pi_1(K)$ by the image of the connecting
homomorphism, and it equals $\pi_1(K)$ only when that homomorphism is
zero. This independently rechecks R4, SA10-C3-04. If $Y$ and $K$ are
aspherical, the same long exact sequence gives $\pi_q(N^*)=0$ for
$q\geq2$. Disconnected fibres are handled by the groupoid formulation;
surjectivity of $D\to B$ must not then be silently asserted.

\paragraph{B1.5 (Proved: the exact criterion for pro-tubes).}
We give a criterion that does not assume a chosen retraction or a
stationary representative. Let $\mathcal N_Y$ be the filtered ordered
set of open neighbourhoods of $Y$, with arrows given by shrinking.
Form the filtered bicategorical colimits in categories
\[
 \begin{aligned}
 \mathcal T&=\operatorname*{2colim}_{N\in\mathcal N_Y}
                         \operatorname{Loc}(N),\\
 \mathcal T^*&=\operatorname*{2colim}_{N\in\mathcal N_Y}
                         \operatorname{Loc}(N\setminus Y).
 \end{aligned}
 \tag{B1.8}
\]
Their objects and arrows have representatives on one neighbourhood;
two representatives agree if they agree after a common shrinking.
There are restriction functors
$r:\mathcal T\to\operatorname{Loc}(Y)$,
$a:\mathcal T\to\mathcal T^*$, and
$b:\operatorname{Loc}(U)\to\mathcal T^*$.
For $E\in\mathcal T$ and $L\in\operatorname{Loc}(U)$, restriction of
a representative morphism to germs along $Y$ defines
\[
 \mu_{E,L}:\operatorname{Hom}_{\mathcal T^*}(aE,bL)
 \longrightarrow
 \operatorname{Hom}_{Y}(rE,i^{-1}j_*L).
 \tag{B1.9}
\]
The natural functor from systems $(L,E,\varphi:aE\to bL)$ to
$\mathcal F$ is an equivalence, compatible with restriction to $U$
and $Y$, if and only if both of the following conditions hold:
$r$ is an equivalence, and every map (B1.9) is bijective.
These are respectively an extension condition for local systems and a
uniform extension condition for their gluing morphisms. For a
cofinal system of actual tube retractions, B1.2 supplies sufficient
geometric hypotheses for both conditions.

\begin{proof}
For a representative $E_N$ and a representative
$\varphi:E_N|_{N'\setminus Y}\to L|_{N'\setminus Y}$ with $N'\subset N$,
adjunction gives $E_N|_{N'}\to j'_*L|_{N'\setminus Y}$.
Its inverse image on $Y$ is (B1.9); further shrinking gives the same
map. This proves that $\mu$ and the asserted functor are well defined.
If $r$ is an equivalence, choose a quasi-inverse and normalize the
systems so that their second entry is an object of
$\operatorname{Loc}(Y)$. The two categories are then Grothendieck
constructions over the same category
$\operatorname{Loc}(U)\times\operatorname{Loc}(Y)$, with fibres the
sets on the two sides of (B1.9). Their arrow conditions are equality
of the composites along the two restrictions. The natural bijections
$\mu$ identify those conditions, proving sufficiency.

Conversely, set $L$ equal to the terminal local system on $U$.
There is a unique possible $\varphi$, so the systems with this
restriction form $\mathcal T$; the corresponding sheaves form
$\operatorname{Loc}(Y)$ by B1.1. Compatibility with restrictions
therefore forces $r$ to be an equivalence. After the preceding
normalization, the comparison has identical base objects and arrows.
Its fibre over a fixed $(L,E)$, with the identity base arrow, is a
discrete category. Both projections allow transport along every
isomorphism of base pairs: conjugate the gluing arrow by the two
given isomorphisms. Thus an equivalence over the base is an
equivalence on these fibres. This is exactly bijectivity
of $\mu_{E,L}$. This proves necessity.
\end{proof}

For groupoid wording, one must additionally impose the monodromy
equivalence on the spaces in a cofinal neighbourhood system and on
$U,Y$, or explicitly define the groupoids as those classifying these
local systems. Ordinary $\Pi_1$ is sufficient on locally path connected,
semilocally simply connected spaces. The pro-representation convention
in (B1.8) keeps the one-stage representability of arrows; it must not
be replaced silently by independent choices of a neighbourhood for
each element of an infinite stalk. In B1.3, $r$ is an equivalence
because every local system near the single point is constant after
shrinking. Precisely the surjectivity of $\mu_{\mathbb N,L}$ fails.

\section{B.2. Shape and the level at which gluing takes place}

\paragraph{B2.1 (Known, with an explicit derivation: the correct shape theorem).}
Let $\mathcal S$ be the $\infty$-category of spaces and
$\operatorname{Top}_\infty$ the $\infty$-category of $\infty$-topoi and
geometric morphisms. The shape $\Pi_\infty\mathcal X\in
\operatorname{Pro}(\mathcal S)$ is characterized by
\[
 \operatorname{Map}(\Pi_\infty\mathcal X,K)
  \simeq \Gamma_{\mathcal X}\Delta_{\mathcal X}K,
       \qquad K\in\mathcal S.
 \tag{B2.1}
\]
For every small diagram whose colimit in $\operatorname{Top}_\infty$
is $\mathcal X$, the canonical map
\[
 \operatorname*{colim}_a\Pi_\infty\mathcal X_a
       \longrightarrow\Pi_\infty\mathcal X
 \tag{B2.2}
\]
is an equivalence of pro-spaces. No stratification or normal
fibration hypothesis occurs here.

\begin{proof}
For a space $K$, the universal property of the slice $\infty$-topos
$\mathcal S_{/K}$ identifies
$\operatorname{Map}_{\operatorname{Top}_\infty}
(\mathcal X,\mathcal S_{/K})$ with the right side of (B2.1).
Mapping a colimit to this slice gives
\[
 \operatorname{Map}(\Pi_\infty\mathcal X,K)
 \simeq\lim_a\operatorname{Map}(\Pi_\infty\mathcal X_a,K).
\]
The fully faithful realization of pro-spaces as accessible left exact
functors $\mathcal S\to\mathcal S$, with variance reversed, detects
equivalences by these evaluations. This proves (B2.2). Equivalently,
shape is a left adjoint, as stated in
\cite[Remark 7.1.6.15]{LurieHTT} and
\cite[Definition 4.2.1]{BGHExodromy}.
\end{proof}

For an actual open cover of $X$, descent for sheaves of spaces gives
the colimit of sheaf $\infty$-topoi indexed by its full descent diagram.
Thus (B2.2) applies to every honest open-cover stage of a pro-tube
construction. If the pieces are locally of singular shape, their shapes
are represented by their singular homotopy types. This gives the
homotopy-colimit statement used in topology. Passing from one stage to
the pro-system is legitimate as a statement about the pro-system of
these colimits; it is not permission to commute an arbitrary infinite
homotopy colimit with a pro-limit. Finite homotopy colimits, indexed
by finite simplicial sets such as a span, can be computed levelwise
after a common cofiltered reindexing. A category with finitely many
objects and arrows need not have finite nerve, so that observation
does not justify truncating an untruncated higher descent diagram
to $\Delta_{\leq2}$. We retain the full stagewise \v Cech diagram.

More explicitly, let $A$ be the cofiltered category of synchronized
neighbourhood choices in A.1, with an arrow $\alpha'\to\alpha$ for
a shrinking. Write $C_q^\alpha$ for the disjoint union of all
$(q+1)$-fold intersections of its actual open cover. The precise
higher statement is the natural equivalence
\[
 \begin{aligned}
 &\left(\operatorname*{colim}_{[q]\in\Delta^{\mathrm{op}}}
       \Pi_\infty\operatorname{Sh}_\infty(C_q^\alpha)\right)_{\alpha\in A}
 \\
 &\hspace{3em}\simeq
 \operatorname{const}_A\bigl(\Pi_\infty\operatorname{Sh}_\infty(X)\bigr)
 \end{aligned}
 \tag{B2.2a}
\]
in $\operatorname{Fun}(A,\operatorname{Pro}(\mathcal S))$.
The comparison is induced by the cover augmentations, and is an
equivalence at every $\alpha$, which proves the displayed natural
equivalence. If every shape in this formula is represented by the
corresponding singular homotopy type, it also says that the resulting
cofiltered system of homotopy colimits is the constant pro-space
$\operatorname{Sing}(X)$. No exchange of the two indexing operations
is made.

The Artin--Mazur qualification also matters. The coherent hypercover
construction models the shape of the \emph{hypercomplete} sheaf
$\infty$-topos: see \cite[Proposition 5.1]{HoyoisGalois}. It yields
the classical pro-homotopy object after the stated loss of coherence.
For the locally contractible spaces used in the examples below these
distinctions do not alter the answer. They cannot be erased for an
arbitrary topos; compare \cite[Example 7.1.6.9]{LurieHTT}.

\paragraph{B2.2 (Refuted: arbitrary bicategorical colimits of ordinary topoi).}
The same sentence is false if ``colimit of topoi'' means a bicategorical
colimit of ordinary Grothendieck topoi, and the right side means a
homotopy colimit of their shapes.

\begin{proof}
Let $G=C_2$ and let $BG$ denote the ordinary topos of $G$-sets. Its
unique geometric morphism $q:BG\to\mathbf{Set}$ has fully faithful
inverse image $q^*$, the trivial-action functor. The bicategorical
pushout of
\[
 \mathbf{Set}\longleftarrow BG\longrightarrow\mathbf{Set}
 \tag{B2.3}
\]
is $\mathbf{Set}$. To verify the universal property, fix any target
topos $\mathcal E$. A pseudocone consists of two geometric morphisms
$a,b:\mathbf{Set}\to\mathcal E$ and an invertible transformation
$q^*a^*\cong q^*b^*$. Full faithfulness of $q^*$ identifies that
transformation, and all transformations of pseudocones, with an
invertible transformation $a^*\cong b^*$ in $\mathbf{Set}$. The category
of two objects with a specified isomorphism between them is equivalent
to the original category by evaluation on the first object. This is
exactly the bicategorical pushout universal property.

The associated sheaf $\infty$-topos of $BG$ is
$\operatorname{Fun}(BG,\mathcal S)$, whose shape is the classifying
space of $G$; see \cite[Example 4.2.3]{BGHExodromy}, or compute global
sections of a constant diagram directly. Hence the homotopy pushout of
the three shapes is $\Sigma BC_2$. Its reduced second integral homology
is $H_1(BC_2;\mathbb Z)=\mathbb Z/2$, while the shape of the
bicategorical pushout is a point. The program
\texttt{check\_B\_finite.py} verifies the normalized bar differentials
$d_1=0,d_2=2,d_3=0,d_4=2$ and the finite group operations.
\end{proof}

This counterexample identifies the exact failure: passage from an
ordinary topos to its associated sheaf $\infty$-topos does not preserve
all bicategorical colimits. The actual open-cover gluings satisfy
descent at every higher level, so they are not affected by this
counterexample. We must prove that higher descent for the diagram
under consideration before applying B2.1.

\paragraph{B2.3 (Known under precise hypotheses: Quinn's input and L2).}
For a metric space $X$ and a closed subset $Y$, put
\[
 H=\operatorname{holink}(X,Y)
  =\{\gamma:[0,1]\to X:\gamma(0)\in Y,
                 \ \gamma(t)\notin Y\ (t>0)\}.
\]
The maps $e_0:H\to Y$ and $e_1:H\to U$ are endpoint evaluations.
If $Y$ is forward tame in $X$, then the ordinary homotopy types satisfy
\[
 X\simeq Y\mathbin{\cup^h}_{H}U.
 \tag{B2.4}
\]
Quinn proves a controlled strict relative version in
\cite[Lemma 2.4(2)]{Quinn1988}, using his controlled holink and
mapping-cylinder convention. We give a separate proof of (B2.4)
that does not identify that cylinder topology with the ordinary
quotient-cylinder topology. Let $N$ range over the open neighborhoods
of $Y$. The numerable open cover $X=N\cup U$ gives, naturally in $N$,
the homotopy-pushout equivalence
$N\cup^h_{N\setminus Y}U\simeq X$.
A fixed tame deformation gives an equivalence
$\{N\}_N\simeq\operatorname{const}Y$ in prospaces: for any target
neighborhood $N$, shrink to the points whose entire deformation lies
in $N$. The same fixed deformation and the explicit endpoint-preserving
path homotopy of C4.02 give
$\{N\setminus Y\}_N\simeq\operatorname{const}H$.
Those formulas need only tameness of this pair; the fibration part
of the HSS definition is not used. The comparison uses the natural
zigzag through paths wholly contained in $N$, so its endpoint maps
and the homotopy between the two maps to $X$ are the original ones.
The finite span colimit in $\operatorname{Pro}(\mathcal S)$ is
computed on these synchronized levels, because filtered colimits
of spaces commute with finite simplicial limits. Replacing the two
pro-germs by their constant equivalents therefore gives
$\operatorname{const}X\simeq
\operatorname{const}(Y\cup^h_HU)$.
The constant embedding is fully faithful, proving (B2.4).
In the usual quotient double-cylinder model, the associated map to
$X$ is induced by the continuous path evaluation
$H\times[0,1]\to X$.
Quinn's Definition 3.1 adds, for each incident pair of strata, that
$e_0$ is a fibration. That is the precise normal-fibration analogue of
151/032, once a deleted tube has been compared with the holink. It is
not the whole definition: forward tameness remains a separate
hypothesis. The fibration condition alone does not imply tameness.
For the Hawaiian earring with $Y$ its common point, $e_0$ is a fibration
because its target is a point. But $Y$ is not tame: a neighbourhood
deformation into $Y$ would null-homotope every sufficiently small
circle in $X$, while retraction of $X$ onto any one of those circles
shows its fundamental loop is nonzero. Consequently L2 must be read
with this additional tameness and comparison data, and its unrestricted
topos-colimit conclusion must be replaced by B2.1. The shape theorem
itself needs neither condition.

\paragraph{B2.4 (Proved: rechecking the two examples in R5).}
For $(\mathbb C,\{0\})$, B1.4 identifies a constructible sheaf with a
set $E$, a $\mathbb Z$-set $M$, and a map $E\to M^{\mathbb Z}$.
Equivalently it is a covariant set-valued functor on the category with
objects $p,q$, automorphism group $\mathbb Z$ at $q$, and a unique map
$p\to q$. The object $p$ is initial; hence the classifying space of
this category, and the shape of its functor topos, is contractible.
This agrees with $\mathbb C$.

For $(S^2,\{p\})$, the complementary plane is contractible. The same
description reduces to a single map of sets $E\to M$. The corresponding
arrow category also has contractible classifying space. Nevertheless
the genuine tube span has shapes
\[
 *\longleftarrow S^1\longrightarrow *,
\]
whose homotopy pushout is $\Sigma S^1\simeq S^2$. Thus the ordinary
constructible-sheaf topos does not retain the sphere's higher homotopy
type. This proves the assertions SA10-C3-05 and SA10-C3-06 without
identifying that functor topos with the required higher gluing.

\section{B.3. Exit paths, truncation, and gerbe gluing}

\paragraph{B3.1 (Known; truncation proved below).}
Assume that $X$ is paracompact, locally of singular shape, and conically
stratified by a poset $P$ satisfying the ascending chain condition.
These hypotheses hold, for example, for a finite-dimensional,
paracompact Hausdorff topologically stratified space with the compact
conical charts used in \cite[Definition 3.1]{Treumann2009}.
Write $\operatorname{Exit}(X)=\operatorname{Sing}_P(X)$.
For every integer $n\geq0$ there are equivalences
\[
 \begin{split}
 \operatorname{Cons}_{\leq n}(X)
 &\simeq \operatorname{Fun}(\operatorname{Exit}(X),\mathcal S_{\leq n})\\
 &\simeq \operatorname{Fun}(h_{n+1}\operatorname{Exit}(X),
                                      \mathcal S_{\leq n}).
 \end{split}
 \tag{B3.1}
\]
Here the left side consists of constructible $n$-truncated sheaves of
spaces, and $h_{n+1}$ retains the objects and the $n$-truncations of
all mapping spaces, with their coherent compositions.
The regularity hypotheses on the strata are substantive. If
``conical'' only means a local product with a cone, allowing an
arbitrary unstratified base, every one-stratum space is conical.
The Cantor-space example after B1.3 then refutes the unqualified
claim already at $n=0$, and also at every higher level by regarding
sets as discrete spaces. For the usual finite Thom--Mather,
Whitney, or conically smooth spaces with manifold strata, the
hypotheses used here hold. Disconnected and empty strata cause no
change: the statement is componentwise and an empty stratum
contributes no objects.

\begin{proof}
The untruncated equivalence is
\cite[Theorem A.9.3]{LurieHA}. Its functor evaluates on stalks as the
corresponding functor evaluates on objects, by
\cite[Proposition A.9.16]{LurieHA}. A constructible object is
$n$-truncated precisely when all these stalks are $n$-truncated:
constructible objects here are hypercomplete, and the assertion also
follows directly from their local conical descriptions. Thus the
equivalence restricts to the first line of (B3.1); this is also recorded
in \cite[Remark A.9.7]{LurieHA}. Mapping spaces in
$\mathcal S_{\leq n}$ are $n$-truncated, since they are limits of
$n$-truncated spaces. The defining universal property of the homotopy
$(n+1)$-category therefore gives the second equivalence, including its
mapping spaces and not merely its objects.
\end{proof}

We give the finite-stratum induction explicitly, including the descent
reduction; this avoids identifying a global undercategory with a single
local link. We prove the untruncated comparison and simultaneously all
its truncated restrictions by induction on the number of nonempty
stratum labels. The empty case is terminal, and the one-stratum case
is ordinary monodromy, \cite[Theorem A.4.19]{LurieHA}.
For an open $F_\sigma$ subset $W\subset X$, let $\mathcal B_W$ be the
category of its conical open $F_\sigma$ charts, as in
\cite[Lemma A.5.16]{LurieHA}. The opposite of the category of such charts
containing any fixed point is filtered, since smaller conical charts
refine intersections. Thus its nerve is contractible. Exit excision
\cite[Theorem A.7.1]{LurieHA} and descent for constructible sheaves give
\[
 \begin{split}
 \operatorname{Fun}(\operatorname{Exit}(W),\mathcal S)
 &\simeq\lim_{V\in\mathcal B_W^{\mathrm{op}}}
              \operatorname{Fun}(\operatorname{Exit}(V),\mathcal S),\\
 \operatorname{Cons}(W)
 &\simeq\lim_{V\in\mathcal B_W^{\mathrm{op}}}\operatorname{Cons}(V).
 \end{split}
 \tag{B3.2}
\]
The same holds for $n$-truncated values, since truncation is local.
The comparison is natural for these restrictions, so an equivalence
on the charts implies an equivalence on $W$ by (B3.2).

It remains to prove the induction step on a chart
$V=Z\times C(K)$, with closed stratum $Z$ and
$V^*=Z\times K\times(0,\infty)$. Here $V^*$ has fewer nonempty
stratum labels. We spell out the categorical induction criterion, so
that no global nearby-cycle stability is presumed.
In $\mathcal C=\operatorname{Fun}(\operatorname{Exit}(V),\mathcal S)$
let $\mathcal C_0$ be functors extended by the initial space from
$\operatorname{Exit}(V^*)$, and let $\mathcal C_1$ be functors sending
every arrow to an equivalence. In
$\mathcal D=\operatorname{Cons}(V)$ let $\mathcal D_0=j_!
\operatorname{Cons}(V^*)$ and let $\mathcal D_1$ be the local systems
pulled back from $Z$. The inclusions of each of these pairs have
complementary right adjoints. On the functor side this follows from
the cone contraction and the dual of \cite[Lemma A.9.18]{LurieHA};
on the sheaf side it is \cite[Lemma A.9.13]{LurieHA}.
The dual of the recollement comparison criterion
\cite[Proposition A.8.14]{LurieHA} now applies for the following
explicit reasons. The restriction
$\mathcal C_0\to\mathcal D_0$ is the comparison on $V^*$, hence an
equivalence by induction. The restriction
$\mathcal C_1\to\mathcal D_1$ is monodromy on $Z$, because radial
contraction identifies the locally constant objects on $V$ with
those on $Z$; the precise sheaf comparison is
\cite[Lemmas A.9.11--A.9.12]{LurieHA}. The comparison preserves
pushouts by \cite[Lemma A.9.15]{LurieHA}. Finally it carries the
$\mathcal C_0$-colocalization of an object of $\mathcal C_1$ to the
corresponding $\mathcal D_0$-colocalization: both are obtained by
restriction to $V^*$ followed by extension by the initial object.
This verifies every hypothesis of the criterion, proving
$\mathcal C\simeq\mathcal D$ and completing the induction.

All these comparisons respect stalks and evaluation, so the full
subcategories whose values are $n$-truncated agree at each induction
step. Mapping into $\mathcal S_{\leq n}$ then factors uniquely through
$h_{n+1}\operatorname{Exit}(X)$ by the already proved universal
property. This proves every truncation level by the same induction.
For infinite ACC posets, the source theorem replaces the finite
induction by rank induction; finite induction alone would not prove
that extension.

The associated closed--open gluing operation is
\[
 (F_U,F_S,\alpha:F_S\longrightarrow i^*j_*F_U).
 \tag{B3.3}
\]
It is the homotopy-pullback version of (B1.2). In an actual conical
chart with contractible base and compact link $K$, its nearby value
is $\lim_{\operatorname{Exit}(K)}F_U$. On the exit side this is the
right Kan extension from the open subcategory. The open cone adds
an initial cone point. For conically smooth charts the exact cone
and sheaf comparisons are \cite[Lemmas 6.1.4 and 6.1.7]{AFTLocal},
taking opposites by its Corollary 1.2.10 to convert enter paths to exit paths;
the general topological conical case is included in the foregoing
recollement argument. Limits of $n$-truncated spaces remain
$n$-truncated, which also makes the compatibility of (B3.3) with
every finite level explicit.

\paragraph{B3.2 (Proved: groupoid stacks and the correction to L3).}
At $n=1$, (B3.1) is the equivalence between constructible stacks in
groupoids and homotopy-coherent $2$-functors from the exit-path
$2$-category to groupoids. It is also the restriction of
\cite[Theorem 7.14]{Treumann2009} from category-valued to
groupoid-valued stacks: the value at each object is the stalk
category, so the condition that these values are groupoids is
preserved in both directions. His Theorem 7.13 proves compatibility
with open-cover descent.

The set-valued level, by contrast, is $n=0$: it retains $\pi_0$ of
each exit mapping space. Thus L3's phrase ``truncated at $\pi_1$ of
the links'' is one level too high if it means that link $\pi_1$
survives as automorphisms of exiting arrows. Those automorphisms
first appear in the stack level. For $(S^2,\{p\})$, the two-stratum
exit category can be represented by two objects and cross mapping
space $S^1$: the open stratum is contractible, and its attaching
holink has the homotopy type of $S^1$. Taking $\pi_0$ makes a single
arrow, while the fundamental groupoid of that mapping space retains
$\mathbb Z$. This is exactly the difference seen by the gerbes.

One should also distinguish two uses of gluing. A general constructible
stack is a closed--open system (B3.2), with a possibly noninvertible
functor to nearby data. An $A$-banded gerbe is locally equivalent to
$BA$ at every point of $X$. Its gluing over an honest open tube and
$U$ uses an \emph{equivalence} on their overlap, as in 151/042.
These gerbes form a subproblem of the constructible-stack theorem;
arbitrary constructible stacks are not all gerbes.

\paragraph{B3.3 (Proved: the gerbe computation, including the quotient by gauges).}
Let $A$ be any discrete abelian group with constant band. Then
\[
 H^2(S^2;A)\cong A.
 \tag{B3.4}
\]
We obtain this by the precise gluing on 151/041--043.

\begin{proof}
Cover $S^2$ by $U=S^2\setminus\{p\}$ and an open disk $N$ about $p$.
Both are contractible, and $U\cap N$ is an annulus. Every $A$-banded
gerbe on either contractible piece is equivalent to the neutral
gerbe $BA$. An equivalence between two copies of $BA$ on an overlap
is an $A$-torsor there; composition is contracted product of torsors.
The isomorphism classes of these torsors are
\[
 H^1(U\cap N;A)=\operatorname{Hom}(\mathbb Z,A)=A.
\]
Changing either neutralization acts by the restrictions of torsors
on $U$ and $N$. Those torsors have trivial isomorphism class because
$H^1(U;A)=H^1(N;A)=0$. Consequently this action identifies no two
distinct elements of $A$. Open-cover $2$-descent is effective, so
every such torsor glues the two neutral gerbes, and two glued gerbes
are equivalent precisely when the torsor classes agree. This proves
(B3.4), using the interpretation of $H^2$ by banded gerbes; the
precise abelian classification used here is
\cite[Theorem 2.7]{DeopurkarGerbes}. The identification of
band-preserving self-equivalences of the neutral gerbe with torsors
is also explicit in its Remark 3.1: the image of the trivial torsor
determines the equivalence, which is contracted product with that
torsor (or its inverse according to the chosen action convention).

In the extension language of the manuscript, both pieces yield
trivial central extensions. An automorphism of the pulled-back
extension $A\times\mathbb Z\to\mathbb Z$, fixing the band and the
quotient, is uniquely $(a,m)\mapsto(a+mh,m)$ for $h\in A$.
Thus the same gluing set is $\operatorname{Hom}(\mathbb Z,A)$.
Choices of basepoints and neutralizations are why this is a
description up to equivalence of gerbes, not a canonical literal
identity between a gerbe and a group extension. This rechecks
SA10-C3-07.
\end{proof}

\section{B.4. The aspherical programme and the coherence it needs}

\paragraph{B4.1 (Proved: the premise in the question is refuted).}
For a two-stratum homotopy-pushout decomposition
$X\simeq Y\cup^h_H U$, suppose every component of $Y,H,U$ is a
$K(\pi,1)$ of CW homotopy type. Then the diagram
\[
 \Pi_1Y\longleftarrow\Pi_1H\longrightarrow\Pi_1U
 \tag{B4.1}
\]
\emph{does} determine the full weak homotopy type of $X$:
\[
 X\simeq B\Pi_1Y\mathbin{\cup^h}_{B\Pi_1H}B\Pi_1U.
 \tag{B4.2}
\]
No extra invariant is required for this span. Thus R5 does not
establish the insufficiency asserted in the last sentence of B.4.

\begin{proof}
For every space $T$ there is a natural map of simplicial sets
$\operatorname{Sing}(T)\to N\Pi_1(T)$ to its fundamental groupoid
nerve. If every component of $T$ is aspherical, this map is a weak
equivalence. Apply it objectwise to the span. A natural objectwise
equivalence of diagrams induces an equivalence on homotopy
colimits; composing with the hypothesized homotopy-pushout
comparison gives (B4.2). This is a natural zigzag of weak
equivalences, so it does not require arbitrary basepoint identifications.

For the sphere, (B4.1) is $1\leftarrow\mathbb Z\to1$, and (B4.2)
is $*\cup^h_{S^1}*\simeq S^2$. It is the set-valued functor topos
that loses this information. On 151/041 Grothendieck explicitly
distinguishes these assertions and withdraws the feared failure of
groupoid-span reconstruction. The prompt's interpretation, not
that sentence of the manuscript, requires correction.
\end{proof}

The hypothesis on $H$ follows from the programme's local hypotheses
when $e_0:H\to Y$ is a fibration with aspherical fibres and $Y$ is
aspherical: the homotopy exact sequence proves asphericity of each
component of $H$. This works with disconnected fibres as well;
groupoids avoid an unjustified surjectivity assertion on $\pi_1$.

\paragraph{B4.2 (Proved: the full groupoid assembly data).}
Let $X\to P$ satisfy the hypotheses of B3.1, and put
$\mathcal E=\operatorname{Exit}(X)$. For every nonempty strict chain
$\sigma=(p_0<\cdots<p_m)$ in $P$, define
\[
 D(\sigma)=\operatorname{Map}_{/P}(\sigma,\mathcal E).
 \tag{B4.3}
\]
This is a space: a natural transformation over the fixed chain has
components inside strata, hence all its components are equivalences.
Restriction along subchains defines a homotopy-coherent functor
$D:\operatorname{sd}(P)^{\mathrm{op}}\to\mathcal S$.
Without any asphericity assumption, the complete assembly formula is
\[
 \operatorname{Sing}(X)\simeq
 \operatorname*{colim}_{\sigma\in\operatorname{sd}(P)^{\mathrm{op}}}
 D(\sigma).
 \tag{B4.3a}
\]
Assume that every stratum $D(p)$ and every total incidence link
$D(p<q)$ is a $1$-type. Then the groupoids
\[
 G_\sigma=\Pi_1D(\sigma),
 \qquad
 G:\operatorname{sd}(P)^{\mathrm{op}}\longrightarrow\operatorname{Gpd},
 \tag{B4.4}
\]
with \emph{all} restriction functors and their coherent comparison
isomorphisms determine the full homotopy type:
\[
 \operatorname{Sing}(X)\simeq
       \operatorname*{colim}_{\sigma\in\operatorname{sd}(P)^{\mathrm{op}}}
                  BG_\sigma.
 \tag{B4.5}
\]
The colimit in (B4.5) is taken in spaces, not in ordinary groupoids.
The incidence-link hypothesis holds when the local homotopy links
are aspherical and their source maps are fibrations over
aspherical strata. In geometric applications, ``all links'' must
mean all incident pairs, not merely adjacent pairs or a union of
different link strata.

\begin{proof}
The Segal equivalence gives
\[
 D(p_0<\cdots<p_m)\simeq
 D(p_0<p_1)\mathbin{\times^h}_{D(p_1)}\cdots
       \mathbin{\times^h}_{D(p_{m-1})}D(p_{m-1}<p_m).
 \tag{B4.6}
\]
This follows by mapping the concatenation of the elementary edges,
which is the chain $\sigma$ in $\operatorname{Cat}_\infty$, into
$\mathcal E$ over $P$; it is
\cite[Lemma 2.7.3]{BGHExodromy}. Homotopy limits of $1$-types are
$1$-types. Thus every value of $D$ is a $1$-type and the natural
objectwise truncation $D\to BG$ is an equivalence of diagrams.

To identify its colimit, extend $D$ along strict chains to all
simplices of $P$. Repeating an index inserts equivalences inside a
stratum, and this extension is the complete Segal nerve of
$\mathcal E$; this is the construction and proof of
\cite[Theorem 2.7.4]{BGHExodromy}. Its geometric realization is
the groupoid completion $|\mathcal E|$. Colimits commute with this
left Kan extension, so
$\operatorname{colim}_{\operatorname{sd}(P)^{\mathrm{op}}}D
\simeq|\mathcal E|$. Finally the natural inclusion
$\operatorname{Exit}(X)\to\operatorname{Sing}(X)$ induces an
equivalence on groupoid completions by
\cite[Corollary A.9.4]{LurieHA}. These three specified comparisons
give (B4.5).

For conical or Quinn pairs, total incidence links are modeled by
the endpoint-variable homotopy-link spaces. The latter's source
fibration has the local normal link as fibre. For popath models
the comparison that fixes both endpoints is
\cite[Theorem 4.9]{Miller2009}. The long exact sequence of that
fibration proves the last assertion of the theorem.
\end{proof}

The proof of (B4.3a) is the colimit and groupoid-completion part of
the same argument, before truncating $D$. This also specifies the
answer if ``links are $K(\pi,1)$'' refers only to the whole normal
link of each stratum. That condition does not imply our pairwise
incidence hypothesis. For example, stratify the closed $4$-ball by
a boundary point $p$, the remaining boundary $S^3\setminus\{p\}$,
and its interior. All three strata are contractible; the whole
normal link at $p$ is a $3$-ball and the normal link along the
remaining boundary is a point. But the incidence link from $p$
into the boundary stratum is $S^2$, not a $1$-type. In this reading,
one retains the spaces $D(\sigma)$, equivalently their full
Postnikov towers with the bonding maps, $k$-invariants and coherent
restriction maps. Formula (B4.3a) proves that this is sufficient in
every case. Formula (B4.6) says that the higher information lies in
the incidence links and their composition maps, not in independent
new spaces for each longer flag. Under the stronger pairwise
asphericity reading, all those towers truncate to the explicit
groupoid data (B4.4).

Formula (B4.4) is the requested additional data in the many-stratum
case. Here $\operatorname{Gpd}$ means its $(2,1)$-category, or
equivalently $\mathcal S_{\leq1}$, and not the category obtained by
discarding natural isomorphisms. Explicitly, the datum consists of
groupoids $G_\sigma$, functors
$F_{\sigma\tau}:G_\tau\to G_\sigma$ for inclusions
$\sigma\subset\tau$, and natural isomorphisms
$\theta_{\sigma\tau\upsilon}:F_{\sigma\tau}F_{\tau\upsilon}
\Rightarrow F_{\sigma\upsilon}$. Identity restrictions are identities,
the corresponding comparisons are identities, and for
$\sigma\subset\tau\subset\upsilon\subset\nu$ we require the equality
\[
 \theta_{\sigma\tau\nu}\circ
          (F_{\sigma\tau}\theta_{\tau\upsilon\nu})
 =\theta_{\sigma\upsilon\nu}\circ
          (\theta_{\sigma\tau\upsilon}F_{\upsilon\nu}).
 \tag{B4.8}
\]
Both sides have source
$F_{\sigma\tau}F_{\tau\upsilon}F_{\upsilon\nu}$ and target
$F_{\sigma\nu}$. These are all the coherence conditions, since
natural transformations between functors of groupoids form sets;
there is no additional layer of higher modifications.
An equivalent Segal presentation is: stratum and incidence groupoids;
composition on triple flags; the natural isomorphisms comparing
successive restrictions; and their coherence on longer flags.
Taking one object per connected groupoid converts functors into
homomorphisms only after choices, and converts the comparison
isomorphisms into conjugating elements. Those elements must be kept.
A diagram merely in the homotopy category of groups or spaces
forgets them. Conversely the full coherent diagram above is a
complete answer because (B4.5) reconstructs $X$, and it is recovered
from $X$ by (B4.3)--(B4.4).

\paragraph{B4.3 (Proved: why coherence cannot simply be discarded).}
The trivial circle bundle $S^2\times S^1\to S^2$ and the Hopf
bundle $S^3\to S^2$ have locally the same aspherical pieces and the
same maps on fundamental groups over a good cover. Their gluing
coherences differ, and their total spaces have different fundamental
groups.

\begin{proof}
Take a good cover of $S^2$ with nerve the boundary of a tetrahedron.
For example, use the open stars of the four vertices in its
geometric triangulation, with a small compatible thickening if
desired. Every nonempty intersection is contractible. Restrict
either oriented circle bundle to every such intersection: the
restriction is trivial and has homotopy type $S^1$. Every inclusion
between these pieces induces the identity homomorphism of
$\mathbb Z$, after orienting the fibre. Therefore the diagram
of individual groups and homomorphisms in the homotopy category
is the constant $\mathbb Z$-diagram for both bundles.

Choose transition maps $g_{ij}:U_i\cap U_j\to S^1$ and real lifts
$\theta_{ij}$ on the contractible pair intersections. On a triple
intersection,
\[
 n_{ijk}=\theta_{ij}+\theta_{jk}-\theta_{ik}\in\mathbb Z
 \tag{B4.7}
\]
is constant. These integers are the natural-transformation
coherence of the groupoid diagram; changing lifts adds a
coboundary. For the trivial bundle the class is zero. For the
Hopf bundle it is a generator: its north--south trivializations
have equatorial transition map of degree one, which is exactly
the winding number measured by (B4.7). This can also be seen
directly from the standard map $(z_0,z_1)\mapsto[z_0:z_1]$
on $S^3\subset\mathbb C^2$, whose two fibre coordinates differ
by the phase of $z_1/z_0$. The program
\texttt{check\_B\_finite.py} checks a tetrahedral representative:
the cocycle $(1,0,0,0)$ pairs to $-1$ with the fundamental cycle
$(-1,1,-1,1)$, so it is not a coboundary.

The total spaces have $\pi_1(S^3)=0$ and
$\pi_1(S^2\times S^1)=\mathbb Z$, respectively. Thus the
discarded coherence changes even the fundamental group of the
assembled space. To view this as a stratified example, lift the
finite face stratification of the tetrahedral sphere to either
bundle. Each stratum is $S^1$ times an open simplex and is
aspherical. Each local normal link is a point, a finite discrete
space, or a circle, hence componentwise aspherical. The same
coherence issue remains in passing from the full flag groupoid
diagram to its homotopy-category shadow.
\end{proof}

We do not claim that the full coherent span in B4.1 fails in this
example: a coarser two-open decomposition records the clutching
integer in its maps of fundamental groups. The loss occurs
precisely when the full coherent diagram is replaced by the list
of its local groups and homomorphisms up to homotopy.
}{}
\IfFileExists{B56.tex}{\section{B.5. Comparison axioms and homotopy descent}\label{b.5.-comparison-axioms-and-homotopy-descent}

AI-produced by GPT-6 Astra (OpenAI, ChatGPT Pro mode), 19 November 2026.
We claim no priority. We read the supplied second-pass transcriptions of
cote 151, scans 003, 005, 011--015, 017--018, and independently inspected
the original scans 003, 005 and 017 at enlarged resolution. The last
subscript in the bottom relation of 005 remains visually unresolved;
its mathematical alternatives are evaluated separately. No independent
ink reading is claimed for the other scans.

Let \(A\) be a small ordinary category, \(\mathcal X\) a Grothendieck
topos, and \(\varphi:A\to\mathcal X\) a functor. Put \(U_a=\varphi(a)\)
and \(\widehat A=\operatorname{Fun}(A^{\mathrm{op}},\mathbf{Set})\). Its
cocontinuous extension is

\[L=\varphi_!: \widehat A\to\mathcal X,\qquad L(F)=\int^{a\in A}F(a)\cdot U_a.\]

Here the dot denotes copower. For a finite diagram \(F:J\to\widehat A\),
write \(c_F:L(\lim_JF)\to\lim_JL(F)\) for the canonical comparison. An
object \(U\) of \(\mathcal X\) represents the slice topos
\(\mathcal X/U\), not its set of global sections. This distinction is
essential for interpreting the requested homotopy colimit.

\subsection{B5.1. The forest statement and the full
condition}\label{b5.1.-the-forest-statement-and-the-full-condition}

\textbf{Status: Proved.} Hyp 1, 2 and \(2'\) of 151/011--014 are
respectively the epimorphicity of

\[L(1)\to1,\qquad L(ya\times yb)\to U_a\times U_b,\]

and

\[L(ya\times_{yc}yb)\to U_a\times_{U_c}U_b\]

for every cospan \(a\to c\leftarrow b\) in \(A\). These are equivalent
to \(c_F\) being epic for all diagrams indexed by the free category on a
finite oriented forest, including the empty forest. This is an
equivalence with a collection of tests; it does not assert that every
indexing category which is not a forest fails. For example, a category
with an initial object has its limit given by evaluation there.

\textbf{Proof.} The category of elements gives
\(L(F)=\operatorname{colim}_{(a,s)\in A/F}U_a\). The coproduct of its
terms maps epimorphically onto that colimit. Thus the image of each
representable comparison is exactly the image of the covering family of
cones written in the manuscript. Hyp \(1'\) is tautological because the
slice contains the identity of its target.

For completeness we prove the extension to arbitrary presheaves. Work
locally in \(\mathcal X\). An element of \(L(F)\) is represented by
\((a,x,s)\) with \(x\in U_a\) and \(s\in F(a)\), and equality is locally
witnessed by a finite zigzag of coend relations. In the internal
category \(E_\varphi\), an object is \((a,x)\) and an arrow
\((a,x)\to(b,z)\) is \(f:a\to b\) such that \(\varphi(f)x=z\). Hyp 2
says any two objects locally have a common predecessor; Hyp \(2'\) says
a cospan locally has a commutative completing square. Every finite
zigzag locally admits a commuting cone: induct on its edges, composing a
cone leg at a forward step and using a completing square at a backward
step. For disconnected diagrams combine the cones by Hyp 2. Only
finitely many covering refinements are used.

Given an element of \(L(G)\times_{L(F)}L(H)\), represent its two entries
locally. Their images are connected by such a zigzag. The cone replaces
them by representatives at a common object with equal images in \(F\),
hence by an element of \(G\times_FH\). This proves the pullback
comparison is epic for arbitrary presheaves; products use the same
argument without the equality zigzag. A tree limit is obtained by
deleting leaves and successively taking pullbacks: a leaf which is a
target contributes no new variable, and a leaf which is a source gives a
pullback over the adjacent vertex. The comparison factors through a
pullback comparison and a base change of the earlier epimorphism.
Products treat components; Hyp 1 treats the empty forest. The converse
uses the empty diagram, the discrete pair, and the cospan. The argument
uses the local meaning of epimorphisms, not enough points. \(\square\)

\textbf{Full condition.} All the finite comparisons are epic if and only
if \(L\) is left exact (flatness of \(\varphi\)).

\textbf{Proof.} Let \(M=\lim_JF\) and \(c:L(M)\to N=\lim_JL(F)\). Add
two cone vertices to \(J\), both carrying \(M\) and its canonical cone.
The enlarged diagram has limit \(M\times_M M=M\). After \(L\) its limit
is \(L(M)\times_NL(M)\), so its comparison is the diagonal of \(c\).
That diagonal is epic by hypothesis and monic automatically, hence
invertible. Therefore \(c\) is monic and, by the original hypothesis,
epic; it is an isomorphism. This also treats empty \(J\). The converse
is immediate. \(\square\)

This independently rechecks SA10-R151a-15, -16, -17, -19, -21 and -22.
The single dot on 151/015 is not the empty diagram: it always yields an
identity comparison. For Hyp 1 it must be the empty diagram, as in
SA10-K151a-05.

\subsection{B5.2. A finite
counterexample}\label{b5.2.-a-finite-counterexample}

\textbf{Status: Refuted, corrected by B5.1.} Hyp 1, 2 and \(2'\) neither
imply full (Hyp) nor suffice for homotopy descent over a point.

\textbf{Proof.} Take \(A=BC_2\) and \(\varphi(*)=\{*\}\). Then \(L\)
takes a \(C_2\)-set to its orbit set. Products compare surjectively by
choosing representatives. For a pullback of equivariant maps,
representatives of two orbits with the same image orbit can be adjusted
by a group element so their images agree; therefore all pullback
comparisons are surjective. On the regular \(C_2\)-set, however, the
identity and translation by the nonidentity element have empty
equalizer, while their maps on the one-point orbit set coincide. Its
equalizer comparison is \(\varnothing\to\{*\}\). Finally,
\(\operatorname{hocolim}_{BC_2}\{*\}=BC_2\). Its normalized simplicial
chain complex over \(\mathbf F_2\) has one generator in degrees 1 and 2
and zero boundary between them, so \(H_1=\mathbf F_2\) and it is not
contractible. The finite program verifies all four cospans, eight
completing squares, the failed equalizer, and the homology obstruction.
\(\square\)

This independently rechecks SA10-R151a-20 and SA10-K151a-04; the example
was already in the supplied material.

\subsection{B5.3. The symmetric group nerve: distinguish three
realizations}\label{b5.3.-the-symmetric-group-nerve-distinguish-three-realizations}

\textbf{Status: Proved, with a correction to ambiguous uses of
realization.} Let \(\mathsf F\) be nonempty finite sets and all maps,
\(A=\mathsf F^{\mathrm{op}}\), and \(G\) a discrete group. Put
\(U(S)=G^S/G\) for simultaneous left multiplication. Hyp 1, 2 and \(2'\)
hold, and

\[\operatorname{hocolim}_{\mathsf F^{\mathrm{op}}}U\simeq BG.\]

For \(G\ne1\), full (Hyp) fails. The ordinary simplicial restriction
also realizes to \(BG\). In contrast, the literal symmetric coend
realization with affine simplices is contractible for finite \(G\).

\textbf{Proof.} We have \(U(1)=*\). Concatenation of representatives
proves the product condition. A cospan in \(A\) is opposite to
\(P\leftarrow R\to Q\) in \(\mathsf F\), whose pushout represents its
pullback in \(A\). If tuple orbits on \(P,Q\) have the same restricted
orbit on \(R\), translate one representative so the restrictions agree.
They descend to a tuple on the pushout. This proves the condition for
every cospan, including noninjective maps. For the two coordinate maps
\(2\rightrightarrows1\) in \(A\), the equalizer is \(1\to2\). Its image
under \(U\) is \(*\to G\), but the equalizer after \(U\) is all of
\(G\).

Put \(T(S)=G^S\). Its category of elements is cofiltered: it is
nonempty, concatenation gives common predecessors, and parallel tuple
maps are equalized by quotienting their finite coordinate set by the
required identifications. The tuple values agree on those
identifications precisely because the two arrows have the same target
tuple. The quotient remains nonempty. A cofiltered category has
contractible nerve: give a common predecessor to the vertices of any
finite subcomplex, then successively equalize the finitely many edge
compatibility equations to obtain a cone. This contracts maps from
finite simplicial spheres.

The \(G\)-action on this element category is free on objects, and on
nerve simplices, since tuples are nonempty. Choosing a representative of
the initial tuple uniquely lifts a whole chain. Thus
\(N(E_U)=N(E_T)/G\), a free contractible simplicial \(G\)-space modulo
\(G\), hence \(BG\). The category-of-elements nerve is the simplicial
replacement computing a homotopy colimit. Restricting to
\(\Delta^{\mathrm{op}}\) gives the homogeneous bar construction
\(G^{n+1}/G\).

For finite \(G\), the symmetric realization of \(T\) is instead the
simplex \(\Delta(G)\): equal labels merge their barycentric weights,
leaving a probability vector on \(G\). Coends commute with orbit
quotients, so the symmetric realization of \(U\) is \(\Delta(G)/G\). The
affine contraction to the invariant barycenter is equivariant, hence
contracts the quotient. For \(C_2\) this is an interval modulo
reflection. \(\square\)

The program verifies 16 product comparisons (input cardinalities at most
4), 1,528 cospans (all three cardinalities at most 3), and 6,871
compatible target pairs for \(C_2\). The pushout proof, not the bounded
computation, proves the all-cardinality statement. The earlier statement
``the symmetric nerve realizes to \(BG\)'' is correct for homotopy
colimit or ordinary simplicial realization; it is false for the naive
symmetric coend.

\subsection{B5.4. Classification of flat symmetric
systems}\label{b5.4.-classification-of-flat-symmetric-systems}

\textbf{Status: Proved.} Flat functors
\(\mathsf F^{\mathrm{op}}\to\mathcal X\) are exactly the systems
\(S\mapsto E^S\) with \(E\to1_{\mathcal X}\) epic. In sets, \(E\) is a
nonempty set.

\textbf{Proof.} Let \(E=\varphi(1)\). The object \(n\) is the product of
\(n\) copies of \(1\) in \(A\). Yoneda and left exactness of \(L\) give
the canonical isomorphism \(\varphi(n)\cong E^n\), naturally for all
coordinate maps. Hyp 1 gives an epimorphism \(\coprod_{n\ge1}E^n\to1\),
which factors through \(E\to1\), proving inhabitedness. Conversely, the
internal tuple category of an inhabited \(E\) is locally nonempty, has
common predecessors by concatenation, and equalizes parallel arrows by
the coordinate quotient used above. It is locally cofiltered. The coend
is internally a filtered colimit of evaluations: finitely many
representatives can be moved to one index, and finitely many equalities
made simultaneously true at a further index. This proves both
surjectivity and injectivity of every finite-limit comparison. All
arguments use finitely many local refinements and apply to any
Grothendieck topos. \(\square\)

This independently proves SA10-R151a-29 without relying on its unchecked
textbook locator.

\subsection{B5.5. Flatness, forest comparisons, and Verdier
hypercovers}\label{b5.5.-flatness-forest-comparisons-and-verdier-hypercovers}

\textbf{Status: Refuted for the proposed ``strictly between''
description.} A hypercover condition agreeing with Verdier's cannot in
general lie between the manuscript's forest conditions and flatness.

For \(A=\Delta^{\mathrm{op}}\), the ordinary Cech nerve
\(U_n=\{0,1\}^{n+1}\) of a cover of a point is a hypercover, and its
realization is the contractible nerve of the indiscrete groupoid on two
objects. Yet Hyp \(2'\) fails. Take in \(\Delta\) the span
\([0]\leftarrow[1]\to[2]\), with the first map constant and the second
selecting endpoints 0,2. The corresponding fiber product contains
\((0,(0,1,0))\). A cocone must have a monotone map \([2]\to[m]\) with
equal endpoint values, hence a constant map. Every tuple lifted from a
completing cone is therefore constant in those three positions. The
displayed element has no lift. In particular full (Hyp) fails as well.

Conversely, \(\Delta[1]\) is a flat functor on \(\Delta^{\mathrm{op}}\):
its extension is evaluation at \([1]\). It is not a hypercover of a
point, because its degree-one matching map omits the ordered pair of
vertices \((1,0)\). Thus even flatness and Verdier's condition are
incomparable for ordinary simplicial indexing.

The failure of ``between'' persists for symmetric indexing. Put
\(U(S)=\operatorname{Cont}(\Delta(S),\mathbf R)\), with maps by affine
precomposition. Every continuous boundary map to \(\mathbf R\) extends
to the simplex by coning from the barycenter and interpolating to zero;
compactness gives continuity at the barycenter. Hence the ordinary
simplicial restriction is a hypercover of a point. For the finite-set
span \(1\leftarrow2\to3\), with the right map selecting vertices 0,2,
take zero on \(\Delta^0\) and

\[b(t_0,t_1,t_2)=t_1(t_0-t_2).\]

They agree on edge 02. Any amalgamating finite-set map identifies
vertices 0,2, so its affine map takes \((1/2,1/2,0)\) and
\((0,1/2,1/2)\) to the same point. A lifted function must have equal
values there; \(b\) has values \(1/4,-1/4\). Thus Hyp \(2'\) fails.
Exact rational computation checks this obstruction on 30 amalgams; the
argument proves it for every target cardinality.

There is nevertheless a one-way implication specific to symmetric
indexing: full (Hyp) implies the hypercover condition. By B5.4 its system
is \(U(S)=E^S\) with \(E\to1\) epic. The simplicial restriction is the
Cech nerve of this epimorphism; every positive matching map is an
isomorphism. The symmetric continuous-simplex example just given proves
the converse false. Thus ordinary flatness and hypercovers are
incomparable, whereas symmetric flatness is strictly stronger than
hypercovers; in either case hypercovers need not satisfy the forest
conditions.

Over sets there is also an explicit ordinary-index classification: flat
functors \(\Delta^{\mathrm{op}}\to\mathbf{Set}\) are precisely nerves of
nonempty linear orders. Indeed a flat functor is a filtered colimit of
the covariant representables
\(\operatorname{Hom}_{\Delta^{\mathrm{op}}}([n],-)=\Delta[n]\). A
filtered colimit of finite ordinals and monotone maps is a linear order:
comparisons and antisymmetry are witnessed at a common stage. Nerves
commute with filtered colimits, since finite strings use finite limits
of sets. Conversely a linear order is the filtered union of its nonempty
finite suborders. The filtered-colimit characterization follows from the
cofiltered element-category argument in B5.4. \(\square\)

This corrects \path{study/readings/151_001-017_INTEGRATION.md} lines
61--68, 78--83, and \path{study/readings/151_001-017.md} lines
303--309. The assertion that the comparison scheme lies strictly between
forest conditions and flatness is not a valid requirement for the
desired Verdier comparison.

\subsection{B5.6. Exact general homotopy-colimit
criterion}\label{b5.6.-exact-general-homotopy-colimit-criterion}

\textbf{Status: Proved under the stated interpretation.} Let
\(\mathscr X\) denote either the associated 1-localic infinity-topos or
its hypercompletion; fix this choice in each formula. Regard \(U_a\) as
0-truncated objects of \(\mathscr X\). Form

\[B_\varphi{}_n=\coprod_{a_0\to\cdots\to a_n}U_{a_0},\qquad
H_\varphi=|B_\varphi|=\operatorname{colim}_{a\in A}U_a.\]

The first face uses \(U_{a_0}\to U_{a_1}\), inner faces compose arrows,
the last face drops the last arrow, and degeneracies insert identities.
The augmentation \(\epsilon:H_\varphi\to1\) is the terminal comparison
of the derived colimit extension
\(\mathbf L\varphi_!:\mathcal P(A)\to\mathscr X\), since
\(1_{\mathcal P(A)}=\operatorname{colim}_{a\in A}y(a)\).

There is a canonical map

\[\operatorname{colim}_{a\in A}\Pi_\infty(\mathscr X/U_a)\longrightarrow\Pi_\infty(\mathscr X).\tag{B5.6.1}\]

It is an equivalence of pro-spaces if and only if for every space \(K\)
the canonical map

\[\operatorname{Map}_{\mathscr X}(1,\Delta K)\longrightarrow
\lim_{a\in A^{\mathrm{op}}}\operatorname{Map}_{\mathscr X}(U_a,\Delta K)\tag{B5.6.2}\]

is an equivalence. It suffices that \(\epsilon\) be an equivalence. Full
ordinary (Hyp) suffices for this stronger assertion even before
hypercompletion; the forest conditions do not, by B5.2. Full (Hyp) is
not necessary, by B5.5.

\textbf{Proof.} A constant object in the slice over \(U\) is
\(U\times\Delta K\to U\), and therefore

\[\operatorname{Map}_{\operatorname{Pro}(\mathcal S)}(\Pi_\infty(\mathscr X/U),K)
=\operatorname{Map}_{\mathscr X}(U,\Delta K).\]

Maps from the colimit of slice shapes to \(K\) give the limit on the
right of (B5.6.2), which is
\(\operatorname{Map}_{\mathscr X}(H_\varphi,\Delta K)\). Shape is the
functor \(K\mapsto\Gamma\Delta K\), and pro-spaces are fully faithfully
represented by these functors. This proves the exact criterion and its
sufficient condition.

Under full (Hyp), B5.1 makes \(L\) left exact; it is cocontinuous and
has right adjoint \(V\mapsto[a\mapsto\operatorname{Hom}(U_a,V)]\). Thus
it is the inverse image of a geometric morphism
\(\mathcal X\to\widehat A\). The universal property of the associated
1-localic infinity-topos (HTT Proposition 6.4.5.7
\cite[Proposition 6.4.5.7]{B56HTT}) first extends this to
\(\mathcal X_\infty\to\mathcal P(A)\). Its inverse image agrees with
\(\varphi\) on representables and hence is the derived colimit extension
into \(\mathcal X_\infty\). It preserves the terminal object. If
\(\mathscr X=\widehat{\mathcal X_\infty}\), compose this inverse image
with hypercompletion, which preserves colimits and finite limits. This
proves \(H_\varphi=1\) for either choice without asserting that the
hypercompletion is itself 1-localic. \(\square\)

In the hypercomplete choice, \(H_\varphi=1\) is equivalent to the
augmented simplicial replacement being a local weak equivalence. After
local fibrant replacement this can be tested by local relative homotopy
fillings of all finite simplicial boundaries (Dugger--Hollander--Isaksen
Proposition 3.1 \cite[Propositions 3.1--3.2]{B56DHI2004}); strict local
fillings characterize a local acyclic fibration (Proposition 3.2). These
are higher coherence conditions, not the ordinary forest comparisons.

\textbf{Exact Artin--Mazur interpretation.} Require \(\mathcal X\)
locally connected, and let
\(q:\operatorname{Pro}(\mathcal S)\to
\operatorname{Pro}(\operatorname{Ho}\mathcal S)\) be the canonical
functor. The hypercover pro-homotopy type is
\[
\operatorname{AM}(\mathcal X)\cong
q\bigl(\Pi_\infty(\widehat{\mathcal X_\infty})\bigr).
\tag{B5.6.3}
\]
Indeed Barnea--Harpaz--Horel Theorem 6.0.4, printed pp.~625--626,
identifies the derived topological realization with the shape of the
\emph{hypercompleted} infinity-topos \cite{B56BHH2017}.
Barnea--Schlank Definition 8.2, Proposition 8.4 and Remark 8.5(2),
printed pp.~62--65 of the opened preprint, identify its image under
\(q\) with the Artin--Mazur construction for every locally connected
topos \cite{B56BarneaSchlank2015}. Their Remark 8.5(3) separately records
the further localization used by Artin--Mazur; applying that localization
preserves all isomorphisms proved here.

Thus (B5.6.1), formed in pro-spaces using the hypercomplete choice and
then sent through \(q\), is the precise single-diagram comparison with
\(\operatorname{AM}(\mathcal X)\). It is an isomorphism in
\(\operatorname{Pro}(\operatorname{Ho}\mathcal S)\) if and only if, for
every space \(K\), the map
\[
\pi_0\operatorname{Map}_{\widehat{\mathcal X_\infty}}(1,\Delta K)
\longrightarrow
\pi_0\!\left(\lim_{a\in A^{\mathrm{op}}}
\operatorname{Map}_{\widehat{\mathcal X_\infty}}(U_a,\Delta K)\right)
\tag{B5.6.4}
\]
is a bijection. To prove this, a pro-space \(P=(P_i)\) satisfies
\(\operatorname{Hom}(qP,K)=\operatorname{colim}_i[P_i,K]
=\pi_0\operatorname{Map}(P,K)\), since filtered colimits commute with
\(\pi_0\). Pro-objects of \(\operatorname{Ho}\mathcal S\) are fully
faithfully detected by their set-valued functors of maps to constant
objects. Apply this to the already formed comparison (B5.6.1).
The limit in (B5.6.4) is a limit of spaces taken \emph{before} \(\pi_0\).
We neither interchange these operations nor claim that \(q\) preserves
the colimit in (B5.6.1). The mapping-space equivalences (B5.6.2) imply
(B5.6.4), but the latter alone is not asserted to recover them.

For the full, possibly nonhypercomplete shape, use (B5.6.2) with
\(\mathscr X=\mathcal X_\infty\). A local weak equivalence or a
hypercover alone supplies the hypercomplete version, and does not by
itself justify that full-shape conclusion. It also suffices for the
protruncated comparison, because truncated constant objects are
hypercomplete. Full ordinary (Hyp) supplies the stronger terminal-object
descent before hypercompletion, as proved above. These are distinct
targets, rather than implicit identifications of pro-homotopy objects
and pro-spaces; see also HTT Example 7.1.6.9 \cite{B56HTT}.

\subsection{B5.7. Verdier matching as comparisons of
representables}\label{b5.7.-verdier-matching-as-comparisons-of-representables}

\textbf{Status: Known, with the precise derivation given.} Define

\[M_nU=\lim_{\varnothing\ne S\subsetneq\{0,\ldots,n\}}U_{|S|-1},\qquad M_0U=1,\]

with face restriction arrows. The conditions \(U_n\to M_nU\) epic for
all \(n\ge0\) are the hypercover conditions. The realization is
infinity-connected in \(\mathcal X_\infty\) and terminal after
hypercompletion. This is HTT Definition 6.5.3.2, Lemma 6.5.3.11 and
Theorem 6.5.3.12 \cite{B56HTT} (or DHI Definition 4.2 and Proposition
3.2).

For \(A=\Delta^{\mathrm{op}}\) and \(n\ge2\), the limit in
\(\widehat A\) of the representables of all proper nonempty faces of
\([n]\) is \(y[n]\). Consequently its comparison \(c_F\) is precisely
\(U_n\to M_nU\). Degrees zero and one must be added separately. For
\(A=\mathsf F^{\mathrm{op}}\) the same statement holds for all
\(n\ge1\), so all positive-degree matching conditions are the
corresponding boundary comparisons, with degree zero \(U(1)\to1\).

\textbf{Proof of the representable identity.} Compatible maps from all
proper faces into an ordinal determine the vertex images. For \(n\ge2\),
every pair of vertices lies in a proper face; the images are therefore
monotone, determining a unique map from \([n]\). At \(n=1\) the two
vertex images need not be monotone. For finite sets, no monotonicity
restriction occurs, and compatible vertex maps determine a unique
function also in the two-vertex case. Yoneda turns these colimit
properties into the asserted limits. \(\square\)

The inclusion \(j:\Delta^{\mathrm{op}}\to\mathsf F^{\mathrm{op}}\) is
homotopy final. For any nonempty \(S\), the comma category
\((S\downarrow j)\) is the category of nonempty ordered tuples in \(S\)
and monotone restriction maps. Choose \(s\in S\). Prepending \(s\)
defines an endofunctor with natural transformations to the identity
(drop the first entry) and to the constant tuple \((s)\) (select the
first entry). Its nerve is contractible. The cofinality criterion HTT
Theorem 4.1.3.1 applies. Thus symmetric matching conditions compute the
homotopy colimit over all finite-set maps, not only the ordinary
simplicial subcategory.

\subsection{B5.8. Every matching degree is independently
needed}\label{b5.8.-every-matching-degree-is-independently-needed}

\textbf{Status: Proved.} For every \(m\ge0\) there is a simplicial set
with all matching maps except the \(m\)th surjective and with
noncontractible realization. The example extends to symmetric indexing
by all nonempty finite sets.

\textbf{Proof.} For \(m=0\), take the empty simplicial set: every
positive matching target is empty. For \(m=1\), take the constant
two-point set: only the diagonal map to its square fails, since all
higher simplex boundaries are connected.

For \(m\ge2\), let \(n=m-1\) and define
\(K_n(k)=Z^n(\Delta^k;\mathbf F_2)\) with cochain pullback maps. A
simplicial map \(T\to K_n\) is precisely a normalized \(n\)-cocycle on
\(T\): its values on \(n\)-simplices form the cochain and the
\((n+1)\)-simplex relation is the cocycle equation. Hence the \(k\)th
matching map is

\[Z^n(\Delta^k;\mathbf F_2)\to Z^n(\partial\Delta^k;\mathbf F_2).\]

For \(k<n\) both groups vanish; for \(k=n\) the target vanishes. At
\(k=n+1\), the target consists of arbitrary values on the \(n+2\)
boundary \(n\)-faces; the image is the codimension-one subspace whose
sum is zero. For \(k>n+1\), every \(n\)- and \((n+1)\)-simplex already
belongs to the boundary, so restriction is an isomorphism. Thus only
degree \(m\) fails.

The normalized Moore complex of \(K_n\) is \(\mathbf F_2\) in degree
\(n\) and zero otherwise: a normalized element vanishes on all specified
faces, and in degree \(n+1\) the cocycle equation forces the remaining
face to vanish; above that degree all \(n\)-faces already lie in
vanishing faces. Its homotopy group \(\pi_n\) is therefore
\(\mathbf F_2\), by Curtis, Theorem 3.7 (printed p.~131)
\cite{B56Curtis1971}. This is also the cocycle model of the Dold--Kan
object of \(\mathbf F_2[n]\); normalization is the equivalence of Stacks
Theorem 14.24.3 (019G) \cite{B56StacksDK}. For the symmetric extension,
define the same cocycles on the abstract simplex on a finite set; an
arbitrary vertex map induces a normalized simplicial chain map, sending
a face with repeated image vertices to zero. Over \(\mathbf F_2\) this
construction is independent of ordering. It gives the same ordinary
simplicial restriction, hence the same homotopy colimit by B5.7. The
finite code checks 32 matching comparisons for \(1\le n\le4\),
\(0\le k\le7\), with precisely the stated failures. \(\square\)

\subsection{B5.9. Component sets alone cannot reconstruct the
homotopy
type}\label{b5.9.-component-sets-alone-cannot-reconstruct-the-homotopy-type}

\textbf{Status: Refuted, corrected by B5.6 and the hypercover
pro-system.} No condition requiring even all ordinary comparisons to be
isomorphisms can ensure \(\operatorname{hocolim}_A\pi_0(U_a)\) models
the Artin--Mazur type without a condition on the pieces.

\textbf{Proof.} Take \(A\) terminal,
\(\mathcal X=\operatorname{Sh}(S^2)\) and \(U_*=1_{\mathcal X}\). Then
\(L:\mathbf{Set}\to\operatorname{Sh}(S^2)\) is constant-sheaf inverse
image, so it preserves finite limits. But \(\pi_0(U_*)=*\), and its
homotopy colimit is a point; the sphere has
\(H^2(S^2;\mathbf Z)=\mathbf Z\). The correct single-diagram formula
uses slice shapes. Component sets can replace them when each canonical
map \(\Pi_\infty(\mathscr X/U_a)\to\pi_0(U_a)\) is an equivalence and
the diagram is homotopy effective: replace each object by that
equivalent pro-space in B5.6. In general Artin--Mazur retains the
cofiltered pro-system of all hypercovers and their component simplicial
sets. Even the constant terminal hypercover is not generally a single
model of the shape. \(\square\)

\subsection{B5.10. An explicit all-index local filling
criterion}\label{b5.10.-an-explicit-all-index-local-filling-criterion}

\textbf{Status: Known criterion, derived here in the required form.} The
exact stronger condition \(H_\varphi=1\) in the hypercomplete
infinity-topos has an explicit test using the original \(U_a\), finite
limits, coproducts and a sequential colimit. It does not require one to
guess boundary objects in the arbitrary category \(A\).

Let \(N=B_\varphi\) be the simplicial replacement from B5.6;
equivalently it is the nerve of the internal element category, whose
objects are \(\coprod_aU_a\) and whose arrow object is
\(\coprod_{a\to b}U_a\). For a finite simplicial set \(K\), let
\([K,N]\) be the object of simplicial maps, expressed as a finite limit
using the nondegenerate simplices of \(K\). Put

\[R_n=\operatorname*{colim}_{r\ge0}[\operatorname{Sd}^{r}\Delta^n,N],\]

with the transition maps induced by the standard last-vertex subdivision
maps. Here colimits and finite limits are in \(\mathcal X\), and
\(\operatorname{Sd}\) is barycentric subdivision. Then the following are
equivalent:

\begin{enumerate}
\def\labelenumi{\arabic{enumi}.}
\tightlist
\item
  \(H_\varphi\to1\) is an equivalence after hypercompletion.
\item
  The simplicial replacement \(N\to1\) is a local weak equivalence.
\item
  For every \(n\ge0\), the following explicit map is an epimorphism in
  \(\mathcal X\):
\end{enumerate}

\[\operatorname*{colim}_{r\ge0}[\operatorname{Sd}^{r}\Delta^n,N]
\longrightarrow
\operatorname*{colim}_{r\ge0}[\operatorname{Sd}^{r}\partial\Delta^n,N].\tag{B5.10.1}\]

For \(n=0\) the right side is terminal. Since filtered colimits commute
with finite limits in a Grothendieck topos, the right side is precisely
\(M_nR\). Thus (3) is an ordinary epimorphic matching-map scheme on the
explicitly derived simplicial object \(R\). It says that every finite
boundary admits a filling locally after sufficiently many subdivisions.
The required subdivision number may vary over the covering family; no
assertion of algorithmic decidability is made.

\textbf{Proof.} Present \(\mathcal X\) by a small site, and regard \(N\)
as a simplicial sheaf there. Apply \(\operatorname{Ex}^{\infty}\)
sectionwise to obtain a simplicial presheaf \(F\). For every object of
the site, \(F\) is a Kan complex and \(N\to F\) is a weak homotopy
equivalence, by Kerodon Propositions 3.3.6.7 and 3.3.6.9 (Tags 00ZH and
00ZK) \cite{B56KerodonExInfinity}. Therefore \(N\to F\) is a local weak
equivalence, and \(F\) is locally fibrant. Dugger--Hollander--Isaksen
Proposition 3.2, reformulated as their equation (3.4) \cite{B56DHI2004},
says that \(F\to1\) is a local weak equivalence if and only if every map
\(F_n\to[\partial\Delta^n,F]\) is a local epimorphism: its local
fibration hypothesis is automatic here. Sheafification preserves finite
limits and colimits. Sheafifying these maps gives precisely (B5.10.1),
by the subdivision--extension adjunction and finite presentation of the
boundaries. This proves (2) equivalent to (3), independently of the
chosen site. The local model category presents the hypercompletion, so
(1) is equivalent to (2). This last restriction matters: local weak
equivalence gives an infinity-connected augmentation in a
nonhypercomplete topos, not necessarily an equivalence there.
\(\square\)

This supplies an effective mathematical description for every small
\(A\), in addition to the direct sufficient condition on the
manuscript's \(c_F\) in B5.1 and the sharper boundary schemes of B5.7.
The weaker forest scheme fails by B5.2--B5.3. Omitting any dimension
from the strict matching scheme fails by B5.8, in ordinary and symmetric
indexing. Conditions for shape equivalence alone can be weaker than
universal local descent; their precise criteria are B5.6.2 for pro-spaces
and B5.6.4 for the original pro-homotopy category. In particular,
(B5.10.1) is necessary and sufficient for the stronger local descent
assertion, not a claimed necessary condition for equality only of shapes.

\subsection{Scope of the comparison
result}\label{scope-of-the-comparison-result}

The answer separates the literal ordinary comparisons, the universal
local descent condition, and the Artin--Mazur target. All finite
ordinary comparisons epic is exactly flatness and suffices even before
hypercompletion; the forest subcollection fails. The boundary
comparisons in ordinary and symmetric indexing give Verdier hypercovers,
with every matching degree indispensable. For arbitrary indexing,
(B5.10.1) is an explicit necessary-and-sufficient higher local descent
criterion after hypercompletion. For equality only of shapes, (B5.6.2)
is the exact criterion; for the original Artin--Mazur pro-homotopy type,
(B5.6.3)--(B5.6.4) give the exact locally connected interpretation.
A universal minimal subset of
the ordinary finite diagram shapes for every possible \(A\) is not
asserted or needed for these results.

\section{B.6. Finite-set presentations}\label{b.6.-presentations-of-the-category-of-nonempty-finite-sets}

Composition is written right to left. Let \(\mathsf F\) have skeleton
\(n=\{1,\ldots,n\}\), \(n\ge1\). The manuscript generators in 151/017
are the permutations \(\mathfrak S_n\), the inclusion
\(\delta_n:n\to n+1\) of the first \(n\) elements, and \(K_n:n+1\to n\)
which fixes \(1,\ldots,n\) and sends \(n+1\) to \(n\). Let \(i_n\)
extend a permutation by fixing the last element, and let
\(t_n=(n\ n+1)\in\mathfrak S_{n+1}\).

\subsection{B6.1. The proposed global presentation is
incomplete}\label{b6.1.-the-proposed-global-presentation-is-incomplete}

\textbf{Status: Refuted, with an explicit finite witness extending to
all degrees.} The four displayed relation families of 151/017 do not
present \(\mathsf F\). Adding both global invariance families boxed in
151/005, even generalized to all degrees, still does not suffice.

The four families, with the inclusions made explicit to fix types, are

\[K_n\delta_n=\mathrm{id},\qquad
\delta_n\sigma=i_n(\sigma)\delta_n\quad(\sigma\in\mathfrak S_n),\]
\[i_{n-1}(\sigma)K_n=K_n i_n i_{n-1}(\sigma)\quad(\sigma\in\mathfrak S_{n-1}),\qquad K_n t_n=K_n.\]

The additional families are

\[(n+1\ n+2)\delta_{n+1}\delta_n=\delta_{n+1}\delta_n,
\qquad K_nK_{n+1}(n\ n+1)=K_nK_{n+1}.\]

\textbf{Countermodel.} Interpret these contravariantly in sets. Put
\(E_1=\{0\}\) and \(E_n=\{0,1\}\) for \(n\ge2\), and let every
permutation act trivially. Put \(p_1=\delta_1^*:E_2\to E_1\) equal to
the unique map and \(d_1=K_1^*:E_1\to E_2\) equal to the inclusion of
\(0\). For every \(n\ge2\), put \(p_n=d_n=\mathrm{id}_{\{0,1\}}\). The
retraction equation is \(p_nd_n=\mathrm{id}_{E_n}\), which holds. All
equivariance and invariance equations hold because every permutation
acts trivially.

In actual finite sets, however, one has the typed equation between maps
\(2\to2\)

\[\boxed{(1\ 2)K_2(1\ 3)\delta_2=\delta_1K_1.}\tag{B6.1.1}\]

Both sides are constant with value \(1\): the middle composite on the
left sends both elements to \(2\), and the leading transposition sends
that value to \(1\). In the contravariant model, the left side is
\(p_2d_2=\mathrm{id}_{E_2}\) and the right side is \(d_1p_1\), constant
with value \(0\). They differ at \(1\). Thus (B6.1.1) does not follow
from the proposed presentation, even after both invariance families are
added. Only the three objects \(1,2,3\) are needed to witness the
failure; the model is defined in every degree. The program checks 412
instances of the stated relations through cardinality six and the
displayed missing relation. The all-degree proof is the preceding
definition, not that finite bound. \(\square\)

Visual inspection of original 151/017 (PDF page 17, printed p.~16)
reads the two lower bounds as \(n\ge2\) and \(n\ge3\), respectively.
The alternative bounds in the supplied transcription cannot fix the failure: the omitted smallest symmetric groups are trivial, and
the countermodel satisfies either reading. The margin's
\(p_{n+1}|\rho_{n+1}\) denotes the idempotent \(\delta_nK_n\). Reading
it as a permutation would be mathematically false, since this map has
image of size \(n\).

\subsection{B6.2. A complete replacement is available in the primary
literature}\label{b6.2.-a-complete-replacement-is-available-in-the-primary-literature}

\textbf{Status: Known, with a direct translation to the manuscript
generators.} Grandis, \emph{Finite sets and symmetric simplicial sets},
§3 and Theorem 4.2 \cite{B56Grandis2001}, gives a complete presentation
of finite cardinals by all elementary cofaces, codegeneracies and
adjacent transpositions, subject to the cosimplicial, symmetric-group
and mixed interchange relations. His §5 restricts the presentation to
nonempty finite cardinals. Thus a complete corrected presentation is
obtained by adding all these translated relations to the manuscript
generators.

Here are explicit substitution words, including their arities.  Write
\(c_i^{(m)}=(i\ i+1\ \cdots\ m)\in\mathfrak S_m\), with
\(c_m^{(m)}=1\).  For the coface \(d_i^{(n)}:n\to n+1\) omitting
\(i\), and the codegeneracy \(e_i^{(n)}:n+1\to n\) identifying
\(i,i+1\), set
\[
 d_i^{(n)}=c_i^{(n+1)}\delta_n\quad(1\le i\le n+1),\qquad
 e_i^{(n)}=c_i^{(n)}K_n\bigl(c_i^{(n+1)}\bigr)^{-2}
       \quad(1\le i\le n).
 \tag{B6.2.1}
\]
Adjacent transpositions are \(s_i=(i\ i+1)\), already present in the
symmetric groups.  The first word fixes positions before \(i\) and
moves every later old position one step to the right.  In the second
word, the source permutation fixes \(k<i\), sends \(i,i+1\) to
\(n,n+1\), and sends \(k>i+1\) to \(k-2\); \(K_n\) merges the last
two, and the output cycle puts the merged value back in position
\(i\).  Thus these are exactly the indicated finite-set maps.
Crucially, the last words recover the old generators literally:
\(d_{n+1}^{(n)}=\delta_n\) and \(e_n^{(n)}=K_n\).

For a completely explicit enlarged relation scheme, use (B6.2.1) in
all the following formulas, whenever their arities are positive and
well typed; dimension superscripts are suppressed:
\[
 d_i d_j=d_{j+1}d_i\ (i\le j),\qquad
 e_j e_i=e_i e_{j+1}\ (i\le j),
\]
\[
 e_jd_i=
 \begin{cases}
 d_i e_{j-1},&i<j,\\
 1,&i=j\text{ or }i=j+1,\\
 d_{i-1}e_j,&i>j+1;
 \end{cases}
\]
\[
 s_i^2=1,\quad s_i s_{i+1}s_i=s_{i+1}s_i s_{i+1},\quad
 s_i s_j=s_j s_i\ (i<j-1),
\]
\[
 \begin{array}{lll}
 s_i d_j=d_j s_i &(i<j-1),&s_i e_j=e_j s_i\quad(i<j-1),\\
 s_i d_i=d_{i+1},&&s_i e_i=e_{i+1}s_i s_{i+1},\\
 s_i d_j=d_j s_{i-1}&(i>j),&s_i e_j=e_j s_{i+1}\quad(i>j),
 \end{array}
 \qquad e_i s_i=e_i.
 \tag{B6.2.2}
\]
These are Grandis's (3), (10), (12), (13), with indices shifted from
zero-based to one-based notation.  The omitted mixed cases follow by
\(s_i^2=1\), exactly as in Grandis.  The group equations are redundant
if the complete groups \(\mathfrak S_n\) and their multiplication laws
are retained as the manuscript prescribes.

We verify carefully that no unspoken defining equations are missing.
Let \(P\) be Grandis's presented category on positive cardinals and
let \(Q\) be the category on the manuscript generators with their group
laws, the valid manuscript equations, and all of (B6.2.2) after the
substitutions (B6.2.1).  Substitution defines a functor
\(F:P\to Q\), because every defining relation of \(P\) has explicitly
been imposed in \(Q\).  This functor is surjective on every hom-set:
its image contains \(\delta_n\) and \(K_n\) literally and contains all
adjacent transpositions, which generate each of the already imposed
symmetric groups.  The finite-set interpretation gives
\(G:Q\to\mathsf F\), because every imposed relation holds for these
functions.  The composite \(GF\) is the standard interpretation of
Grandis's generators and is an isomorphism by Theorem 4.2 and §5.

For faithfulness of \(G\), choose lifts \(u,v\) in \(P\) of two
parallel arrows of \(Q\).  Equality of their \(G\)-images gives
\(GF(u)=GF(v)\), hence \(u=v\), and hence equality of the arrows in
\(Q\).  Fullness follows from the same composite.  Thus
\(Q\cong\mathsf F\), with exactly the requested old generators.
This proof does not assume that generator substitutions are inverses
without checking their defining equations.  The auxiliary script
\texttt{check\_B6\_translation.py} verifies all 168 coface and
codegeneracy words for \(1\le n\le12\); the explicit coordinate
argument above proves the formulas in every arity.  We make no
minimality claim for this enlarged relation list.

Exact locator: Grandis, \emph{Theory and Applications of Categories}
\textbf{8} (2001), no. 8, 244--252; §3, printed pp.~247--248; Theorem
4.2, p.~249; §5, p.~250. The full primary PDF was opened and downloaded.
No novelty is claimed for this replacement.

\subsection{B6.3. The corrected low-degree presentation is
complete}\label{b6.3.-the-corrected-low-degree-presentation-is-complete}

\textbf{Status: Proved for the corrected 151/005 reading; refuted for
the explicitly identified alternatives.} The final relations of 151/005
present the full subcategory \(\mathsf F_{\le3}^{\mathrm{op}}\),
provided the last subscript is \(1\). They do not do so if that
subscript is \(3\), or if the erroneous relation on 151/003 is retained
uncorrected.

Write the three objects \(E_0,E_1,E_2\) for cardinalities \(1,2,3\),
respectively, and abbreviate

\[p=p_0:E_1\to E_0,\quad d=\delta_0:E_0\to E_1,
\quad q=p_{12}:E_2\to E_1,\quad e=\delta_1:E_1\to E_2.\]

The groups are \(\mathfrak S_2\) and \(\mathfrak S_3\). In the
coordinate interpretation, let

\[\sigma(x,y)=(y,x),\quad \rho(x,y,z)=(y,z,x),
\quad\sigma_1(x,y,z)=(x,z,y),\quad\sigma_3(x,y,z)=(y,x,z).\]

Thus \(\sigma_i\) fixes coordinate \(i\). The corrected list is

\[pd=1,\quad \sigma d=d,\quad qe=\sigma,\quad q\rho e=dp,
\quad q\rho^2e=1,\]
\[\sigma_1e=e,\quad q\sigma_3=\sigma q,\quad ged=ed\ (g\in\mathfrak S_3),
\quad pq\sigma_1=pq,\]

together with the group equations. The coordinate maps

\[p(x,y)=x,\quad d(x)=(x,x),\quad q(x,y,z)=(x,y),\quad e(x,y)=(y,x,x)\]

satisfy them. These maps and the permutations generate all maps between
sets of cardinalities at most three, by elementary inclusion and
retraction factorization. Hence the presented category \(Q\) maps
surjectively on hom-sets to \(\mathsf F_{\le3}^{\mathrm{op}}\).

We now prove injectivity by a terminating reduction and exact hom-set
bounds. First combine adjacent permutations. A path which rises from
\(E_1\) to \(E_2\) and immediately returns has a segment \(qge\). Every
\(g\in\mathfrak S_3\) is \(\rho^i h\) with
\(h\in\langle\sigma_1\rangle\). Since \(he=e\), the three relations
\(q\rho^ie\in\{\sigma,dp,1\}\) eliminate this segment. A path rising
from \(E_0\) to \(E_1\) and returning has a segment
\(p\sigma^\epsilon d=p d=1\). These replacements terminate: order words
first by the number of \(q,e\) transitions and then by the number of
\(p,d\) transitions. A first replacement lowers the first count; a
second lowers the second without increasing the first.

Every surviving path descends in height and then ascends, with at most
one valley. Paths through \(E_0\) are constant maps: on the output side
\(\sigma d=d\) and \(ged=ed\) erase all output permutations. On the
input side, there is one form from \(E_0\), at most two forms
\(p\sigma^\epsilon\) from \(E_1\), and at most three forms \(pqh\) from
\(E_2\), because \(pq\sigma_1=pq\). A path with valley \(E_1\) from
\(E_2\) to \(E_2\) is \(ge\sigma^\epsilon qh\). The stabilizer relation
\(\sigma_1e=e\) leaves three choices for the coset of \(g\); the
equation \(q\sigma_3=\sigma q\) moves the two-element left stabilizer of
\(h\) into the central \(\sigma^\epsilon\), leaving three choices for
its coset. Thus there are at most \(3\cdot2\cdot3=18\) forms of rank
two. The same calculation gives at most six rank-two maps each from
\(E_1\) to \(E_2\) or from \(E_2\) to \(E_1\). Finally there are at most
two invertible forms at \(E_1\) and six at \(E_2\).

Collecting these bounds, with source cardinality indexing rows and
target cardinality columns, gives

\[\begin{pmatrix}1&1&1\\2&4&8\\3&9&27\end{pmatrix}.\]

The entries are exactly
\(a^b=|\operatorname{Hom}_{\mathsf F^{\mathrm{op}}}(a,b)|\). Since the
functor was already surjective on each hom-set, each bound is attained
and it is bijective on each hom-set. It is the identity on objects,
hence an isomorphism of categories. The program enumerates all 56 maps,
verifies this matrix, checks the rank decomposition and every corrected
relation. The reduction proof supplies completeness, which checking the
relations in the target alone would not establish. \(\square\)

The doubtful alternatives must remain separate. On 151/005, the last
relation is transcribed
\(p_0p_{12}\sigma_{\text{unsure }1|3}=p_0p_{12}\). The \(1\) reading is
the corrected presentation just proved. The \(3\) reading says that
first-coordinate and second-coordinate projections agree, and is false
already on \((0,1,0)\). On 151/003, the transcribed \(\sigma_3e=e\) is
false: the two maps have coordinate values \((x,y,x)\) and \((y,x,x)\).
The final list on 151/005 instead has \(\sigma_1e=e\), which is true.
Retaining every literal 151/003 equation together with 151/005 therefore
gives a false presentation. These checks confirm the mathematical
correction SA10-K151a-01 and complete the earlier partial verification
SA10-R151a-08. Independent inspection of original 151/003 (PDF page 3, printed p.~2)
confirms the boxed subscript \(3\); inspection of the middle box on
151/005 (PDF page 5, printed p.~4) confirms subscript \(1\).
The final 151/005 subscript is a dense short mark and remains visually
unresolved between the two supplied readings.  Thus the mathematical
conclusions for both readings are retained; no paleographic certainty
is claimed for that final mark.

\subsection{Source and computation
record}\label{source-and-computation-record}

The companion files are \texttt{CHECK\_EVENTS\_B56.tsv},
\path{references_B56.bib}, \path{code_B56/check_B56.py},
\path{code_B56/check_B56.out},
\path{code_B56/check_matching_independence.py}, and its
\texttt{.out} file. CPU timings and complete finite outputs are in the
output files. Every bounded check is separated above from its all-size
proof. Sources consulted but not opened successfully are not used as
premises.
}{}
\part{Cote 150: deployment, presentation, and control data}
% AI-produced by GPT-6 Astra (OpenAI), 2026-11-19. Generated by code_C/build_C_sections.py.
\section{C.1. Deployment, simultaneous cones, and exit categories}\label{c.1.-deployment-simultaneous-cones-and-exit-categories}

\subsection{C1.01 --- Refuted: weak conical data do not imply fibrancy}\label{c1.01-refuted-weak-conical-data-do-not-imply-fibrancy}

Let Y=R, E=S\ensuremath{{}^{1}}, M=D\ensuremath{{}^{2}}. Embed E as \ensuremath{\partial}M, give it its usual collar, and let p:E\ensuremath{\to}Y be constant with value 0. Form X=M \ensuremath{\amalg}\_E C(p), where

C(p)=(E\ensuremath{\times}{[}0,1{]} \ensuremath{\amalg} Y)/((e,0)\ensuremath{\sim}p(e)),

and E\ensuremath{\times}1 is identified with \ensuremath{\partial}M. Thus X is a sphere attached to a line at one point. Set X\ensuremath{{}_{0}}=Y and X\ensuremath{{}_{1}}=X. The open complement of Y is R\ensuremath{{}^{2}}. The collar supplies exactly the conical neighborhood of 150/023 used inductively at /036. No local triviality of p occurs in those data. The closed family is finite, increasing, exhaustive, and satisfies the required intersection identity.

Its prescribed deployment has \ensuremath{\Sigma}\ensuremath{{}_{0}}=R, \ensuremath{\Sigma}\ensuremath{{}_{1}}=D\ensuremath{{}^{2}}, \ensuremath{\Sigma}\ensuremath{{}_{01}}=S\ensuremath{{}^{1}}, source p, and target the boundary embedding. All embedding, intersection, and triple conditions hold or are vacuous. The source is proper because S\ensuremath{{}^{1}} is compact. It is not a Hurewicz or Serre fibration: the path \ensuremath{\alpha}(t)=t in R starts at p(e)=0 but cannot lift through a constant map for t\textgreater0.

This does not refute modern conical stratifications. This two-stratum decomposition is not locally conically trivial along Y at 0. It refutes the implication from the weaker \S{}8 data. The margin of /039 explicitly restricts condition (1) to the equisingular case.

\textbf{Correct implication.} If the maps \ensuremath{\Sigma}\ensuremath{{}_{ij}}\ensuremath{\to}\ensuremath{\Sigma}\ensuremath{{}_{i}} are locally trivial bundles, or fibrations stable under the relevant base changes and composites, every initial-segment map \(\Sigma_{d\prime}\)\ensuremath{\to}\(\Sigma_d\) is a fibration. The Segal pullback identity expresses it as a succession of base changes of these source maps. This is an additional hypothesis, not a consequence of arbitrary \S{}8 conical structure.

\subsection{C1.02 --- Proved under fibrancy: the complete Segal object}\label{c1.02-proved-under-fibrancy-the-complete-segal-object}

Let C be an internal category in compactly generated spaces, graded by a finite poset I. Suppose C\ensuremath{{}_{0}}=\ensuremath{\coprod}\ensuremath{{}_{i}}\ensuremath{\Sigma}\ensuremath{{}_{i}}, C\ensuremath{{}_{1}}=\ensuremath{\coprod}\ensuremath{{}_{i}}\ensuremath{\le}\ensuremath{{}_{j}}\ensuremath{\Sigma}\ensuremath{{}_{ij}}, \ensuremath{\Sigma}\ensuremath{{}_{ii}}=\ensuremath{\Sigma}\ensuremath{{}_{i}} consists exactly of identities, and each strict source \ensuremath{\Sigma}\ensuremath{{}_{ij}}\ensuremath{\to}\ensuremath{\Sigma}\ensuremath{{}_{i}} is a Hurewicz fibration. The ordinary nerve has Segal maps that are homeomorphisms onto iterated ordinary pullbacks. A pullback along a Hurewicz fibration computes a homotopy pullback. Thus its nerve is a Segal object in the infinity-category of spaces.

The functor to N(I) takes equivalences to identities. Every equivalence therefore lies over i=i. That summand consists only of identities, and all identities are equivalences. Hence \(C_0\to N(C)_1^{\mathrm{eq}}\) is an equivalence, proving completeness.

If ``Segal space'' is defined to require Reedy fibrancy, take a Reedy fibrant replacement. The raw nerve need not be Reedy fibrant. For the identity-only category on S\ensuremath{{}^{1}} it is constant, and its first matching map is S\ensuremath{{}^{1}}\ensuremath{\to}S\ensuremath{{}^{1}}\ensuremath{\times}S\ensuremath{{}^{1}}, the diagonal. A pair of unequal paths starting at the same point cannot lift through this diagonal.

This proves completeness, not the geometric identification with Exit(X). That identification requires a separate comparison; an equivalence of realizations alone does not imply it.

\subsection{C1.03 --- Proved: the simultaneous-cone construction}\label{c1.03-proved-the-simultaneous-cone-construction}

This supplies an explicit completion of the postponed construction, rather than attributing unspecified transition maps to /043. Let \(\Sigma:\operatorname{Drap}(I)^{op}\to\mathrm{Top}\) be the flag diagram of the deployed category, including flags of length zero. For a nonempty flag \(e=\{i_0<\cdots<i_n\}\), write \(e_{\le a}=\{i\in e:i\le a\}\). A pair \((d,e)\) means \(d\subset e\setminus\{\max e\}\); the empty \(d\) is allowed. Put \(a_0<\cdots<a_p=\max e\) for the elements of \(d\cup\{\max e\}\), so \(p=|d|\), not necessarily the length of \(e\).

Define the simultaneous mapping cylinder \[
 C(d,e)=\operatorname{hocolim}\bigl(
 \Sigma_{e_{\le a_p}}\longrightarrow\cdots\longrightarrow
 \Sigma_{e_{\le a_0}}\bigr),
\] using the ordinary thick simplicial mapping-cylinder model. Explicitly it is the quotient of \[
 \coprod_{k=0}^{p}\Sigma_{e_{\le a_k}}\times\Delta^k
\] that identifies the face omitting the last vertex with its image under \(\Sigma_{e_{\le a_k}}\to\Sigma_{e_{\le a_{k-1}}}\), and makes the analogous identifications on all nested subfaces. Faces omitting a nonfinal vertex remain boundary faces. In particular \(C(\varnothing,e)=\Sigma_e\); for one cut it is the mapping cylinder of the indicated projection, called a mapping cone in the manuscript's convention.

There are three kinds of maps, generating the transition diagram.

\begin{itemize}
\tightlist
\item
  If \(d\subset d'\subset e\setminus\{\max e\}\), the inclusion of the selected vertices gives \(C(d,e)\to C(d',e)\).
\item
  For \(f=d\cup\{\max e\}\subset e\), the restriction maps \(\Sigma_{e_{\le a}}\to\Sigma_{f_{\le a}}\) give \(C(d,e)\to C(f\setminus\{\max f\},f)\).
\item
  For each selected \(a\in d\), the initial-prefix cylinder already present as a face has its canonical map \(C(d\cap e_{<a},e_{\le a})\to C(d,e)\). These maps identify internal prefix blocks with their standalone objects in the diagram.
\end{itemize}

All relations among these maps are those induced by inclusions of flags and functoriality of \(\Sigma\). This can equivalently be taken as the image diagram with the maps above, imposing every commuting square that follows from the displayed inclusions.

The colimit of this diagram is canonically homeomorphic to the ordinary realization of the strict nerve: \[
 \operatorname{colim}_{(d,e)}C(d,e)
 \;\cong\;
 \left(\coprod_{\varnothing\ne e\in\operatorname{Drap}(I)}
 \Sigma_e\times\Delta^{|e|-1}\right)/\sim .
\]

\textbf{Proof.} A continuous map out of the right-hand quotient is a family of continuous maps \(f_e:\Sigma_e\times\Delta^{|e|-1}\to Y\), compatible with every face restriction. The full block \(C(e\setminus\{\max e\},e)\) contains the final-vertex-deletion faces internally; the third transition family identifies each such prefix face with its standalone prefix block, including iterated final deletions. Its other faces correspond to \(C(d,e)\). The first transition map identifies these as the corresponding faces of the full block; the second imposes precisely the restriction to the selected flag \(d\cup\{\max e\}\). Every face is either final-vertex deletion or contains the final vertex. Thus cocones out of the cone diagram and maps out of the realization are the same, naturally in the target topological space. The quotient topologies agree by this universal property. Empty \(d\) supplies the vertex identifications; flags of one element supply the object spaces. For \(I=\varnothing\), both constructions are empty. Identities can instead be inserted to use the ordinary unital nerve; its degenerate simplices collapse and give the same realization.

\textbf{Geometric recovery.} Suppose the finite conical space is supplied with recursive compatible collars in the sense actually used in its construction: cutting a minimal stratum \(A\) gives a space \(X'\), with retained outer boundary \(H\), and \[
 X=X'\cup_H C(H\to A),
\] while the same collar systems induce the smaller deployments on \(X'\) and \(H\). Induction is performed simultaneously for spaces with any already retained outer boundary faces. Then the realization is homeomorphic to \(X\), relative to those faces.

For completeness, let \(C\) be the deployed category and \(i\) a minimal label. Let \(D\) be its restriction to the other labels. Define \(M\) by \(N_nM=\coprod_{j_0<\cdots<j_n}\Sigma_{i,j_0,\ldots,j_n}\). The face maps make \(M\) the deployment of \(H\). Every strict simplex of \(C\) either lies in \(D\) or has \(i\) as its unique first vertex. Hence its realization has a canonical homeomorphism \[
 |C|\cong |D|\cup_{|M|} C(|M|\to\Sigma_i).
\] On a simplex containing \(i\), write its barycentric coordinates as \(t_i=1-r,\ t_{j_k}=r u_k\). At \(r=0\) all points are identified by the source to \(\Sigma_i\); at \(r=1\) they attach to \(D\) by forgetting \(i\). These formulas and their inverse are compatible with every face, so define continuous inverse maps of the quotient spaces. The induction hypothesis identifies \(|D|,|M|\), and their map with \(X',H\), and its boundary inclusion. The last display becomes the given conical pushout for \(X\). This proves the assertion. Fibrancy is not needed for this topological reconstruction; it is needed in the separate exit-category comparison.

The notation \(p=\operatorname{long}(d')\) in /043 conflicts with the preceding \(2^n\) count and its own marginal description \(d=\{\max d_1,\ldots,\max d_p\}\). The original scan has now been visually checked: the inconsistent notation is genuinely written. The construction above explicitly corrects it to \(p=|d|\). An internal audit also caught that the third transition family is essential: without it, for the constant one-point diagram on \(0<1\), the standalone \(\Sigma_0\) remains an extra isolated point. With it the colimit is the required interval. This correction is part of the construction, not a claim that those transitions were supplied in the manuscript. The attached exact cell-quotient program checks this counterexample: omitting prefix maps gives three vertices and one edge in two components, while all three transition families give two vertices and one edge. For constant diagrams on chains of zero through six vertices it checks the complete simplex face counts. The recorded run took 0.024328412 seconds.

% Independent C.1 bridge supplied by agent /root/homotopy_levels.
% AI-produced by GPT-6 Astra (OpenAI, ChatGPT Pro mode), 2026-11-19.

\paragraph{C1.B1 (Proved: the directed collage at a minimal stratum).}
Let $P$ be a finite poset and let $\mathcal C_\bullet$ be a complete
Segal object in spaces, conservative over $N(P)$. Let $i\in P$ be
minimal, let $A=\mathcal C_i$ (an $\infty$-groupoid), and let
$\mathcal D=\mathcal C|_{P\setminus\{i\}}$. Let $\mathcal M$ be the
category of arrows of $\mathcal C$ whose source is in $A$ and whose
target is in $\mathcal D$. Morphisms in $\mathcal M$ are commutative
squares, with the source-side arrow necessarily an equivalence.
Write $q:\mathcal M\to A$ and $p:\mathcal M\to\mathcal D$ for the
source and target functors. Then there is a canonical pushout formula
in $\operatorname{Cat}_\infty$
\[
 \mathcal C\simeq
 A\mathbin{\coprod}_{\mathcal M\times\{0\}}
 (\mathcal M\times[1])
 \mathbin{\coprod}_{\mathcal M\times\{1\}}\mathcal D.
 \tag{C1.B1}
\]
The maps from the two endpoints are $q$ and $p$, respectively.
The comparison is induced by the tautological natural transformation
from the source to the target of an arrow.

\begin{proof}
We verify the universal property on mapping spaces into an arbitrary
$\infty$-category $\mathcal K$. A functor from the right side to
$\mathcal K$ is a pair $F_A:A\to\mathcal K$,
$F_D:\mathcal D\to\mathcal K$, together with a natural transformation
\[
 \eta:F_Aq\Longrightarrow F_Dp
       \quad\hbox{in }\operatorname{Fun}(\mathcal M,\mathcal K).
 \tag{C1.B2}
\]
Endpoint identifications in this description are equivalences; hence
the description is a homotopy pullback of spaces of functors and
arrows, not a set of strict diagrams.

We show that these data are exactly a functor $\mathcal C\to\mathcal K$.
Restriction gives $F_A,F_D$, and evaluation on each crossing arrow
gives $\eta$. Naturality of $\eta$ follows by applying the functor to
the corresponding square in $\mathcal C$.
Conversely, use $F_A,F_D$ on the two full subcategories and use
$\eta_m$ on a crossing arrow $m$. There is no arrow from
$\mathcal D$ to $A$. Thus a composable string has at most one
crossing arrow. Precomposition with an arrow in $A$ and
postcomposition with an arrow in $\mathcal D$ are precisely the
source and target actions encoded by the squares in $\mathcal M$;
(C1.B2) makes their images commute with composition.

For clarity, the same argument includes all higher coherence, as can
be checked on the complete Segal nerves. A mixed simplex has a unique
last vertex in $A$. Its part before that vertex is a simplex in the
groupoid $A$, so it transports the crossing arrow by an equivalence.
The remaining string is the corresponding simplex of $\mathcal M$
followed by its target simplex in $\mathcal D$. Specifying its image
is therefore specifying the cone whose edges are the components of
$\eta$ on that simplex. The face maps are exactly restriction,
source transport, and target composition. The definition of a
natural transformation supplies these cones compatibly on every
simplex of $\mathcal M$, including their faces and degeneracies.
The Segal equivalences identify the spaces of such extensions with
these spaces of compatible cones; their inverse choices have
contractible fibres. Thus restriction and the construction just
given are inverse equivalences of mapping spaces, naturally in
$\mathcal K$. Yoneda proves (C1.B1).
\end{proof}

\paragraph{C1.B2 (Proved: identification of the incidence category).}
Suppose the deployment is a finite graded topological category with
object spaces $\Sigma_j$, arrow spaces $\Sigma_{jk}$ for $j<k$,
and identities only within each object space. Suppose its source
maps $s_{jk}:\Sigma_{jk}\to\Sigma_j$ are fibrations and the flag
identities and control equations are
\[
 \Sigma_{ijk}=s_{jk}^{-1}(\Sigma_{ij})\subset\Sigma_{jk},
 \qquad s_{ik}=s_{ij}s_{jk}\quad\hbox{on }\Sigma_{ijk}.
 \tag{C1.B3}
\]
Here $\Sigma_{ij}\hookrightarrow\Sigma_j$ is the target inclusion,
and $s_{jk}$ in the second equation is restricted to
$\Sigma_{ijk}\to\Sigma_{ij}$. Define a topological category $M_i$
with objects $\coprod_{j>i}\Sigma_{ij}$, arrows
$\coprod_{i<j<k}\Sigma_{ijk}$, and the induced source, target and
composition maps. Then, after passing to complete Segal objects,
$M_i$ is the $\mathcal M$ of C1.B1. Its target functor
$p:M_i\to\mathcal D$ is a left fibration. Its source functor
$q:M_i\to\Sigma_i$ is well defined by the control equation.

\begin{proof}
The target functor forgets $i$ and regards $\Sigma_{ij}$ as a
subspace of $\Sigma_j$. The flag identities identify the lifting
comparison in each degree with
\[
 (M_i)_n\longrightarrow (M_i)_0
                  \mathbin{\times^h}_{\mathcal D_0}\mathcal D_n.
 \tag{C1.B4}
\]
This is an equivalence. Indeed, the strict fibre product is the
specified iterated flag intersection; it computes the homotopy
fibre product because the source projection of each composable
string is a composite of pullbacks of the given fibrations.
The comparisons in (C1.B4), together with the Segal equivalences,
are precisely the left-fibration lifting condition for simplicial
spaces. Equivalently, an outgoing arrow from $i$ has, for each
further arrow in $\mathcal D$ from its endpoint, a contractible
space of lifts obtained by composition. The source in $\Sigma_i$
stays fixed because of the second identity in (C1.B3).

The same fibre-product calculation identifies the simplices with
those of the category of crossing arrows: in a square between two
such arrows the source-side path is an equivalence in the space
$\Sigma_i$, while the target-side arrows compose in $\mathcal D$.
Replacing strict fibres by homotopy fibres accounts for precisely
these paths. Thus the canonical comparison is an equivalence of
left fibrations over $\mathcal D$, compatible with $q$, and hence
an equivalence of the required crossing-arrow categories.
\end{proof}

In the terminology of \cite[Definition 1.1.4 and Theorem 1.1.7]{B56Haine2023},
the flag diagram is a d\'ecollage: its adjacent-pair Segal comparisons
are equivalences. That theorem identifies such diagrams with
$\infty$-categories conservative over $P$. It confirms the categorical
presentation used here; it does not replace the geometric comparison
with $\operatorname{Exit}(X)$ proved in C1.B3 below.

\paragraph{C1.B2a (Proved: coherent removal of retained outer faces).}
Suppose a finite set $R$ of already removed labels is a lower ideal
in the original poset. Give every retained outer face a product
collar, with all collar coordinates and collar transports commuting
at intersections and with the source maps of the remaining deployment
(whose labels lie outside $R$).
For a deployed space $Z$ write $Z^\circ_R$ for the space with these
outer faces deleted, keeping the remaining stratification.
There is a stratum-preserving homotopy $h^Z_t:Z\to Z$, with
$h^Z_0=\operatorname{id}$ and $h^Z_1(Z)\subset Z^\circ_R$, natural
for the source maps and incoming face inclusions among those remaining
labels. It restricts to a
homotopy of $Z^\circ_R$. Therefore
$Z^\circ_R\hookrightarrow Z$ induces an equivalence on exit categories
when the retained outer faces are ignored as separate strata.

\begin{proof}
Normalize each collar coordinate to $u\in[0,1)$. Choose a continuous
function $f:[0,1)\to(0,1)$ with $f(u)=u$ for $u\geq2/3$ and
$f(0)>0$, and with $f(u)\geq u$. Replace that coordinate at time $t$
by $(1-t)u+tf(u)$ and leave all other coordinates fixed. Extend by
the identity outside the collar. The maps agree at the collar edge.
The finite family of these operations commutes by the product-collar
hypothesis, so their composite is a globally defined homotopy.
It preserves every remaining stratum, since only the ignored outer
coordinates change. At positive time all affected zero coordinates
become positive; at every time originally positive coordinates stay
positive. Compatibility of coordinates and transports makes the
homotopies natural for the stated diagram. The time-one map and
inclusion are thus stratum-preserving homotopy inverses. Applying
exit paths sends such a homotopy to a natural equivalence: each
point's trajectory stays in one stratum, hence is an equivalence,
and the product homotopy supplies its coherence.
\end{proof}

The inclusion of a deleted outer face into its parent is not asserted
to commute with an inward push: its image moves off that face.
Such an incidence is represented by the positive-radius collar map
and its collar homotopy, as in the next paragraph.

For the next minimal label $i$ in $P\setminus R$, let
$q:H_i\to A_i=\Sigma_i$ be its boundary-link bundle. Its old outer
faces are exactly the inverse images of those of $A_i$. Indeed an
$r$-face with $r<i$ pulls back by (C1.B3); an $r$ incomparable with
$i$ has disjoint incoming face; and $i<r$ is impossible because
$R$ is a lower ideal. Hence
\[
 (H_i)^\circ_R=q^{-1}((A_i)^\circ_R).
 \tag{C1.B4a}
\]
In particular a proper bundle $q$ remains proper on these interiors.
The complementary cut space $W$ has the additional outer face $i$.
Its comparison from $(H_i)^\circ_R$ to $W^\circ_{R\cup\{i\}}$
is the positive-radius $i$-collar section, followed by the old
collar pushes. All old pushes commute with this collar, so this
map agrees with the pushed incidence map and the comparison
homotopy is the specified $i$-collar motion. Thus deleting outer
faces preserves the entire gluing diagram coherently, while
preserving properness of the link bundle. We never declare an
ignored boundary to be a separate exit stratum.

\paragraph{C1.B3 (Proved under recursive compatible collar hypotheses).}
Assume, in addition to C1.B2, that the deployment is realized by
proper conically smooth, stratified locally trivial source bundles
and mutually compatible conical collars. Require explicitly that
cutting a minimal stratum produces a collared space $W$
with retained outer boundary $H$, that $H\to\Sigma_i$ is such a
proper bundle in collars, and that both $W$ and $H$ inherit deployments of the
same class with fewer stratum labels. Require that their induced
collars commute with all source maps. This is the recursive
equisingularity condition needed in the induction; mere fibrancy
in C1.B2 is insufficient for this geometric assertion.
Then the complete Segal object of the deployment presents
$\operatorname{Exit}(X)$, naturally with respect to the boundary
inclusions and source maps used in the deployment.

\begin{proof}
Induct on the number of nonempty labels. With one label, the
deployment is its stratum space, presenting its fundamental
$\infty$-groupoid, which is its exit category. Remove a minimal
label $i$. By the boundary-compatible induction hypothesis,
$\mathcal D\simeq\operatorname{Exit}(W)$ and
$M_i\simeq\operatorname{Exit}(H)$, compatibly with
$H\hookrightarrow W$ and $H\to\Sigma_i$.
The punctured cone collar is $H\times(0,\epsilon)$, whose exit
category is $\operatorname{Exit}(H)$ because the interval carries
the trivial stratification and is contractible. Use C1.B2a to remove all previously retained outer faces;
this identifies exit categories and restricts the link bundle by
base change, hence preserves properness. The same collar homotopy
identifies $\operatorname{Exit}(X)$ with the exit category of its
old outer interior. The fibrewise
cone comparison on these interiors gives
\[
 \operatorname{Exit}(C(H\to\Sigma_i))
 \simeq \Sigma_i\mathbin{\coprod}_{M_i\times\{0\}}
                       (M_i\times[1]).
 \tag{C1.B5}
\]
This is exactly \cite[Lemma 2.13]{C_Volpe}, together with
$\operatorname{Exit}(H\times[0,\infty))\simeq
M_i\times[1]$; the boundary point is the smaller stratum.
The geometric cone-collar gluing is an open-cover gluing after
slightly enlarging the collar and shrinking the complementary
cut piece. We first apply C1.B2a: delete the old outer faces
in $H$ and $\Sigma_i$, and all outer faces in $W$. The source
bundle remains proper by (C1.B4a), so Volpe is applied only to
these ordinary conically smooth interiors. The map to the
complementary open piece is the positive-radius collar section
specified in C1.B2a. Its coherent collar homotopies identify this
diagram with the displayed diagram on $W,H,\Sigma_i$. Exit excision,
equivalently \cite[Corollary 2.17]{C_Volpe} or
\cite[Theorem A.7.1]{LurieHA}, therefore gives
\[
 \operatorname{Exit}(X)\simeq
 \operatorname{Exit}(W)
       \mathbin{\coprod}_{\operatorname{Exit}(H)}
 \operatorname{Exit}(C(H\to\Sigma_i)).
 \tag{C1.B6}
\]
Substitute the two induction comparisons and (C1.B5), and compare
with the identical universal pushout (C1.B1). This yields the
equivalence of full exit categories. The construction uses the
source and inclusion maps and their tautological transformations,
so the comparisons commute with the retained boundaries at the
next induction step. The naturality is part of the pushout
universal properties, not an additional choice of independent
equivalences of individual spaces. This completes the induction.
\end{proof}


\subsection{C1.04 --- Proved under the corrected hypotheses: incidence groupoids}\label{c1.04-proved-under-the-corrected-hypotheses-incidence-groupoids}

Here are the construction and the category in which the moduli version holds. For finite-depth smooth Thom--Mather data, perform the successive cuts in a linear order in which every lower stratum is cut before the strata above it, retaining every new boundary. At the stage of label \(i\), the cut-open closure of its stratum is \(\Sigma_i\). The face of \(\Sigma_j\) over an incidence \(i<j\) is \(\Sigma_{ij}\), and its source is the retained link-bundle projection to \(\Sigma_i\). A longer flag \(d=(i_0<\cdots<i_n)\) defines \[
 \Sigma_d=\Sigma_{i_0i_n}\cap\cdots\cap\Sigma_{i_{n-1}i_n}.
\] Removing a nonfinal entry gives the inclusion of intersections. Removing a final entry gives the appropriate source projection restricted to that intersection. Projection control identifies all composites with the same source and target, so these maps define the contravariant flag functor. The inverse-image identities make its Segal squares cartesian. Proper stratified link bundles and their base changes prove the initial-segment fibrancy; the incoming target faces have their specified commuting collars. The empty-index case gives the empty space and the empty category. A one-stratum deployment is its fundamental infinity-groupoid.

The hypotheses needed in C1.B3 hold for this construction: a cut or an intersection of faces retains the same iterated boundary fibrations, their proper source bundles, and the restrictions of the commuting collars. C2.05 derives this closure property from the iterated-fibration resolution theorem. In the conically smooth formulation it is supplied by AFT, \emph{Local structures on stratified spaces}, Theorem 7.3.8 (full unzip), Proposition 7.3.10 (the blow-up square), and Proposition 8.2.5 (the tubular presentation) \cite{C_AFT}. Thus C1.B3 identifies the resulting complete Segal object with the full exit category in that class. C1.01 explains why the weaker literal mapping-cone data in \S{}8 do not imply this result without equisingularity.

For \(\overline{\mathcal M}_{g,n}\), use the following finite incidence category, retaining automorphisms. Its objects are stable graphs \(\Gamma\) of type \((g,n)\), with labelled legs, vertex genera, and stability at every vertex. An arrow is a specified edge contraction \(c:\Gamma\to\Gamma'\), together with its identification of the contracted graph with \(\Gamma'\); graph isomorphisms are arrows as well. Composition contracts the union of the specified edge sets. Contracting a loop raises the genus of its vertex; contracting an edge joining two vertices merges them and adds their genera. The open stratum is the orbifold \[
 \mathcal S_\Gamma=
 \left[\prod_{v\in V(\Gamma)}\mathcal M_{g(v),n(v)}
          \big/\operatorname{Aut}(\Gamma)\right].
\] Equivalently one may keep half-edge labels throughout and then form the indicated finite groupoids. Fixed contractions preserve which node branches are involved; an orbit-poset does not.

Use the cutting construction in smooth universal deformation charts. Looijenga's Theorem 9.13 \cite{C_Looijenga} supplies the chart and its finite automorphism action; its boundary parameters \(q_e\) are indexed by the actual nodes. The ordinary real oriented blow-up writes \(q_e=r_e\exp(i\theta_e)\). To obtain the full deployment with a boundary hypersurface for every strict stratum incidence, further blow up its corner faces successively in increasing dimension. Locally this replaces the face lattice by its flag lattice. In a normal coordinate space with \(r\) nodes, the outgoing link before resolution is the coordinate-stratified sphere \(S^{2r-1}\); the source fibre of the full deployment is its recursively resolved link. For \(r=1\) this is a circle. For \(r>1\), radial-ratio coordinates and their faces are also retained, so identifying this full link with only \((S^1)^r\) would discard data. The ordinary oriented blow-up and the full flag deployment are related by these explicit further corner blow-ups.

For a composable chain of specified contractions, take the corresponding flag-face space in these charts and its groupoid of chart identifications and graph automorphisms. Denote the resulting orbifold by \(\boldsymbol\Sigma_d\). Face and source maps are the preceding geometric maps with their equivariant identifications. In simplicial degree \(k\), take the disjoint union over contraction chains of length \(k\), modulo their isomorphism groupoids. Thus fibre products are homotopy fibre products of groupoids/spaces, not orbit-set fibre products. Local transverse product collars and proper link bundles show the Segal comparisons are equivalences. The maps glue because changes of deformation chart preserve the labelled normal-crossings incidences and their resolved boundary fibrations. The source and target maps therefore define a Segal object in orbifold homotopy types.

There are two equivalent descriptions of its completion. In the classifying-space construction take the classifying space of the entire groupoid of chains in every degree, with isomorphisms at every vertex and commuting identifications on each arrow. Replacing only the object set by classifying spaces, while leaving the raw arrow sets unchanged, would not suffice. For example, for one object with automorphism group \(G\), the groupoid of length-\(k\) strings is \[
 G^k\mathbin{/\mkern-6mu/}G^{k+1}.
\] The action changes the trivializations at each vertex. It is transitive with stabilizer a diagonal copy of \(G\), so its classifying space is \(BG\) in every degree. The degree-zero degeneracy to the equivalence subspace in degree one is therefore an equivalence. This is the Rezk classifying-diagram construction. Applied to the entire incidence diagram, with homotopy fibre products for compositions, it gives a complete-Segal model. Equivalently take the Rezk completion of the raw incidence/action-category nerve. A raw group nerve with one discrete object is not already complete when its automorphism group is nontrivial, and discarding that group is not completion.

Choose the compatible boundary fibrations and collar data once on the orbifold resolution, and pull them back to each atlas chart. These data can be constructed equivariantly: normal-crossings blow-ups preserve incidence labels, invariant metrics and connections exist by finite-group averaging and orbifold partitions of unity, and the iterated collar construction proceeds on the preserved boundary fibrations. A change of chart consequently preserves the pulled-back structures. Independent arbitrary collar choices are not being asserted to agree. The chartwise source maps are representable proper locally trivial bundles, so their pullbacks compute the needed homotopy fibre products after passage to atlas chain groupoids.

Define the orbifold exit category by descent from an orbifold atlas: take the colimit in infinity-categories of the exit categories of its simplicial atlas and its chart identifications. On a finite quotient chart \([U/G]\), this is \(\operatorname{Exit}(U)\mathbin{/\mkern-6mu/}G\). The atlas definition is independent of the atlas. For a surjective stratum-preserving local diffeomorphism \(V\to U\), cover \(U\) by opens admitting stratum-preserving local sections. A section supplies an extra degeneracy in the augmented \v{C}ech diagram, making its image under the exit functor a split augmented simplicial diagram. Its colimit is the exit category of that open. Exit open-cover excision \cite[Theorem A.7.1]{LurieHA} and interchange of colimits globalize the equivalence to \(U\). Two atlases have a common refinement given by their fibre product; applying the preceding argument in both directions of its bisimplicial atlas identifies the two colimits. For \([U/G]\), the atlas nerve is \(U\times G^\bullet\), giving precisely the homotopy-orbit category just displayed, including ineffective isotropy.

The comparison of C1.B3 is natural for the face maps, source maps, and isomorphisms. It is consequently an equivariant equivalence on every deformation chart and on every atlas intersection. Taking the displayed colimit gives \[
 \widehat{N\boldsymbol\Sigma}\simeq
 \operatorname{Exit}_{\mathrm{orb}}(\overline{\mathcal M}_{g,n}),
\] where the hat denotes complete-Segal localization. The proof is categorical descent of chartwise equivalences; equality of classifying-space realizations alone is not used. Realization gives the homotopy type of the orbifold, with its finite isotropy. On a level cover on which all stabilizers act freely, the same construction gives the ordinary space and its ordinary exit category.

For coarse moduli spaces, finite orbit quotients commute with the explicit topological colimits of C1.03, so their cone realization still recovers the coarse space. However quotienting can destroy source fibrancy and can make the stated strata nonmanifolds; C3.03 gives a concrete interior obstruction in genus four. Thus the orbifold exit comparison cannot simply be transferred to the boundary-only coarse stratification. One may instead refine its strata by the isotropy/quotient singularity type and apply the ordinary conical theorem to that refined structure. Fixed contractions alone do not perform that refinement.

For the basic case \(\overline{\mathcal M}_{0,4}\cong\mathbb P^1\), the ordinary space has three point strata. Its open object is a sphere with three open discs removed, its point objects are those three points, and its three arrow spaces are the three boundary circles. The source on each circle maps to its point object; the target is the corresponding boundary embedding. Gluing the three mapping cylinders fills the three holes by discs, giving \(S^2\). Thus the cut-open convention does satisfy the presentation in this case; taking the open stratum or a boundary circle as its point object would be a different datum. This verifies the specific inherited \((0,4)\) comparison independently.

Empty labels may be deleted, and disconnected strata may be kept as disjoint unions. A grading need not be injective on components: a strictly increasing grading can merge incomparable components, but their actual incidences and disjoint incoming branches must remain in the diagram. The orbifold formulation allows ineffective actions as well. A stabilizer that acts trivially on a deformation chart still contributes its classifying space; quotienting by its kernel would change the stack and its exit category. In particular the vertex moduli factors retain their own automorphisms in addition to graph automorphisms. No divisibility or dense-order assumption on the indexing order is made or needed.

Finally this is a homotopy and exit-category presentation, not the missing topos gluing theorem of A.1. The successive cylinders retain the entire face-incidence diagram and its coherence. The manuscript's pair/triple stones omit actual iterated neighbourhood intersections, as the counterexample of A.1 shows. Taking all complete flag-intersection pro-topoi and their \v{C}ech descent maps gives the topos reconstruction there; taking only the homotopy types or only the global stratum-poset forgets information needed for arbitrary sheaves. These are different functors applied to the geometric descent diagram, and their colimit claims have been proved in their respective categories.

\section{C.2. Projection control and Thom--Mather data}\label{c.2.-projection-control-and-thommather-data}

\subsection{C2.01 --- Refuted: the bare intersection axioms imply transitivity}\label{c2.01-refuted-the-bare-intersection-axioms-imply-transitivity}

The ill-typed b\ensuremath{{}_{ij}}\ensuremath{{}^{-1}}(\ensuremath{\Sigma}\ensuremath{{}_{ij}}) is interpreted as s\ensuremath{{}_{jk}}\ensuremath{{}^{-1}}(\ensuremath{\Sigma}\ensuremath{{}_{ij}}) inside \ensuremath{\Sigma}\ensuremath{{}_{jk}}. The bare conditions say:

\begin{enumerate}
\def\labelenumi{\arabic{enumi}.}
\tightlist
\item
  b\ensuremath{{}_{ij}}:\ensuremath{\Sigma}\ensuremath{{}_{ij}}\ensuremath{\to}\ensuremath{\Sigma}\ensuremath{{}_{j}} is a closed embedding.
\item
  \ensuremath{\Sigma}\ensuremath{{}_{ik}}\ensuremath{\cap}\ensuremath{\Sigma}\ensuremath{{}_{jk}} is empty for incomparable i,j.
\item
  \ensuremath{\Sigma}\ensuremath{{}_{ik}}\ensuremath{\cap}\ensuremath{\Sigma}\ensuremath{{}_{jk}}=s\ensuremath{{}_{jk}}\ensuremath{{}^{-1}}(\ensuremath{\Sigma}\ensuremath{{}_{ij}}), for i\textless j\textless k.
\end{enumerate}

They do not imply s\ensuremath{{}_{ik}}=s\ensuremath{{}_{ij}}s\ensuremath{{}_{jk}} on the common subset.

\textbf{Finite example.} Let I=\{0\textless1\textless2\}, every space be \{0,1\}, every b the identity, s\ensuremath{{}_{01}}=s\ensuremath{{}_{12}}=id, and s\ensuremath{{}_{02}}(z)=1\ensuremath{-}z. Conditions 1--3 hold; all subsets in condition 3 are the whole set. But (0,0)\textless(1,0)\textless(2,0), while no arrow (0,0)\ensuremath{\to}(2,0) exists. The analogous failure occurs at 1. The exact program code\_C/check\_cote150.py verifies all six arrows and both missing composites.

\textbf{Smooth compact proper enhancement.} Let A=S\ensuremath{{}^{1}} and g:A\ensuremath{\to}A be rotation by \ensuremath{\pi}. Take \ensuremath{\Sigma}\ensuremath{{}_{0}}=A, \ensuremath{\Sigma}\ensuremath{{}_{1}}=A\ensuremath{\times}{[}0,1{]}, \ensuremath{\Sigma}\ensuremath{{}_{2}}=A\ensuremath{\times}{[}0,1{]}\ensuremath{{}^{2}}. Let \ensuremath{\Sigma}\ensuremath{{}_{01}}=A\ensuremath{\times}\{0,1\}\ensuremath{\subset}\ensuremath{\Sigma}\ensuremath{{}_{1}}. In coordinates (z,u,v) on \ensuremath{\Sigma}\ensuremath{{}_{2}} put

\ensuremath{\Sigma}\ensuremath{{}_{12}}=A\ensuremath{\times}{[}0,1{]}\ensuremath{\times}\{0,1\}, \ensuremath{\Sigma}\ensuremath{{}_{02}}=A\ensuremath{\times}\{0,1\}\ensuremath{\times}{[}0,1{]}.

All targets are inclusions. Set s\ensuremath{{}_{01}}(z,u)=z, s\ensuremath{{}_{12}}(z,u,v)=(z,u), and s\ensuremath{{}_{02}}(z,u,v)=g(z). They are proper smooth bundles with fibres two points, two points, and two intervals, respectively. Opposite faces are transverse; their intersection A\ensuremath{\times}\{0,1\}\ensuremath{{}^{2}} is exactly s\ensuremath{{}_{12}}\ensuremath{{}^{-1}}(\ensuremath{\Sigma}\ensuremath{{}_{01}}). Every boundary component is accounted for. The maps respect base collars: s\ensuremath{{}_{12}} is a product projection and \ensuremath{\Sigma}\ensuremath{{}_{0}} has no boundary. Thus d)-e), properness, equisingularity, and the nonsingular condition do not repair the omission. At every corner s\ensuremath{{}_{01}}s\ensuremath{{}_{12}}(z,u,v)=z but s\ensuremath{{}_{02}}(z,u,v)=g(z).

\textbf{Precision about the manuscript.} It does explicitly call (41) a ``diagramme commutatif hexagonal.'' Its left commutativity is the needed equation. The false inference is ``donc'' from the bare conditions; if commutativity of (41) is included as an axiom it is already written on the page. The inherited SA10-K150e-02 should be qualified accordingly.

\subsection{C2.02 --- Proved: the projection-control repair}\label{c2.02-proved-the-projection-control-repair}

Add to 1--3 the identity

s\ensuremath{{}_{ik}}\textbar \ensuremath{\Sigma}\ensuremath{{}_{ijk}}=s\ensuremath{{}_{ij}}\ensuremath{\circ}(s\ensuremath{{}_{jk}}\textbar \ensuremath{\Sigma}\ensuremath{{}_{ijk}}).

The restriction on the right has target \ensuremath{\Sigma}\ensuremath{{}_{ij}} by condition 3. For x\ensuremath{\in}\ensuremath{\Sigma}\ensuremath{{}_{i}},y\ensuremath{\in}\ensuremath{\Sigma}\ensuremath{{}_{j}} define x\textless y iff i\textless j, y\ensuremath{\in}\ensuremath{\Sigma}\ensuremath{{}_{ij}}, and x=s\ensuremath{{}_{ij}}(y). Irreflexivity follows from strict labels. If x\textless y\textless z, with i\textless j\textless k, then z\ensuremath{\in}\ensuremath{\Sigma}\ensuremath{{}_{jk}} and s\ensuremath{{}_{jk}}(z)=y\ensuremath{\in}\ensuremath{\Sigma}\ensuremath{{}_{ij}}. Condition 3 puts z in \ensuremath{\Sigma}\ensuremath{{}_{ijk}}; the new identity gives s\ensuremath{{}_{ik}}(z)=x. Hence the relation is transitive.

At most one predecessor of a point has a given label, since b\ensuremath{{}_{ik}} is an embedding. Two predecessor labels must be comparable by condition 2. If i\textless j, condition 3 and projection control order the corresponding predecessors. Thus every lower interval is totally ordered and its order is induced by I.

Conversely, suppose this relation is transitive and each lower interval has the order induced by I. For z\ensuremath{\in}\ensuremath{\Sigma}\ensuremath{{}_{ik}}\ensuremath{\cap}\ensuremath{\Sigma}\ensuremath{{}_{jk}} with i\textless j, its predecessor of label i must precede its predecessor of label j, giving the new identity. Transitivity gives the reverse inclusion in condition 3. Thus the subset identity and the projection identity are exactly the needed algebraic facts. Longer compositions are associative because an ordered category has at most one arrow between two fixed objects.

\subsection{C2.03 --- Proved: one radius level does not determine control}\label{c2.03-proved-one-radius-level-does-not-determine-control}

In a positive-radius tubular chart write coordinates (r,y), where r\textgreater0 is a lower radial function and y\ensuremath{\ge}0 is normal to an intermediate stratum. The controlled retraction is \ensuremath{\pi}(r,y)=(r,0). Fix \(\epsilon>0\) and choose explicitly \(h(y)=-y/(4(1+y))\), so \(-1/4<h(y)\le0\), with equality only at \(y=0\). Set

\ensuremath{\pi}\ensuremath{\prime}(r,y)=(r+(r\ensuremath{-}\ensuremath{\epsilon})h(y),0).

It is a retraction. The new radial coordinate is positive because \(r(1+h(y))-\epsilon h(y)>0\) for \(r>0\). Its derivative in \(r\) is \(1+h(y)>3/4\), so \((\pi\prime,y)\) is a diffeomorphism onto a variable-width neighbourhood of the base and is still a tubular coordinate map after shrinking. Both maps agree on r=\ensuremath{\epsilon} and have the same inverse image of that level, because

r+(r\ensuremath{-}\ensuremath{\epsilon})h(y)\ensuremath{-}\ensuremath{\epsilon}=(r\ensuremath{-}\ensuremath{\epsilon})(1+h(y)).

But r\ensuremath{\pi}\ensuremath{\prime} is not r off that level. Hence a level-set incidence condition is strictly weaker than rho-control. If one fixed map preserves every small level, the full equation follows by choosing \ensuremath{\epsilon} equal to the radial value of each point. Coherent radial collars or all levels are required to recover the functional equation from deployment data.

\subsection{C2.04 --- Refuted: the bare recapitulation is Thom--Mather data}\label{c2.04-refuted-the-bare-recapitulation-is-thommather-data}

Take one label and \ensuremath{\Sigma}\ensuremath{{}_{0}} a Cantor set. Every a)-e) and projection-control condition is vacuous. A Thom--Mather stratum must be a manifold; a Cantor set is not. Thus the bare completed recapitulation is not such an equivalence. The established comparison concerns suitably smooth manifolds with corners, compatible iterated boundary fibrations, and radial collar information, not arbitrary spaces and one isolated level. Even adding compact manifold object and arrow spaces and proper Hurewicz source fibrations is insufficient: the two-stratum examples C5.06 satisfy the triple and projection conditions vacuously and have nonmanifold fibres. Thus the regular source-bundle hypotheses in C2.05 have mathematical content.

\subsection{C2.05 --- Known, with derivation: the smooth resolution theorem}\label{c2.05-known-with-derivation-the-smooth-resolution-theorem}

The exact smooth replacement uses the whole boundary-fibration structure, not just one incidence level. Take a finite deployment whose nonempty object spaces are smooth manifolds with corners. Require every incoming \(\Sigma_{ij}\subset\Sigma_j\) to be a collective embedded boundary hypersurface, with their union the whole boundary, with incomparable incoming faces disjoint, and with the specified transverse product collars. Require the sources to be proper smooth locally trivial bundles with compact manifold-with-corners fibres, compatible with all base and fibre faces. On a triple require both the subset identity and projection control of C2.02. Require that the boundary of a source fibre consists exactly of the intermediate-label faces, and choose compatible radial collar germs.

In these hypotheses the data are precisely a system of manifolds with iterated boundary fibrations. To see the correspondence without suppressing labels, on a maximal piece \(\Sigma_m\) the hypersurface labelled \(i\) is \(\Sigma_{im}\), its base is \(\Sigma_i\), and its fibration is \(s_{im}\). At an intersection \(i<j<m\), the diagram is \[
 \Sigma_{ijm}\xrightarrow{s_{jm}}\Sigma_{ij}\subset\Sigma_j,
 \qquad
 s_{im}=s_{ij}s_{jm}.
\] This is exactly the commuting boundary-fibration diagram. Its inverse-image identity says that the horizontal boundary is the pullback of the base boundary. All other boundaries are the stated vertical fibre boundaries. Higher intersections are recovered by iterating these identities, so no further arbitrary associativity datum remains. Conversely such an iterated-fibration system supplies all the \(\Sigma_i,\Sigma_{ij}\), their intersections, sources, and collars. If two top pieces have a common lower stratum, retain the common base and use the disjoint union of their faces as one collective hypersurface. Separating that collective hypersurface would instead normalize a disconnected-link singularity.

Albin, \cite[Theorem 6.3(a)--(c), p.48]{C_Albin2016}, gives resolution and blowdown as inverse constructions for compact Thom--Mather spaces and manifolds with these iterated fibrations, after choosing compatible collars. Its paragraph following the theorem states explicitly that without fixed collars the blowdown is defined only up to isomorphism. The collective-hypersurface qualification and normalization warning are on p.49. The theorem's proof reduces to ALMP, \cite[Propositions 2.3 and 2.5, pp.253--257]{C_ALMP}; Propositions 2.6--2.7, p.258, give the isomorphism comparison.

For finite-depth, second-countable, locally compact spaces without compactness, the precise modern formulation is Albin--Banagl--Piazza, \emph{Smooth atlas stratified spaces}, \cite[Theorems 2.10 and 2.13, pp.24--28]{C_ABP}. The former proves the resolution/blowdown equivalence for its smooth atlas category, with maps preserving the regular part and the interior respectively; the latter identifies the geometric objects with Thom--Mather spaces. We use its object and isomorphism statement. We do not identify arbitrary continuous stratified maps, arbitrary conically smooth maps, and its liftable smooth maps.

Here is why the cited theorem gives precisely the needed control identities. In the collar of a boundary fibration \(H_i\to\Sigma_i\), write a point as \((h,r_i)\) and collapse \(r_i=0\) along that fibration. Define \(\pi_i(h,r_i)=s_i(h)\) and \(\rho_i(h,r_i)=r_i\). On an overlapping lower collar \(k<i\), source compatibility makes \(\pi_i\) preserve the \(r_k\) coordinate, so \(\rho_k\pi_i=\rho_k\); triple commutativity gives \(\pi_k\pi_i=\pi_k\). Away from the zero radius the product coordinates and smooth bundle maps make \((\pi_i,\rho_i)\) a stratumwise submersion. These are Mather's two control identities and submersion requirement \cite[section 7, p.15]{C_Mather}. Conversely cutting a controlled tube at \(r_i=\epsilon_i\) gives the incoming face and its source. The two control equations restrict respectively to the subset identity and projection identity, and the controlled collar flow supplies the omitted radial germs. Iterating gives the cited resolution.

A single chosen \(\epsilon\) records only the corresponding level diagram. To claim equivalence of \emph{control data} one must retain the radial collars, or equivalently all sufficiently small levels with their compatible identifications. C2.03 proves why this extra datum cannot be discarded. Different collar choices give isomorphic resolved geometric objects; they need not give equal numerical radius functions. Empty strata and empty hypersurfaces may be deleted. Disconnected strata are handled componentwise or with collective faces; one must not silently replace a disconnected fibre by its components.

\section{C.3. Corners, flags, and moduli incidences}\label{c.3.-corners-flags-and-moduli-incidences}

\subsection{C3.01 --- Refuted: precylindrical alone implies a manifold}\label{c3.01-refuted-precylindrical-alone-implies-a-manifold}

The definition 150/028 asks for collared predecessor embeddings, transverse collars, and Boolean intervals. It does not require strata to be manifolds. Let K be the middle-thirds Cantor set, M=K\ensuremath{\times}{[}0,\ensuremath{\infty}), X\ensuremath{{}_{0}}=K\ensuremath{\times}\{0\}, and X\ensuremath{{}_{1}}=M. The interval \{0\textless1\} is Boolean. The required collar is the identity K\ensuremath{\times}{[}0,\ensuremath{\delta})\ensuremath{\to}M. The family is finite and satisfies (14). Thus it is precylindrical.

No point has a locally connected neighborhood, because arbitrarily small clopen separations in K persist in the product. Every Euclidean orthant is locally connected. Thus M is not a manifold with corners.

\textbf{Exact local correction.} A finite precylindrical structure is a manifold-with-corners structure precisely when the transverse bases of its local collar products are manifolds of the required dimensions. A local product \(B\times[0,\delta)^k\) becomes an orthant chart if B\ensuremath{\cong}\(\mathbb R^d\) locally. Conversely an orthant chart preserving the specified faces restricts to a manifold chart on their full intersection. The requirement to preserve the specified faces matters: TOP has exotic collared factors of ambient manifolds.

For conically smooth finite-depth spaces with manifold strata, AFT \cite[Theorem 7.3.8]{C_AFT} gives the full unzip as a smooth manifold with corners. The local step replaces \(\mathbb R^d\)\ensuremath{\times}C(L) by \(\mathbb R^d\)\ensuremath{\times}{[}0,\ensuremath{\infty})\ensuremath{\times}L; induction resolves L. This supplies the corrected theorem in that exact category, not for arbitrary \S{}8 spaces.

\subsection{C3.02 --- Proved: Boolean intervals of flags}\label{c3.02-proved-boolean-intervals-of-flags}

For a maximal label i, deployed faces are \(\Sigma_d\) for flags d ending at i, with the reverse-inclusion order. For e\ensuremath{\subset}d with the same last entry, an intermediate flag is exactly e\ensuremath{\cup}A for A\ensuremath{\subset}\(d\setminus e\). Every such subset is a flag. Hence the interval {[}d,e{]} is the Boolean lattice of subsets of \(d\setminus e\) (with orientation adjusted to the chosen ordering convention). Transverse collars supply one independent normal coordinate for each added flag entry.

The finite program checks every naturally labelled poset of size at most five: counts 1,1,2,7,40,357 for sizes 0,\ldots,5, comprising 4,135 flags and 10,603 intervals. The argument above proves the assertion in all finite sizes; these counts only check the implementation of the combinatorial description.

For a stratum with several local branches, a global label does not identify a branch. The face index must retain the local incidence or boundary germ. For a normal crossings divisor with k local branches, faces are locally indexed by subsets of those branches. Monodromy and automorphisms may permute them, so the global description is an incidence category or an orbifold groupoid. An orbit-poset loses that action.

\subsection{C3.03 --- Refuted on coarse moduli: contractions remove isotropy}\label{c3.03-refuted-on-coarse-moduli-contractions-remove-isotropy}

This is a counterexample on the coarse space as well as an orbifold qualification. Let C be the smooth projective completion of

y\ensuremath{{}^{2}}=x(x\ensuremath{-}1)(x\ensuremath{-}3)(x\ensuremath{-}7)(x\ensuremath{-}12)(x\ensuremath{-}20)(x\ensuremath{-}31)(x\ensuremath{-}45)(x\ensuremath{-}64).

There are nine distinct finite branch points and infinity, so Riemann--Hurwitz gives genus 4. The exact program checks all 10\ensuremath{\cdot}9\ensuremath{\cdot}8=720 Möbius maps specified by the images of (\ensuremath{\infty},0,1). Only the identity preserves the whole branch set. This is exhaustive even over C: the unique map carrying three rational projective points to three rational projective points is defined over Q. The hyperelliptic quotient is unique here: two distinct degree-two maps to \(\mathbb P^1\) would give a birational image of bidegree \((2,2)\) in \(\mathbb P^1\times\mathbb P^1\), whose arithmetic genus is one. If their joint map instead had degree two, the image would have bidegree \((1,1)\), making the maps equivalent. Both alternatives exclude distinct degree-two quotients for this genus-four curve. Every automorphism consequently descends to that quotient, with its two-element deck group as kernel. Thus Aut(C)=\{1,\ensuremath{\iota}\}, \ensuremath{\iota}(x,y)=(x,\ensuremath{-}y).

A neighborhood of {[}C{]} in M\ensuremath{{}_{4}} is a finite quotient of the smooth universal deformation germ, with tangent \(H^1(C,T_C)\). A finite holomorphic action is linearizable: average the initial coordinate map \(z\) as \(|G|^{-1}\sum_{g\in G}(dg_0)^{-1}z\circ g\). The average is equivariant, with identity derivative, so it is a coordinate map. Serre duality identifies the dual tangent representation with \(H^0(C,2K_C)\).

An invariant basis is \(x^k\)(dx/y)\ensuremath{{}^{2}} for 0\ensuremath{\le}k\ensuremath{\le}6, and an anti-invariant basis is \(x^k\) dx\ensuremath{{}^{2}}/y for 0\ensuremath{\le}k\ensuremath{\le}1. At a finite branch point use x\ensuremath{-}a=t\ensuremath{{}^{2}},y=t\ensuremath{\cdot}unit, so regularity is immediate. At infinity x=t\ensuremath{{}^{-2}},y=t\ensuremath{{}^{-9}}\ensuremath{\cdot}unit; the two displayed differentials have orders 12\ensuremath{-}2k and 3\ensuremath{-}2k, respectively. The bounds are exactly those making the orders nonnegative. Independence is clear in the function field, and 7+2=9=3g\ensuremath{-}3, so the forms span. The quotient germ is consequently

C\ensuremath{{}^{7}}\ensuremath{\times}(C\ensuremath{{}^{2}}/\{\ensuremath{\pm}1\}).

The second factor is the cone on RP\ensuremath{{}^{3}}. The product link is S\ensuremath{{}^{13}}*RP\ensuremath{{}^{3}}\ensuremath{\cong}\ensuremath{\Sigma}\ensuremath{{}^{14}}RP\ensuremath{{}^{3}}. The integral cellular complex of RP\ensuremath{{}^{3}} has d\ensuremath{{}_{2}}=2,d\ensuremath{{}_{1}}=d\ensuremath{{}_{3}}=0; hence H\ensuremath{{}_{1}}(RP\ensuremath{{}^{3}};Z)=Z/2. Local homology at the quotient origin therefore has Z/2 in degree 16, besides its top Z in degree 18. That is impossible for a topological manifold. This point lies in the interior, so blowing up the Deligne--Mumford boundary leaves the local obstruction untouched. Indexing contractions cannot remove it.

\textbf{Correct moduli statement.} A universal deformation chart B of a stable marked curve has normal crossing boundary, with branches indexed by the actual nodes e. Replacing q\_e by r\_e exp(i\ensuremath{\theta}\_e) gives a manifold-with-corners chart, acted on by the finite automorphism group. These charts define an orbifold with corners. On an actual torsion-free level cover they define a manifold with corners. Fixed contractions record the subsets of nodes that remain unsmoothed, not a removal of isotropy.

The precise source is Looijenga \cite[Theorems 8.11 and 9.13; Proposition 9.16]{C_Looijenga}: deformation tangent, stable charts, and orbifold real oriented blow-up, respectively.

\subsection{C3.04 --- Proved: the exact local corner criterion}\label{c3.04-proved-the-exact-local-corner-criterion}

First distinguish the literal \S{}8 input from a collection of unrelated cone charts. For the finite recursive conical data actually defined on 150/035--036, retaining the boundaries always gives a precylindrical structure in the manuscript's collar sense. This assertion does not require source fibrancy. The induction is explicit in the original /036, visually checked: after cutting the minimal strata, (22) contains the old collared boundaries together with every new link boundary, as a mutually transverse family transverse to the remaining strata; the next conical datum is required on this entire system.

Here is the inductive verification. A cut replaces each selected cone direction by its radius interval and retains its zero-radius boundary. Conical transversality to an old boundary says in coordinates that this cut is a product with the old collar coordinate; hence each old collar restricts to the cut piece and commutes with the new collar. Distinct simultaneously cut minimal strata are disjoint, so their new faces are disjoint. By the recursive requirement the next cuts are products with all the retained coordinates. This proves the transverse-collar axioms at every stage and on every removed base piece \(\Sigma_i\). Two incoming labels that are incomparable cannot meet: their original closed strata intersect only in lower strata, and, after those lower strata have been cut, a common positive-radius link point would give a point of both original strata outside that lower intersection. Thus all nonempty face intersections have flag labels. Their subface intervals are the Boolean intervals proved in C3.02. The intersection of two such closed faces is the face for the union of their labels when that union is an admissible flag, and is empty otherwise; local incidence branches are retained separately. Including the full pieces and deleting empty faces verifies the finite closed-cover and intersection axioms (14). This proves precylindricity for each cut-open piece and hence for the maximal pieces comprising the full unzip of the literal recursive \S{}8 datum. Merely choosing separate conical neighborhoods without this retained-boundary compatibility is not that datum.

The remaining question is exactly when this precylindrical output is a manifold with its specified corners. For any of TOP, PL, or smooth, fix a finite precylindrical structure with the stated compatible collars and incidence branches. At a point \(x\), let \(F_x^\circ\) be the open face obtained by intersecting precisely the boundary germs through \(x\), and deleting all smaller faces. Let \(k(x)\) be the number of those independent germs. The collar axioms provide a germ \[
 (X,x)\cong (F_x^\circ,x)\times([0,\infty)^{k(x)},0)
\] respecting each specified face. The resulting space is a manifold with these corners in the chosen category if and only if every \(F_x^\circ\) is a manifold in that category and the collar transition maps have that regularity. Its dimension is locally constant; for a single prescribed dimension \(n\), also require \(\dim F_x^\circ+k(x)=n\) for every point. Empty outputs satisfy the convention vacuously.

Indeed, a local manifold chart on \(F_x^\circ\), followed by the collar product, is exactly an orthant chart of dimension \(\dim F_x^\circ+k(x)\). Conversely a corner chart preserving all the prescribed faces restricts on their full intersection to a manifold chart on \(F_x^\circ\). Both implications are local; compatibility gives an atlas. Hausdorffness and second countability are required separately for the usual definition of manifold and are inherited in the finite resolution settings considered here.

This is an exact criterion rather than an assertion that every collared topological factor of a manifold is itself a manifold. In TOP that cancellation assertion is false by C5.03--04. Even the open cone on \(F=\Sigma H\) from C5.03 is homeomorphic to \(\mathbb R^5\), while its prescribed conical cutting produces \(F\times[0,\infty)\), which is not a manifold with boundary. Thus ordinary manifoldness of the original strata cannot alone validate an arbitrary topological choice of conical data.

For recursive finite-depth smooth or PL conical data whose resolved link pieces and base strata are manifolds and whose transitions lift compatibly to their resolutions, the criterion holds inductively: replacing a local cone \(B\times C(L)\) by \(B\times[0,\infty)\times\widetilde L\) gives its required orthant charts. The same is a sufficient criterion in TOP, with the full products checked as above. In the smooth Thom--Mather case it holds by C2.05. Finiteness can be replaced by locally finite finite-depth data with locally finite charts and collars; an arbitrary non-locally-finite family does not define an ordinary finite-dimensional corner chart at an accumulation point.

\section{C.4. Pro-neighbourhoods, holinks, and cohomology}\label{c.4.-pro-neighbourhoods-holinks-and-cohomology}

\subsection{C4.01 --- Proved for finite metric HSS: the germ comparisons}\label{c4.01-proved-for-finite-metric-hss-the-germ-comparisons}

Use a finite stratification satisfying the frontier condition, with closed lower unions and locally closed strata \(S_i=X_i^*\). This is the finite-index scope explicitly adopted at /038--039. Assume the metric space is homotopically stratified in Quinn's sense. Whitney and Thom--Mather examples fall in this scope. Put \(O_i=X\setminus\bigcup_{k<i}S_k\), so \(S_i\) is closed and pure in the open subspace \(O_i\). Let \(V_i\) be the prospace of open neighborhoods of \(S_i\) in \(O_i\), ordered by shrinking. These neighborhoods are also open in \(X\). Finite intersections define \(V_d\).

Quinn, \cite[Proposition 3.2, p.469]{Quinn1988}, gives one fixed nearly strict, stratum-preserving neighborhood deformation \(h_i:N_i\times[0,1]\to O_i\), with \(h_i(x,0)=x\), \(h_i(x,1)\in S_i\), and \(h_i(s,t)=s\) for \(s\in S_i\). Until the final time it stays in the starting stratum. Reversing Quinn's interval convention gives the notation used here. The theorem applies to \(O_i\), whose induced stratification is again HSS: restrictions of pairwise tameness and holink fibrations give the defining conditions, or directly use Quinn 2.4 Corollary (2)--(3).

\textbf{Single-stratum germs.} Given an open \(U\supset S_i\), compactness of the time interval and continuity of \(h_i\) imply that \[
 W_U=\{x\in N_i:h_i(x,[0,1])\subset U\}
\] is open and contains \(S_i\). The fixed endpoint retraction \(r_i=h_i(-,1)\), restricted to \(W_U\), and the restricted homotopy prove that \(S_i\to V_i\) is a pro-homotopy equivalence. The same endpoint map is used at all levels, so these are compatible pro maps, not unrelated levelwise equivalences. This proves (51). If \(Z\) is a locally closed union of strata containing \(S_i\), the homotopy on \(W_U\cap Z\) stays in \(Z\), including its endpoint, proving both arrows of (56).

\textbf{Flag germs.} Let \(d=(i_0<\cdots<i_n=j)\), and fix target neighborhoods \(U_{i_k}\) for every flag entry. First shrink \(N_j\) so that its whole fixed homotopy stays in \(U_j\). For \(k<n\), metric normality and closedness of \(S_{i_k}\) in \(O_{i_k}\) give open neighborhoods \(U'_{i_k}\) with \(\overline{U'_{i_k}}^{\,O_{i_k}}\subset U_{i_k}\). Set \(A=\bigcap_{k<n}U_{i_k}\) and \(A'=\bigcap_{k<n}U'_{i_k}\). Because \(O_j\subset O_{i_k}\), one has \(\overline{A'}^{\,O_j}\subset A\). Let \[
 W=\{x\in N_j:h_j(x,[0,1])\subset A\}.
\] It is open and contains \(S_j\cap A\). Therefore \(F=\overline{A'}^{\,O_j}\setminus W\) is closed in \(O_j\) and disjoint from \(S_j\). Choose \(U'_j=N_j\setminus F\). Then \(A'\cap U'_j\subset W\), so its whole homotopy lies in \(A\cap U_j\). These restrictions of the one fixed \(h_j\) define a pro-homotopy retraction \[
 V_d\longrightarrow V_d\cap S_j.
\] It is the identity on \(S_j\), giving both inverse homotopies. If \(Z\) is a locally closed union of strata containing \(S_j\), the same homotopy restricted to \(Z\) stays there. The triangle through \(V_d\cap S_j\) proves (57), in fact under the weaker requirement \(S_j\subset Z\); the manuscript assumes \(X_j\subset Z\).

These are homotopy equivalences in prospaces, represented by maps and homotopies after cofinal shrinking; in particular they are isomorphisms in \(\operatorname{Pro}(\mathrm{HoTop})\). No assertion is made that taking a limit in ordinary spaces or topoi preserves them. A coherent pro-space formulation uses the same fixed homotopies, so all refinements come from restriction.

\textbf{Bands and (69).} Let \(J,L\subset I\) be convex bands satisfying \(J\prec L\): every member of \(J\) is below some member of \(L\), and every member of \(L\) is above some member of \(J\). The interval hull \[
 I'=\{i:\exists j\in J,\ell\in L,\ j\le i\le\ell\}
\] is convex; \(J\) is a lower set and \(L\) an upper set in \(I'\). For example, if \(j_0\le i\le j\) and \(j_0,j\in J\), convexity puts \(i\) in \(J\); the upper statement is dual. This proves the manuscript's order lemma, including uniqueness of the interval hull.

For a finite normal stratification choose a cofinal system of neighborhoods with \(U_i\cap U_k=\varnothing\) when \(i,k\) are incomparable. Indeed the two strata are disjoint closed subsets after removing both lower frontiers; normality separates them, and finitely many such shrinkings can be intersected. Thus any nonempty finite neighborhood intersection has flag labels. The space \(V_J\cap V_L\) is covered by the finitely many \(V_j\cap V_\ell\), \(j\in J,\ell\in L\). If \(j\le\ell\), the maximum label is in \(L\). If \(\ell<j\), then convexity and domination imply both labels belong to \(J\cap L\); again the maximum is in \(L\). The same remains true for every finite intersection of these pieces.

Apply the flag result to each such intersection, with \(Z=X_L^*\), which contains its largest stratum. It gives a pro-homotopy equivalence on every intersection after restricting to \(X_L^*\). Finite open-cover homotopy descent then gives \[
 V_J^L=V_J\cap X_L^*\ \longrightarrow\ V_J\cap V_L
\] as a pro-homotopy equivalence, proving (69). To justify descent concretely, use the finite poset of nonempty subsets of the finite cover-piece label set, with opposite inclusion for its intersection diagram. Its nerve has finitely many nondegenerate simplices. A partition of unity gives a homotopy inverse from the union to this finite homotopy colimit. The natural restriction map of intersection diagrams is an equivalence in prospaces at every object by the fixed-homotopy proof above; its inverses need not be chosen as separate natural maps. Finite homotopy colimits preserve this equivalence and commute with a synchronized cofiltered presentation: on mapping spaces this is the fact that filtered colimits commute with finite limits. Only finitely many source refinements are required. This argument does not truncate the unbounded ordinary \v{C}ech nerve, and would require a separate justification for an infinite index family.

This argument does not require a strictly commuting system of Mather retractions. The inherited R7 assertion that rho-control is necessary for this pro-homotopy statement is stronger than needed.

\subsection{C4.02 --- Proved: the relative holink comparison}\label{c4.02-proved-the-relative-holink-comparison}

% Independent proof supplied by /root/homotopy_levels, 2026-11-19.
% Replacement proof for C4.02; retain the existing section identifier.
For possibly overlapping dominated convex bands put
\(Y=X_J^*\), \(L_0=X_L^*\), and \(Z=Y\cup L_0\). Define the relative path space
\[
 P_{J,L}=\{\gamma\in C([0,1],Z):\gamma(0)\in Y,
                    \ \gamma(t)\in L_0\text{ for every }t>0\}.
\]
It contains \emph{all} such paths: paths starting in \(Y\cap L_0\) need not be constant. In particular it contains all constant paths on that intersection. Give it the compact-open topology. The set \(Y\) is a closed pure union of strata in the induced finite metric HSS \(Z\), so \cite[Proposition 3.2]{Quinn1988} supplies a fixed neighborhood \(N\) of \(Y\) and a nearly strict stratum-preserving deformation
\(h:N\times[0,1]\to Z\), with \(h_0=\mathrm{id}\), \(h_1=r:N\to Y\), and \(h_t|_Y=\mathrm{id}\). For \(t<1\), a point outside \(Y\) remains in its original stratum.

For each open neighborhood \(U\) of \(Y\) in \(Z\), let \(P_U\) be the paths in \(P_{J,L}\) whose whole image lies in \(U\). There is a natural zigzag
\[
 \operatorname{const} P_{J,L}
 \ \longleftarrow\ \{P_U\}_{U\supset Y}
 \ \xrightarrow{e_1}\ \{U\cap L_0\}_{U\supset Y}=V_J^L.
\]
The equality on the right means equality of pro-germs: restrictions of the unions of the finitely many individual stratum neighborhoods form a cofinal system of neighborhoods of \(Y\) in \(Z\). Indeed each such union is open and contains \(Y\); conversely an open neighborhood of \(Y\) in \(Z\) is the restriction of an ambient open set containing \(Y\), and intersecting that ambient open set with the prescribed open complements of the lower frontiers gives the required individual neighborhoods.

\textbf{The left arrow.} The path space is metrizable, using the uniform metric of any bounded compatible metric on \(Z\), and hence is paracompact. The open subsets
\[
 C_n=\{\gamma:\gamma([0,1/n])\subset U\},\qquad n\ge1,
\]
cover \(P_{J,L}\). Choose a locally finite partition of unity subordinate to a refinement of this cover; for each member of the refinement retain its integer \(n\). If its partition functions are \(\lambda_a\) and integers are \(n_a\), put
\(\epsilon(\gamma)=\sum_a\lambda_a(\gamma)/n_a\).
This is positive and continuous. Every active \(1/n_a\) is an admissible initial interval for \(\gamma\); since their finite maximum is also admissible, \(\gamma([0,\epsilon(\gamma)])\subset U\). Thus
\[
 S\gamma(u)=\gamma(\epsilon(\gamma)u)
\]
is a continuous map \(P_{J,L}\to P_U\). The homotopy
\[
 K_s\gamma(u)=\gamma\bigl((1-s+s\epsilon(\gamma))u\bigr)
\]
joins the identity to inclusion followed by \(S\), remains in \(P_{J,L}\), and restricts to a homotopy in \(P_U\) when \(\gamma\in P_U\). It fixes the initial point. Consequently every inclusion \(P_U\hookrightarrow P_{J,L}\) is a homotopy equivalence. These inclusions are a natural transformation, so the left arrow is a pro-homotopy equivalence. The endpoints are not fixed by this shortening homotopy; their coherent comparison in \(L_0\) is the explicitly displayed terminal segment of \(\gamma\).

\textbf{The right arrow.} For a target neighborhood \(U\), put
\[
 W_U=\{x\in N:h(x,[0,1])\subset U\}.
\]
Compactness of \([0,1]\) makes \(W_U\) an open neighborhood of \(Y\). Define
\[
 s_U:W_U\cap L_0\longrightarrow P_U,
 \qquad s_U(x)(u)=h(x,1-u).
\]
For \(x\in L_0\setminus Y\), the positive-time points of this reversed path remain in the stratum of \(x\), hence in \(L_0\). For \(x\in Y\cap L_0\) the path is constant. Thus it belongs to the relative path space in either case, and \(e_1s_U=\mathrm{id}\). All the maps \(s_U\) are restrictions of the same formula using the fixed \(h\); for finitely many targets a common smaller neighborhood gives literal agreement. They therefore define one pro map.

For the other composite, take \(\gamma\in P_{W_U}\) and define, for \(t>0\),
\[
 H_t(\gamma)(u)=
 \begin{cases}
 h\bigl(\gamma(t),1-u/t\bigr),&0\le u\le t,\\
 \gamma(u),&t\le u\le1,
 \end{cases}
 \qquad H_0(\gamma)=\gamma.
\]
The formulas agree at \(u=t\). Every positive-time point lies in \(L_0\), including when \(\gamma(t)\in Y\cap L_0\), because then \(h\) fixes that point. Every path lies in \(U\), and its endpoint is always \(\gamma(1)\). At \(t=1\) the path is \(s_U(e_1\gamma)\). Continuity away from \((t,u)=(0,0)\) is immediate. At that corner, \(\gamma(t)\to\gamma(0)\in Y\), and continuity of \(h\), compactness of its time interval, and \(h(\gamma(0),v)=\gamma(0)\) uniformly in \(v\) prove continuity. The same argument with \(\gamma\) varying gives continuity of the adjoint map on \(P_{W_U}\times[0,1]^2\), hence continuity into the compact-open path space. This proves the second inverse homotopy after shrinking. All refinements restrict this one fixed formula. The right arrow is consequently a pro-homotopy equivalence, in fact over \(L_0\).

The zigzag thus gives a canonical equivalence of pro-homotopy types. Its arrows literally commute with the endpoint maps to \(L_0\); the start map on the middle term is evaluation at zero. Under the inverse supplied by \(h\), the induced map to \(Y\) is \(r\), and the displayed \(H_t\) supplies its coherent comparison with the initial point of any representative path.

When \(J\cap L=\varnothing\), this relative path space is the conventional strict holink of \((Z,Y)\), recovering Quinn's usual formulation. When \(J=L\ne\varnothing\), that conventional strict holink is empty, whereas \(V_J^J=X_J^*\). The relative path space instead becomes \(C([0,1],X_J^*)\), which contracts by endpoint evaluation to the constant paths. Thus the literal strict-holink claim needs disjointness; the all-paths relative formulation above is valid with overlap.


\subsection{C4.03 --- Proved: the finite cover spectral sequence}\label{c4.03-proved-the-finite-cover-spectral-sequence}

Fix finite dominated bands, the interval hull \(I'\), and \[
 D=\{d\subset J\cup L:d\text{ a nonempty flag},\
             \min d\in J,\ \max d\in L\}.
\] Give \(D\) arrows \(d\to e\) when \(d\subset e\). Thus \(V_e\subset V_d\), and cohomological restriction is covariant in \(D\). For an abelian sheaf \(F\) on the neighborhood germ \(V_J^L\), interpret its cohomology as the filtered colimit of cohomology of representatives. There is a first-quadrant, strongly convergent spectral sequence \[
 E_2^{p,q}=R^p\!\lim_{d\in D}H^q(V_d^{X'},F)
 \ \Longrightarrow\ H^{p+q}(V_J^L,F).
\] This is the typed meaning of the category-cohomology notation in (73).

First choose cofinal representatives for which every upper-band neighbourhood \(U_l\) lies in the open upper union \(X_L^*\subset X'\). Their union is then exactly \(X_L^*\). Consequently the union of the pairwise cover pieces \(U_j\cap U_l\) is exactly \(V_J\cap X_L^*=V_J^L\) at that stage. All flags in \(D\) have maximum in \(L\), so their intersections are actual opens of this union. Thus no homotopy invariance of arbitrary sheaf cohomology is being invoked. Choose these representatives for the finite open cover and a flabby resolution of \(F\). Form the double complex whose degree \((p,q)\) is the product of \(\Gamma(V_{d_p},\mathcal I^q)\) over chains \(d_0\subset\cdots\subset d_p\), with the usual alternating restriction differential horizontally. Its horizontal augmentation computes sections on the union. Exactness can be checked locally inside a covering piece: adjoining that piece to an intersection gives the common refinement and contracts the augmented \v{C}ech complex; empty intersections contribute zero. More precisely use the augmented complex of sheaves obtained by direct image from the open intersections. On a covering open set, contraction is by adjoining its flag; a nonflag union has empty intersection and contributes zero. This local contraction proves exactness without identifying the stalk of an open direct image outside its open set with zero. With a flabby resolution, all the open direct images are flabby. Flabbiness makes all restriction terms acyclic for sections, so the total complex computes sheaf cohomology on the union.

Taking vertical cohomology first yields the derived-limit complex of the coefficient diagram \(d\mapsto H^q(V_d,F)\), proving the displayed \(E_2\)-page. Taking filtered colimits over representative neighborhoods is exact in abelian groups and commutes with the finite products here. Finally \(D\) has bounded chain length, so the horizontal filtration is finite; there is no convergence or product-completion ambiguity. Empty \(D\) gives the empty union and zero groups.

\subsection{C4.04 --- Refuted: HSS alone gives conical collars}\label{c4.04-refuted-hss-alone-gives-conical-collars}

Quinn's example \cite[p.452]{Quinn1988} is the shrinking hedgehog \[
 X=\bigvee_{n\ge1}[0,1/n],\qquad Y=\{0\},
\] with the metric \(|s-t|\) on one arm and \(s+t\) on distinct arms. It is compact metric. The radial contraction \(h((n,s),t)=(n,(1-t)s)\) makes \(Y\) tame, strictly off \(Y\) until the final time; the holink map to a point is a fibration. Thus the two-step filtration is homotopically stratified in Quinn's general definition.

It has no mapping-cylinder neighborhood of \(Y\). Every neighborhood contains all sufficiently short arms, including their terminal endpoints, and these endpoints converge to \(Y\). In a mapping-cylinder neighborhood with radial coordinate \(r\), a point with \(0<r<1\) lies on an embedded interval through itself, obtained by varying \(r\) in both directions. A terminal endpoint of an arm cannot be an interior point of any embedded interval in \(X\). Hence every such endpoint would have to lie on the outer level \(r=1\). Continuity is contradicted by their convergence to the apex \(r=0\). The same argument applies after making the cylinder smaller.

Therefore \(\Sigma_d\) in the actual collar/mapping-cylinder sense need not exist for a general HSS. Equations (51), (56), (57), (69), the holink comparison, and the finite spectral sequence have the independent pro-homotopy proofs above. An equality \(V_d\simeq\Sigma_d\) requires an actual compatible conical deployment, or must define a replacement \(\Sigma_d\) using the full iterated holink diagram. The example uses Quinn's general strata, not manifold strata; no claim is made here that it satisfies an additional manifold-strata hypothesis.

\subsection{C4.05 --- Proved with proper controlled collars: germ and face}\label{c4.05-proved-with-proper-controlled-collars-germ-and-face}

% Independent replacement proof supplied by /root/homotopy_levels, 2026-11-19.
% Replacement body for C4.05; retain its statement identifier and heading.
Use the actual finite proper conical deployment with compatible iterated fibre collars of C2.05, in the metric setting of C4.01. We spell out the radial data used in this assertion. Each stratum \(S_i\) has a fixed control tube \(N_i\), a retraction \(\pi_i:N_i\to S_i\), and a radius \(\rho_i:N_i\to[0,1)\) with zero set \(S_i\). Each closed subtube \(0\le\rho_i\le a<1\) is proper over \(S_i\times[0,a]\) under \((\pi_i,\rho_i)\). This follows from the proper compact link bundle in its mapping-cone presentation: the product of that bundle with \([0,a]\), followed by the cone quotient, is proper.

On the resolution \(\beta_j:\Sigma_j\to\overline{S_j}\), the interior is identified with \(S_j\). If the original construction retains a cut core instead, its collar stretching gives the same identification: append an outer collar with coordinate in \((-1,0]\), combine it with an inner collar \([0,a)\), and choose a monotone homeomorphism \((-1,a)\to(0,a)\) equal to the identity near \(a\). Extend by the identity outside the collar. At corners apply these reparametrizations in the commuting product coordinates. Retained outer faces in intermediate induction diagrams are handled by C1.B2a. The compatible iterated fibre collars give the following precise simultaneous coordinates. For
\(d=(i_0<\cdots<i_{n-1}<j)\), put
\[
 F_d=\bigcap_{k<n}\Sigma_{i_kj}=\Sigma_d,
 \qquad F_d^\circ=F_d\setminus
       \bigcup_{\ell\notin\{i_0,\ldots,i_{n-1},j\}}
                  (F_d\cap\Sigma_{\ell j}).
\]
Thus all remaining incoming boundary faces, including intermediate ones not selected by \(d\), are removed in \(F_d^\circ\). After fixing compatible collar widths, the intersection of the selected collar neighborhoods in the interior is identified with
\[
 c_d:F_d^\circ\times(0,a)^n\xrightarrow{\cong}
           S_j\cap\bigcap_{k<n}N_{i_k}.
\]
Here and below \(c_d\) includes the blowdown on the interior. The widths may first be normalized by the given collar coordinates. Their compatibility with the boundary fibrations gives continuous maps
\[
 p_k:F_d^\circ\times(0,a)^k\longrightarrow S_{i_k}
\]
with
\[
 \rho_{i_k}c_d(z,r)=r_k,
 \qquad
 \pi_{i_k}c_d(z,r)=p_k(z,r_0,\ldots,r_{k-1}).
 \tag{C4.05a}
\]
The absence of later radii in the second formula is essential: the later collars are vertical for the earlier source fibration. These are the simultaneous product and fibre-compatibility axioms of the actual resolved collar system, equivalently the compatible radial control constructed in C2.05. Pairwise commutativity of arbitrary retractions alone would not supply (C4.05a).

Here is the translation from the resolved data of C2.05. On the face \(\Sigma_{i_kj}\), the faces with labels \(i_\ell<i_k\) are the inverse images of the corresponding boundary faces of the base \(\Sigma_{i_k}\); this is the flag inverse-image identity. Their collars are carried to the base collars. The faces with labels \(i_\ell>i_k\) are vertical faces of the source fibre, and their compatible collars leave that source map fixed. Applying the commuting collar maps successively therefore carries the source only through the earlier coordinates, exactly as in (C4.05a). The transverse product-collar axiom identifies intersections of the chosen collar neighborhoods with the simultaneous collar of their face intersection, after shrinking the finite collar system. It also preserves every nonselected face. Removing all boundary faces to pass to \(S_j\) consequently removes the nonselected faces from the base \(F_d\), giving precisely \(F_d^\circ\), and requires each selected radius to be positive. This proves the displayed coordinate description from the full iterated fibre-collar data, rather than assuming it for arbitrary retractions.

\textbf{A cofinal family consisting of actual neighborhoods.} For every positive continuous function \(\epsilon_i:S_i\to(0,a)\), let
\[
 T_i(\epsilon_i)=
 \{x\in N_i:\rho_i(x)<\epsilon_i(\pi_i(x))\}.
\]
These are genuine open neighborhoods of \(S_i\) in the required open complement of its lower frontier. They form a cofinal family. Indeed, for a prescribed neighborhood \(U_i\), choose a fixed smaller closed subtube \(K_i=\{\rho_i\le a'\}\), \(0<a'<a\). The set \(K_i\setminus U_i\) is closed; properness makes its image under \((\pi_i,\rho_i)\) a closed subset of \(S_i\times[0,a']\) disjoint from the zero section. There is a positive continuous function \(\epsilon_i<a'\) whose subgraph misses that image. To construct it explicitly, choose a cover of \(S_i\) by open sets \(B_\alpha\) and numbers \(b_\alpha>0\) for which \(B_\alpha\times[0,b_\alpha]\) misses the closed image. A locally finite subordinate partition of unity gives \(\epsilon_i(y)=\frac12\sum_\alpha\lambda_\alpha(y)b_\alpha\), capped at \(a'/2\). Each active bound is admissible at \(y\), hence so is their weighted average. Then \(T_i(\epsilon_i)\subset U_i\).

Consequently the germ \(V_d\cap S_j\) is represented by the actual intersections
\[
 H_\epsilon=S_j\cap\bigcap_{k<n}T_{i_k}(\epsilon_{i_k}).
\]
The last neighborhood, that of \(S_j\) itself, imposes no restriction after intersecting with \(S_j\). In the fixed simultaneous coordinates, (C4.05a) identifies \(H_\epsilon\) with
\[
 \left\{(z,r):
 0<r_k<\epsilon_{i_k}\bigl(p_k(z,r_0,\ldots,r_{k-1})\bigr)
 \quad(0\le k<n)\right\}.\tag{C4.05b}
\]
These are recursively varying radius regions. The family is cofinal in both required senses: every member is exactly an actual intersection of the indicated neighborhood representatives, and the preceding properness argument puts some such member inside every prescribed tuple of neighborhoods. We do \emph{not} claim that arbitrary positive functions on the entire resolved face or on fibre directions produce cofinal ambient neighborhood systems.

\textbf{The homotopy equivalence and its compatibility.} Each region (C4.05b) is homeomorphic to \(F_d^\circ\times(0,1)^n\). The homeomorphism is given successively by
\[
 r_k=\epsilon_{i_k}\bigl(p_k(z,r_0,\ldots,r_{k-1})\bigr)u_k,
 \qquad 0<u_k<1.
\]
Its inverse divides by the same strictly positive continuous functions, successively in increasing \(k\). Projection to \(z\) is therefore a homotopy equivalence. All the transition inclusions for smaller control functions commute with this one fixed projection. We obtain a natural objectwise equivalence from the cofinal diagram of the \(H_\epsilon\) to the constant space \(F_d^\circ\), and hence an equivalence of prospaces, with full coherent compatibility. There is no inference merely from unrelated isomorphisms in \(\mathrm{HoTop}\).

Finally \(F_d^\circ\hookrightarrow F_d\) is a homotopy equivalence. The finitely many remaining commuting collars can be pushed inward, using in each coordinate a function positive at zero and equal to that coordinate outside a smaller collar. They preserve the selected faces because their collars commute. The resulting global homotopy has image in \(F_d^\circ\) at its final time and preserves \(F_d^\circ\) at every time. Thus
\[
 V_d\cap S_j\simeq F_d^\circ\simeq F_d=\Sigma_d.
\]
C4.01 supplies \(V_d\simeq V_d\cap S_j\). Empty face intersections have an empty cofinal neighborhood intersection; the same statement then holds with the empty space. For a one-element flag the assertion reduces to the inward-collar equivalence between the resolved piece and its stratum.

The rectangular version must be qualified. An arbitrary neighborhood of an intermediate stratum can impose \(r_1<r_0^2\) near a lower stratum, so no product rectangle with two fixed positive radii is contained in it. Formula (C4.05b), with radius functions pulled back from the actual stratum bases, is what makes the comparison valid. The resulting equivalence is canonical in the homotopy category for the fixed compatible deployed collar germs; the contractible choices of smaller radius functions and the displayed natural projection specify its compatible homotopies.


\subsection{C4.06 --- Proved: the canonical normal-crossings comparison}\label{c4.06-proved-the-canonical-normal-crossings-comparison}

The notation here is checked in the original scans 147/173--177 \cite{C_Grothendieck147}. Their notation is a filtered normal bundle over \(B=D_{d_*}^*\), with a pair \(d_0''\le d_0'\) of entries of the flag and an exact sequence \[
 0\longrightarrow A=\check N_{d_*,d_0'}
 \longrightarrow E=\check N_{d_*,d_0''}
 \xrightarrow{q}Q=\check N_{d_*,(d_0',d_0'')}
 \longrightarrow0.
\] This is (286)--(287). The normal-crossings branches give \(Q\) a local decomposition into complex lines, with transitions allowed to permute those lines. Write \(Q^*\) for the complement of all coordinate hyperplanes. Thus its fibre is \((\mathbb C^*)^r\), not \(\mathbb C^r\setminus\{0\}\). For rank zero, the empty coordinate arrangement has complement the whole zero bundle; the torus fibre is a point. The same argument then applies without a puncture, and an empty base gives the empty equivalence. Page175(290) defines \(E^*=q^{-1}(Q^*)\); (291) describes \(E^*\to Q^*\) as a torsor for the pulled-back additive bundle \(A\).

The quotient map \(q:E^*\to Q^*\) is a canonical homotopy equivalence. Choose a Hermitian splitting \(E=A\oplus Q\); then \((a,v)\mapsto((1-t)a,v)\) contracts the kernel and proves the claim. Splittings form an affine space, so their homotopies identify all inverse choices. The map \(q\) itself does not depend on a splitting.

The sector tube \(V^{d_0'}_{d_*,d_0''}\) is, by definition of the normal-crossings incidence sector, the neighborhood germ of \(B\) in the \(E\)-directions, with precisely the coordinates belonging to \(Q\) required to be nonzero. Its pro-homotopy type has a canonical comparison with \(E^*\). Here is a construction and a check of its ambiguity. For a simple normal crossings divisor, use its divisor line bundles with defining sections. Choose an ordinary smooth tubular projection onto the intersection \(B\) and connections on those line bundles. Parallel transport identifies their restrictions on a small neighborhood with their pullbacks from \(B\); the tuple of defining sections then has differential the normal-coordinate isomorphism along \(B\). The inverse function theorem, after variable-radius shrinking, gives an adapted tubular germ from the direct sum of the normal lines to the original space. It takes every coordinate hyperplane and every flag subbundle to the corresponding divisor incidence. This constructs the required sector comparison.

Normalize each adapted tubular germ to have identity derivative on the intrinsic normal bundle, by composing with the inverse of its fibrewise normal derivative. This preserves the coordinate arrangement and the flag subbundles. Arbitrary normal-bundle automorphisms are not being declared homotopic to the identity. If a kernel block has no chosen line decomposition, retain its coordinates from an ordinary tubular chart and use the defining sections in the quotient block; the same inverse-function argument respects every flag subbundle.

If \(\Phi,\Psi\) are two adapted tubular germs with the same normal derivative, their transition \(g=\Psi^{-1}\Phi\) fixes the zero section and has identity normal derivative. Fibre dilations give \[
 g_t=\delta_{1/t}\,g\,\delta_t\quad(t>0),\qquad g_0=\mathrm{id}.
\] Taylor's formula in normal coordinates shows this is a continuous, indeed smooth, family of germs through \(t=0\); its tangential displacement tends to zero and its normal derivative tends to the identity. Every \(g_t\) preserves the coordinate-hyperplane arrangement. Restricting to sufficiently small variable-radius neighborhoods makes the family defined for all \(t\in[0,1]\). The same construction works for parameter families. Hence the induced comparison on the sector has a unique homotopy class, compatible with refinements and changes of tubular choices. This supplies the canonical meaning that the manuscript explicitly leaves unresolved below (294).

The small-sector pro-germ also compares explicitly with the entire punctured bundle. For a bundle norm and a positive continuous radius function \(\eta\) on the base, the inclusion \(\{v\in E^*:\|v\|<\eta(\pi v)\}\hookrightarrow E^*\) is a homotopy equivalence. An inverse up to homotopy is \[
 v\longmapsto \frac{\eta(\pi v)}{\eta(\pi v)+\|v\|}v.
\] Interpolation of its strictly positive scalar factor with \(1\) supplies both inverse homotopies and preserves every coordinate nonvanishing condition. Such variable-radius disk neighborhoods are cofinal by the tubular chart and paracompactness of the base. Their natural inclusions therefore identify the sector prospace with the constant space \(E^*\), coherently under shrinking.

Combining it with \(q\) gives the precise statement \[
 V^{d_0'}_{d_*,d_0''}
 \ \simeq\
 \left.\check N_{d_*,(d_0',d_0'')}^{\,*}\right|_{D_{d_*}^*}.
\] The ambiguity is homotopy, not a distinguished homeomorphism of neighborhoods. A Hermitian norm further retracts each punctured line to its unit-circle bundle, giving the corresponding torus bundle; metrics form a convex space of choices.

For normal crossings with branch monodromy, perform the construction on the finite local branch-label cover. The normal line bundles and the quotient map descend equivariantly, as does the homotopy-class comparison, because the coordinate arrangement and the dilation are invariant under branch permutations. On a Deligne--Mumford orbifold chart finite-group averaging supplies invariant auxiliary metrics and connections. Thus the result is canonical in the incidence/orbifold category. Passing to an unlabelled orbit-poset and discarding stabilizers would lose data.

\section{C.5. The consequence margins in TOP, PL, and smooth}\label{c.5.-the-consequence-margins-in-top-pl-and-smooth}

\subsection{C5.01 --- Proved: the component order under local path connectivity}\label{c5.01-proved-the-component-order-under-local-path-connectivity}

Let P be a Hausdorff locally path-connected topological poset, with closed order R\ensuremath{\subset}P\ensuremath{\times}P and projections s,b:R\ensuremath{\to}P. Assume b is a proper local immersion and s a proper Hurewicz fibration. The diagonal is clopen in R: it is closed by Hausdorffness and open by local injectivity of b at identities. Then distinct comparable points belong to different path components; the induced relation on \ensuremath{\pi}\ensuremath{{}_{0}}(P) is a partial order. If every lower interval is totally ordered, the target map on arrows between any two fixed components is a closed embedding.

\textbf{Proof.} Fibres of b are discrete and compact, hence finite. Put N(y)=\textbar b\ensuremath{{}^{-1}}(y)\textbar, including the identity. This is locally bounded. Choose neighborhoods U\ensuremath{{}_{1}},\ldots,U\_N of the points over y on which b is injective. The image of the closed complement of their union is closed and avoids y. In its open complement W every fibre has at most N points.

For fixed x the upper interval P\_\ensuremath{\ge}x=b(s\ensuremath{{}^{-1}}(x)) is compact. A finite subcover of the preceding neighborhoods bounds N on it. Suppose x\textless y lie in one path component C. Lift through s any path from x to another point of C, starting at the arrow x\textless y. Its target remains in C, and its arrow remains strict because identities form a clopen subset. Hence every point of C has a strict successor in C. Construct arbitrarily long chains x=x\ensuremath{{}_{0}}\textless x\ensuremath{{}_{1}}\textless\ldots\textless x\_n.~Their endpoints lie in P\_\ensuremath{\ge}x by transitivity but have N(x\_n)\ensuremath{\ge}n+1, a contradiction.

For component arrows C\ensuremath{\to}D and D\ensuremath{\to}E, lift a path in D from the source of the second to the target of the first, giving composable representatives. Thus the component relation is transitive. If C\ensuremath{\le}D\ensuremath{\le}C, the resulting composite lies within C, so its endpoints coincide; antisymmetry in P then identifies the middle point as well, proving C=D.

Finally, two arrows from C to one target z have sources in the totally ordered lower interval of z. The sources are comparable and in one component, hence coincide. Thinness then identifies the arrows. The restricted target map is proper, because components are clopen, and a proper injective map to a Hausdorff space is a closed embedding. This proves all assertions.

This repairs the ordered-space application of 150/070--071 without the questionable ``rare subset lowers dimension'' step on /069. It does not prove the same statement for an arbitrary category with nontrivial idempotents; the order hypothesis is used in forming arbitrarily long strict chains.

\subsection{C5.02 --- Refuted in all three categories: automatic transversality}\label{c5.02-refuted-in-all-three-categories-automatic-transversality}

Use \(I=\{0<1<2\}\), \[
 \Sigma_0=\{*\},\quad
 \Sigma_1=\mathbb R\times[0,\infty),\quad
 \Sigma_2=\mathbb R\times[0,\infty)^2.
\] Write \((x,t,r)\) on the last space and take \[
 \Sigma_{01}=\mathbb R\times\{0\},\quad
 \Sigma_{12}=\{r=0\},\quad
 \Sigma_{02}=\{t=0\}.
\] All targets are face inclusions and \(s_{01},s_{02}\) are constant. Put \[
 s_{12}(x,t,0)=(xt,t)
\] for the smooth example. For a PL example replace \(xt\) by \(f(x,t)=\max(-t,\min(x,t))\), a finite piecewise-linear function.

The incoming faces are closed, properly embedded, and transversely collared. Their intersection is \(\{t=r=0\}=s_{12}^{-1}\Sigma_{01}\), and projection control holds because both maps to \(\Sigma_0\) are constant. Thus the induced relation is a topological order with finite totally ordered lower intervals; every bare algebraic and face condition in the proposed implication holds. The source is not assumed fibrant in the general Scholium or its /062 margin.

Yet \(s_{12}\) is not transverse to the base boundary in the collar-product meaning explicitly used at /025--027 and /031. Its boundary restriction maps the whole \(x\)-axis to \((0,0)\). If a collar-product description existed near that axis, the image of a small collar neighborhood would lie in a one-dimensional collar through that point. In either example the image of every sufficiently small neighborhood of \((0,0)\) contains a nonempty open two-dimensional wedge: in the smooth example \(|u|<\delta t\), and in the PL example \(|u|<t\). A one-dimensional embedded collar cannot contain an open subset of a two-dimensional manifold, by invariance of dimension. Thus the specified collar-product transversality fails even topologically. This terminology must be retained: the smooth map does satisfy ordinary differential transversality to the target boundary, which is a weaker requirement here.

A sufficient extra hypothesis is local triviality of the entire boundary diagram over the base collar, with its trivializations preserving all its faces. Pulling that diagram back along the collar's product coordinate then gives compatible product collars. This is the strong equisingularity qualification in /039; an unspecified use of ``fibrant'' must not be substituted for it.

Even adding raw proper Hurewicz fibrancy does not imply the smooth product-collar condition. Take \(\Sigma_0=\{*\}\), \(\Sigma_1=[0,1]\), \(\Sigma_2=[0,1]^2\), \(\Sigma_{01}=\{0,1\}\), \(\Sigma_{12}=[0,1]\times\{0,1\}\), and \(\Sigma_{02}=\{0,1\}\times[0,1]\). Let \(s_{12}(u,v)=u^2\) and let the other sources be constant. Every source is proper and Hurewicz: on each component \(u\mapsto u^2\) is a homeomorphism. All target collars, intersections, and projection-control identities hold. But the derivative of \(s_{12}\) in the normal direction at \(u=0\) is zero, invariant under smooth collar changes; a smooth product over that collar has nonzero normal derivative. Thus the smooth source transversality fails. The assertion in the inherited reading that the written fibrancy hypothesis itself supplies local triviality must use its stronger, category-specific bundle interpretation.

\subsection{C5.03 --- Refuted in TOP: automatic manifold fibres}\label{c5.03-refuted-in-top-automatic-manifold-fibres}

Let \(H\) be the Poincare homology three-sphere and \(F=\Sigma H\) its suspension. Then \(F\) is a compact polyhedron which is not a topological manifold: near either suspension vertex the punctured neighborhood is \(H\times(0,\epsilon)\), whose nontrivial fundamental group cannot become that of a punctured four-ball. More precisely the nested punctured cone neighborhoods all induce isomorphisms on that nontrivial group, whereas a four-manifold has a basis of punctured simply-connected ball neighborhoods.

Cannon's double-suspension theorem gives \(\Sigma^2H\cong S^5\). Removing the two outer suspension points gives \(F\times\mathbb R\cong S^4\times\mathbb R\): the image of those two points can be sent to the north and south poles by a homeomorphism of \(S^5\). Consequently \(E=F\times S^1\) is a closed topological five-manifold, since it is locally \(F\times\mathbb R\).

Take \(\Sigma_0=S^1,\ \Sigma_1=E\times[0,1]\), with incoming face \(\Sigma_{01}=E\times\{0,1\}=\partial\Sigma_1\), and source projection to \(S^1\). This is a proper locally trivial topological bundle, with fibre \(F\amalg F\). Both object spaces are manifolds and the target is the whole boundary. All triple, order, and transversality conditions are vacuous. The fibre is not a manifold. This disproves the mixed-fibre implication at /051 and /063 even with genuine bundles, not merely homotopy fibrations.

The precise primary statements used here are Cannon \cite[Theorem 1.5, p.833; Corollary 3.6, p.838]{C_Cannon}.

\subsection{C5.04 --- Refuted in TOP: automatic manifold faces}\label{c5.04-refuted-in-top-automatic-manifold-faces}

Use the same compact polyhedron \(F=\Sigma H\) as in C5.03, and put \[
 \Sigma_0=\{*\},\qquad \Sigma_1=[0,1],\qquad
 \Sigma_2=F\times[0,1]^2.
\] Let \(\Sigma_{01}=\{0,1\}\), \(\Sigma_{12}=F\times[0,1]\times\{0,1\}\), and \(\Sigma_{02}=F\times\{0,1\}\times[0,1]\). In coordinates \((z,u,v)\), the source \(s_{12}\) is projection to \(u\), and \(s_{01},s_{02}\) are constant to the point.

The square is locally a plane at its interior points and a half-plane at every boundary point, including its vertices. Since \(F\times\mathbb R\) is a five-manifold by C5.03, \(\Sigma_2\) is a compact topological six-manifold with boundary. Its whole boundary is \(F\times\partial([0,1]^2)\cong F\times S^1\), itself a five-manifold. It is the union of the two specified collective faces. These faces have their evident transverse product collars, and \[
 \Sigma_{12}\cap\Sigma_{02}
 =F\times\{0,1\}^2=s_{12}^{-1}(\Sigma_{01}).
\] The source \(s_{12}\) respects the base collars, and projection control holds because the other two source maps are constant. Every object \(\Sigma_i\) is a manifold with boundary. All source maps are now proper Hurewicz fibrations: \(s_{12}\) is a product bundle with compact fibre \(F\amalg F\), and the other targets are points. Thus even adding equisingularity does not repair this TOP face assertion.

Each incoming face is a disjoint union of two copies of \(F\times[0,1]\), which is not a manifold with boundary. At positive interior interval coordinates the space is a five-manifold. At an endpoint its local homology vanishes in all degrees, since the relative local homology of the half-interval factor vanishes. Hence, if the product were a five-manifold with boundary, its intrinsic boundary would be exactly \(F\times\{0,1\}\). The boundary of a manifold with boundary is a manifold, contradicting the proved singularity of \(F\). The full family therefore satisfies the surrounding order, properness, source-fibrancy, target-collar and source-transversality conditions while its incoming faces are not manifolds. This refutes both the /067 margin and the version of /051(d) that tries to deduce manifold flag spaces from c) and the object-manifold part of d).

\subsection{C5.05 --- Proved for PL bundles and smooth submersions}\label{c5.05-proved-for-pl-bundles-and-smooth-submersions}

Work with the entire completed incidence system of C2.02, its transverse target collars, and its full boundary decomposition. Impose the following regularity on every source map. The positive PL and smooth assertions require these category-specific meanings of fibrancy. In PL assume the source is a locally trivial PL bundle and every collar and transition is PL. In the smooth case assume every source is a proper smooth submersion with the stated product behavior on horizontal faces. Together with full projection control this gives the vertical-face regularity proved below; alternatively one may start with the corresponding smooth bundles of corner manifolds. A map which is merely smooth and a Hurewicz fibration need not have the required differential regularity.

First, in PL, a locally finite polyhedron \(F\) such that \(F\times\mathbb R^m\) is a PL manifold is itself a PL manifold. Use a product triangulation and choose a top-dimensional simplex in the second factor. The link of its interior product with a vertex of \(F\) is the link of that vertex in \(F\). In a PL manifold that link is a PL sphere or ball, according as the point is interior or boundary. The link criterion gives the conclusion, at every vertex and then throughout \(F\). Applying this in a local bundle trivialization over an interior base point proves that the fibre is a PL manifold with boundary/corners; local triviality identifies the other fibres with it.

In the smooth case first verify tangential surjectivity on every vertical face, using the whole source-map system. On \(\Sigma_{ikj}\), with \(i<k<j\), projection control gives \[
 s_{ij}|_{\Sigma_{ikj}}
   =s_{ik}\circ(s_{kj}|_{\Sigma_{ikj}}).
\] The horizontal pullback and product-collar condition makes \(s_{kj}|_{\Sigma_{ikj}}:\Sigma_{ikj}\to\Sigma_{ik}\) a submersion, while \(s_{ik}\) is another source submersion. Their composite is therefore a submersion. Iterating this argument on flags gives surjectivity on every intersection of vertical faces. Thus one can choose horizontal lifts tangent to all those faces; in a product corner chart take such lifts simultaneously and combine them with a partition of unity. Horizontal base-boundary coordinates are separated by the specified product collars. The corner submersion theorem now gives fibre charts by fixing the base coordinates, so the fibre is a manifold with corners. Properness allows these lifts to be integrated over a common base neighbourhood, giving compatible local trivializations. An isolated smooth submersion lacking these vertical-face checks would not justify that conclusion.

In both categories a local chart on an arrow space \(\Sigma_{ij}\) therefore separates into base coordinates, horizontal base-boundary coordinates, fibre coordinates, and vertical fibre-boundary coordinates. The horizontal faces are exactly \(\Sigma_{kij}\), \(k<i\), by the inverse-image identity; the vertical faces are exactly \(\Sigma_{ikj}\), \(i<k<j\), by the full boundary decomposition and the ordering of the incoming labels. Thus \[
 \partial\Sigma_{ij}
 =\bigcup_{k<i}\Sigma_{kij}
   \ \cup\!
   \bigcup_{i<k<j}\Sigma_{ikj},
\] while the intrinsic boundary of each source fibre consists only of the latter, intermediate-index pieces. These formulas include all corners through their intersections. This proves the mixed-fibre and face-boundary conclusions of /051, /063, and /067 under these precise PL and smooth hypotheses.

For /067 alone, no source bundle is needed if the incoming collars and all ambient spaces are PL/smooth: the same local link criterion in PL, or the smooth collar chart in the smooth category, proves that each transverse face intersection is a manifold with corners. The boundary formula follows by deleting exactly the other incoming faces through each point. The TOP counterexample C5.04 explains why this cancellation step cannot be copied into bare topology.

For /062, the implication from the bare order and face axioms is false in every category, by C5.02. Strong local triviality of the entire face diagram over the base collars is a sufficient correction in every category. For /070--071 in the ordered-space situation, C5.01 proves the claimed component order and makes the extra monomorphism assumption redundant whenever components are path-connected and open (in particular for manifolds and locally path-connected PL spaces), in all three categories, without the manuscript's finite-dimensionality hypothesis.

\subsection{C5.06 --- Refuted in PL and smooth: raw fibrancy suffices}\label{c5.06-refuted-in-pl-and-smooth-raw-fibrancy-suffices}

The full construction and proof, independently audited, follows. It gives a closed six-manifold mapping properly to a circle with one nonmanifold fibre; in the smooth model there is one critical point.

% Independent contribution by /root/hypercovers_finite_sets, 2026-11-19.
% Suggested integration as C5.06--C5.07. Citation keys are C5H-prefixed.
\label{sec:C5H}

\paragraph{Scope and conclusion.}
The mixed assertions of 150/051 and 150/063 are false if
``fibrant'' means only Hurewicz fibration, even when all the maps are PL,
or, separately, all the maps are smooth, and all object spaces are compact
manifolds with boundary.  The source map in the counterexample below is
proper.  The positive statements for PL bundles and smooth source systems with the face regularity proved in C5.05
therefore require their explicitly stronger hypotheses.  This counterexample
does not affect the separate boundary-face assertion on 150/067.

\paragraph{The counterexample (proved).}
There exist a closed PL six-manifold $E_{\mathrm{PL}}$ and a proper
surjective PL map
\[
 p_{\mathrm{PL}}:E_{\mathrm{PL}}\longrightarrow S^1
\]
which is a Hurewicz fibration and has a fibre which is not a topological
manifold.  There also exist a closed smooth six-manifold $E_{\mathrm{sm}}$
and a smooth map $p_{\mathrm{sm}}:E_{\mathrm{sm}}\to S^1$ with the same
properties.  The smooth example has just one critical point.  The two
examples are constructed separately; the smoothing homeomorphism is not
asserted to be PL.

\paragraph{A contractible five-manifold with nontrivial boundary group.}
Take a smooth compact contractible Mazur four-manifold $A$ such that
$G=\pi_1(\partial A)\ne1$.  An opened primary construction is Zeeman,
\emph{On the dunce hat}, printed p.~347, the passage preceding Theorem~(5),
which describes the handle construction, records the nontrivial boundary
group, and gives the original locator Mazur~1961, p.~227
\cite[p.~347]{C5HZeeman1964}.  Here Zeeman's dotted-manifold notation denotes
the boundary.  The handle attachment can be made smooth and corners rounded.
Choose a smoothly embedded closed three-disc $D^3\subset\partial A$ and set
\[
 W=(A\times S^1)\mathop{\cup}_{D^3\times S^1}(D^3\times D^2),
 \qquad H=\partial W.
 \tag{C5H.1}
\]
Round the corners of this attachment.  Van Kampen gives
$\pi_1(W)=\mathbb Z*_{\mathbb Z}1=1$, since the attaching circle generates
$\pi_1(A\times S^1)$.  In the reduced Mayer--Vietoris sequence the only
possibly nonzero positive homology map is the isomorphism
$H_1(D^3\times S^1)\to H_1(A\times S^1)$.  Thus $W$ is acyclic.
A simply connected acyclic finite CW complex is contractible, by the
Hurewicz and Whitehead theorems, so $W$ is contractible.

The boundary is
\[
 H=\bigl((\partial A\setminus\operatorname{int}D^3)\times S^1\bigr)
       \mathop{\cup}_{S^2\times S^1}(S^2\times D^2).
 \tag{C5H.2}
\]
Removing a three-disc does not change the fundamental group of the closed
three-manifold $\partial A$.  The second van Kampen calculation consequently
reads
\[
 \pi_1(H)=(G\times\mathbb Z)*_{\mathbb Z}1\cong G\ne1.
\]
This construction is a boundary spin, not stabilization by an interval.
In particular it does not make the invalid inference that
$\partial(A\times I)$ retains $\pi_1(\partial A)$.

Write $M=DW$ for the double of $W$.  This is a simply connected smooth
homology five-sphere.  More specifically, the rounded manifold $W\times I$
is contractible and its boundary is $DW$.  Apply the smooth $h$-cobordism
theorem after deleting an interior six-disc: the resulting six-dimensional
cobordism has simply connected boundary components $S^5$ and $DW$, and their
inclusions are homology equivalences.  It is a product.  Hence
\[
 W\times I\cong D^6,\qquad DW\cong S^5.
 \tag{C5H.3}
\]
The dimension restriction is satisfied.  This is exactly the implication
``contractible, simply connected boundary $\Rightarrow$ smooth disc'' in
Milnor's Proposition~A, pp.~108--109, following Theorem~9.1, p.~107
\cite[\S9]{C5HMilnor1965}.

\paragraph{The quotient map and its exceptional fibre.}
Use one copy of $W\subset M=DW$ and fix $0\in S^1$.  Form
\[
 E=(M\times S^1)/(W\times\{0\}),\qquad
 p([x,t])=t,
 \qquad Y=M/W.
 \tag{C5H.4}
\]
Only the entire set $W\times\{0\}$ is identified to one point $v$; all
other equivalence classes are singletons.  The quotient is compact
metrizable.  All fibres except $p^{-1}(0)$ are copies of $M$, whereas
$p^{-1}(0)=Y$.  An outward collar of $H=\partial W$ identifies a neighbourhood
of the image of $W$ in $Y$ with an open cone on $H$.
The local homology there is that of an interior point of a five-manifold,
because $H$ is a homology four-sphere.  But the punctured conical
neighbourhoods have fundamental group $G\ne1$, and all inclusions between
smaller and larger punctured conical neighbourhoods induce isomorphisms.
This contradicts local simple connectivity at an interior point of a
five-manifold.  The local homology also excludes a boundary point.  Thus
$Y$ is not a topological manifold, with or without boundary.

\paragraph{The Hurewicz lifting property: complete control argument.}
We prove this for the ordinary quotient (C5H.4), without assuming its
collapse map $f:M\to Y$ is PL.  A small closed collar enlargement $A_\eta$ of
$W$ is a compact contractible manifold, and
\[
 B_\eta=A_\eta/W\cong C(H)
\]
is a compact polyhedron.  Choose a contraction of $W$ within itself.
The pair $(A_\eta,W\cup\partial A_\eta)$ is a cofibration pair, since all
these subspaces have compatible collar triangulations.  Extend the
contraction, together with the stationary homotopy on $\partial A_\eta$,
to a homotopy $H_t:A_\eta\to A_\eta$ starting at the identity.  It satisfies
$H_t(W)\subset W$ for all $t$ and $H_1|W$ is constant.  Therefore $H_1$
factors through a map $s_\eta:B_\eta\to A_\eta$, fixed on the outer
boundary.  Extend $s_\eta$ outside $B_\eta$ by the inverse of $f$, obtaining
$s_\eta:Y\to M$.

The homotopy $H_t$, extended by the identity outside $A_\eta$, gives
$\mathrm{id}_M\simeq s_\eta f$, supported in $A_\eta$.  Because
$H_t(W)\subset W$ at every time, $fH_t$ descends to a homotopy
$\mathrm{id}_Y\simeq f s_\eta$, supported in $B_\eta$.  Both are fixed on
their outer boundaries.  These are the required controlled inverse
homotopies; they use only the homotopy extension property of a finite
polyhedral pair.

Equip $E$ with any compatible metric.  Given $\epsilon>0$, first choose
$\eta$ and then $\delta>0$ so small that
\[
 \operatorname{diam}\bigl(q(A_\eta\times[-\delta,\delta])\bigr)<\epsilon,
 \tag{C5H.5}
\]
where $q:M\times S^1\to E$ is the quotient and the signed coordinate is
chosen in a small arc about $0$.  Such choices exist because these compact
sets shrink to $q(W\times\{0\})=\{v\}$.  Also arrange
\[
 \sup_{x\in M,\ |t|\le\delta}
    d\bigl(q(x,t),q(x,0)\bigr)<\epsilon,
\]
by uniform continuity of $q$.  The maps between the fibres at $0$ and $t$
are $f$ and $s_\eta$, placed in these respective levels.  Outside the
supports they move points less than $\epsilon$ by the latter inequality;
inside the supports they and both inverse homotopies move points less than
$\epsilon$ by (C5H.5).  At every nonzero base point, ordinary product
identifications of nearby fibres give the same property, with stationary
inverse homotopies.

Thus $p$ is strongly regular in the precise sense of Ferry's definition on
p.~373.  Its fibres are compact ANRs: $M$ is a manifold and $Y$ is a finite
polyhedron, obtained by attaching the cone on $H$ to the complementary
collared manifold.  The base $S^1$ is complete and finite dimensional.
Ferry's Theorem~1, pp.~373--374, therefore proves that $p$ is a Hurewicz
fibration \cite{C5HFerry1978}.  This establishes the full lifting property,
not merely a quasifibration or a statement about homotopy-equivalent fibres.

\paragraph{A PL total space and a PL map.}
Triangulate the smooth pair $(M,W)$ compatibly and choose a PL outward
collar $H\times[0,a]$.  In $M\times(-2a,2a)$ take the polyhedral
neighbourhood
\[
 N=(W\times[-a,a])\ \cup\
   \{((z,s),t):z\in H,\ 0\le s\le a,\ s+|t|\le a\}.
 \tag{C5H.6}
\]
It is a genuine neighbourhood of $W\times\{0\}$, including an outward
collar; the inadequate set $W\times[-a,a]$ alone would touch $W\times\{0\}$
at its boundary.  The boundary $L=\partial N$ is the double of $W$ with an
extra collar, hence is a PL five-sphere by (C5H.3).  Let
$h:L\to[-a,a]$ be the restriction of the coordinate $t$, a PL function.
Replace $N/(W\times\{0\})$ by the abstract cone $C(L)$ and put
\[
 \widetilde p(r,\theta)=r h(\theta),\qquad \widetilde p(v)=0.
 \tag{C5H.7}
\]
This is the linear cone extension of a PL function and is PL.  Attach the
cone to $(M\times S^1)\setminus\operatorname{int}N$ by the identity on $L$,
keeping the original projection outside.  Since $C(L)$ is a PL six-disc,
the resulting space $E_{\mathrm{PL}}$ is a closed PL six-manifold and
$p_{\mathrm{PL}}$ is PL.

Here is an exact height-preserving identification with (C5H.4).  For a
boundary point $(x,\pm a)$ with $x\in W$, send
\[
 (r,(x,\pm a))\longmapsto q(x,\pm ar).
\]
For a side-boundary point $((z,s),t)$ with $s+|t|=a$, send
\[
 (r,((z,s),t))\longmapsto q((z,rs),rt).
\]
The formulas agree on their common faces and send the whole cone vertex
to $v$.  Every point off $v$ has a unique preimage: use $r=|t|/a$ over $W$,
and $r=(s+|t|)/a$ on its outward collar.  The continuous bijection between
compact Hausdorff spaces is a homeomorphism, identity on $L$, and it
preserves height exactly.  Extending it outside $N$ identifies
$p_{\mathrm{PL}}$ with $p$.  Hurewicz lifting and properness therefore pass
to the PL model.

\paragraph{A smooth total space and a smooth map.}
Choose a smooth outward collar coordinate $s$ and smooth functions
$\alpha,\beta$ which vanish for $s\le0$, are positive for $s>0$, are flat
at zero, and satisfy $\alpha'(s)>0$ for $s>0$.  On $W$ take $\alpha=0$ and
the horizontal vector field to be zero.  For small $c>0$ put
\[
 N_{\mathrm{sm}}=\{\alpha(s)+t^2\le c\},\qquad
 V=\beta(s)\partial_s+t\partial_t,
 \tag{C5H.8}
\]
with the first summand zero on $W$.  Choose the collar long enough that
$\alpha$ exceeds $c$ before its outer end.  The boundary $L_{\mathrm{sm}}$
is smooth; the derivative of $\alpha(s)+t^2$ along $V$ is
$\alpha'(s)\beta(s)+2t^2$, positive except on $W\times\{0\}$.
Both halves of $L_{\mathrm{sm}}$ are diffeomorphic to $W$, so it is
$DW\cong S^5$.  The function $h=t|L_{\mathrm{sm}}$ is smooth and has zero
as a regular value: at $t=0$, $\partial_t$ is tangent to $L_{\mathrm{sm}}$
and $dh(\partial_t)=1$.

Every forward orbit from $N_{\mathrm{sm}}\setminus(W\times\{0\})$ exits
through $L_{\mathrm{sm}}$ in finite time.  For $t\ne0$ this follows from
$V(t)=t$; for $t=0,s>0$ it follows from $\beta(s)>0$.
Backward orbits converge uniformly towards the compact set
$W\times\{0\}$: $t$ decreases exponentially and the exterior collar
coordinate decreases monotonically to zero.  The flow consequently gives
a diffeomorphism
\[
 L_{\mathrm{sm}}\times(0,1]\longrightarrow
 N_{\mathrm{sm}}\setminus(W\times\{0\}),\qquad
 (\theta,r)\longmapsto\Phi_{\log r}(\theta),
\]
and after collapse a homeomorphism $C(L_{\mathrm{sm}})\cong
N_{\mathrm{sm}}/(W\times\{0\})$.  It is the identity on the boundary and
satisfies the exact formula $p(r,\theta)=rh(\theta)$.
Using a diffeomorphism $L_{\mathrm{sm}}\cong S^5$ gives a smooth disc
structure at the cone point, compatible off that point with the original
smooth structure through the flow chart.  The flow extends across the
boundary, so these charts glue smoothly to the exterior.

Choose a strictly increasing smooth $g:[0,1]\to[0,1]$ with $g(0)=0$,
$g(r)=\exp(-1/r^2)$ near zero, $g(r)=r$ near one, and $g'(r)>0$ for $r>0$.
The radial homeomorphism $(r,\theta)\mapsto(g(r),\theta)$ is the identity
near the boundary.  Precompose $p$ with it.  The resulting local function
is
\[
 p_{\mathrm{sm}}(r,\theta)=g(r)h(\theta).
 \tag{C5H.9}
\]
In ordinary Euclidean cone coordinates this is smooth and flat at the
origin: derivatives of the angular factor have at most power growth in
$r^{-1}$ and are dominated by the flat radial factor.  Off the origin it
is a submersion: use $g'(r)h(\theta)\ne0$ when $h\ne0$, and $dh\ne0$ when
$h=0$.  The domain homeomorphism preserves the Hurewicz property and
properness, and the exceptional fibre still has local model $C(H)$.
This proves the smooth assertion of the counterexample.

\paragraph{Insertion into the two-stratum manuscript axioms.}
For either of the two categories, put
\[
 \Sigma_0=S^1,\qquad \Sigma_1=E\times[0,1],\qquad
 \Sigma_{01}=E\times\{0,1\}.
\]
The target is the whole-boundary embedding
$\Sigma_{01}=\partial\Sigma_1\hookrightarrow\Sigma_1$.
The source is $p$ on each of its two components.  All spaces are compact
manifolds in the chosen category; the source is proper and Hurewicz;
$\partial\Sigma_0=\varnothing$; there are no strictly composable pairs or
higher faces.  Thus boundary-transversality and triple compatibility
conditions impose no extra restriction.  The source fibre over $0$ is
$Y\amalg Y$, which is not a manifold.  This is the requested counterexample
to deriving the fibre-manifold assertion on 150/051 and 150/063 from mere
proper Hurewicz fibrancy and object-manifold hypotheses.

\paragraph{Auxiliary finite certificate.}
Rolfsen's printed p.~356, \S11E, Exercise~3, records the boundary presentation\cite{C5HRolfsen1976}
\[
 \langle x,y\mid y^5=x^7,\ y^4=x^2yx^2\rangle.
\]
It has a concrete nontrivial permutation quotient:
$x=(0\ 1\ 2\ 3\ 4\ 5\ 6)$ and $y=(0\ 3\ 1\ 4\ 2)$.
The script \texttt{code\_B56/check\_C5\_boundary\_group.py} checks both
relations, examines all 504 five-cycles for this $x$, finds 14 solutions,
and verifies that this pair generates a group of order 2520, namely $A_7$.
The recorded run took 0.006783651 seconds.  The geometric proof uses the
opened Zeeman construction; this additional finite check does not substitute
for a derivation of a handle presentation.


\subsection{C5.07 --- Refuted in TOP: connected components suffice}\label{c5.07-refuted-in-top-connected-components-suffice}

The word ``connected'' in /068--071 cannot be replaced silently by ``path-connected.'' Even all the stated finiteness, properness, local-immersion, and collared-target conditions permit a strict arrow within one connected component. The following example also satisfies the stronger bordant-target condition of /063, not just the nowhere-dense condition of /069.

Let \(K\subset\mathbb R^2\) be a compact, nondegenerate hereditarily indecomposable continuum, for example the plane pseudoarc constructed by R. H. Bing, \emph{Concerning hereditarily indecomposable continua}, Example 1, printed pp.43--44 \cite{C_Bing1951}. Every path in \(K\) is constant. Indeed, a nonconstant path has a nondegenerate image continuum. Uniform continuity partitions its domain into finitely many consecutive intervals whose image continua all have diameter smaller than that whole image. The image would then be a finite union of proper continua with connected successive unions. Indecomposability rules this out, by induction on the number of pieces.

Choose a continuous map \(q:[1,\infty)\to\mathbb R^2\) whose cluster set at infinity is exactly \(K\), with \(\operatorname{dist}(q(t),K)\to0\). Here is an explicit existence construction. Fix a point of \(K\). For each positive integer \(n\), choose a finite \(1/n\)-net in \(K\). Its open \(1/n\)-neighbourhood in the plane is connected and hence path-connected. On \([n,n+1]\), let \(q\) traverse a piecewise linear loop, based at the fixed point, which visits every point of that net and stays in that neighbourhood. Concatenation is continuous. These visits show that every point of \(K\) is a cluster point; the neighbourhood bound excludes any other cluster point.

Set \[
 C=\{(q(t),1/t):t\ge1\}\subset\mathbb R^3,
 \qquad P=C\cup(K\times\{0\}).
\] The height coordinate identifies \(C\) with the ray \([1,\infty)\). The space \(P\) is its compact closure, hence is a connected compact metric space of finite covering dimension at most three. The dimension bound can be checked directly: for any finite open cover of this compact subset of \(\mathbb R^3\), choose a sufficiently fine triangulation of \(\mathbb R^3\). Its vertex open-star cover restricts to a refinement on \(P\), by the Lebesgue number of the given cover; each point lies in at most four such stars, the vertices of its carrier simplex. This is the defining covering-dimension bound. Its single connected component is open. Write \(b=(q(1),1)\); its neighbourhood with height greater than \(1/(1+\epsilon)\) is a collar \([0,\epsilon)\) of the point \(b\). Choose \(a\in K\times\{0\}\).

The path component of \(a\) in \(P\) is the singleton \(\{a\}\). To verify this, suppose a path starting in the remainder enters \(C\). An interval component of the inverse image of \(C\) has a finite left endpoint at which the height is zero. On that interval the corresponding ray parameter tends to infinity at its left endpoint. By the intermediate value theorem, the image of every sufficiently short interval after that endpoint contains an entire tail \([T,\infty)\) of ray parameters. Continuity of the path would therefore force the whole tail of \(q\) to converge to the same remainder point. This contradicts its nondegenerate cluster set \(K\). A path remaining in the remainder is constant by the preceding indecomposability argument.

Give \(P\) the order with relation \[
 R=\Delta_P\cup\{(a,b)\}\subset P\times P.
\] This is a closed partial order. Its strict relation consists of one point, which is a clopen component of \(R\); the other component is \(\Delta_P\cong P\). Every lower interval is totally ordered and has at most two elements. Both projections \(s,b:R\to P\) are proper because the domain is compact. The target projection is a local embedding: it is the identity on the diagonal component and the embedding of a point on the strict component. The strict target image \(\{b\}\) is nowhere dense and collared. Thus its family of incoming boundary germs is transverse; there is at most one such germ at any point.

The source projection is a Hurewicz fibration, including its full lifting property for arbitrary parameter spaces. Given a map \(g:Z\to R\) and a homotopy \(H:Z\times[0,1]\to P\) with \(H(z,0)=s(g(z))\), let \(Z_e\) be the inverse image of the isolated strict arrow. This is clopen in \(Z\). On \(Z_e\), each path \(H(z,-)\) starts at \(a\), so it is constant; lift it constantly to the strict arrow. On \(Z\setminus Z_e\), lift by \((H(z,t),H(z,t))\in\Delta_P\). These continuous lifts agree with the initial map and glue across the clopen partition. Source transversality to the target collar is also satisfied: on the identity component it is the identity, while the strict component maps to \(a\), outside a chosen collar of \(b\).

Nevertheless \(a<b\) lie in the same connected component of \(P\). The map to the one-point set of connected components is not strictly increasing. This refutes the unrestricted TOP conclusion of /070 and therefore the equivalence claimed in /071 in that scope. The additional monomorphism condition for distinct components is vacuous here and cannot repair it. No doubtful transcription is used: the original /070--071 explicitly says connected components, as does the open-component assumption on /068.

The exact correction proved in C5.01 is to use open, path-connected components; in particular it applies to manifolds and locally path-connected PL spaces. Another sufficient correction, which can use connected components without local path-connectedness, is to require the proper source map to be open. Here is the additional proof. Suppose a strict arrow has source and target in the same open connected component \(D\). The strict arrows whose targets lie in \(D\) form a clopen subspace of \(R\). Their source image is open by openness of \(s\), and closed by properness; it meets \(D\), so it contains all of \(D\). Thus every point of \(D\) has a strict successor in \(D\). Starting at one point gives arbitrarily long strict chains in its compact upper interval. Proper local injectivity of the target projection makes the cardinalities of lower intervals locally bounded, hence bounded on that compact upper interval, a contradiction. The same open-and-closed image argument aligns middle objects to prove transitivity of the component relation. Cycles would give a strict chain returning to a component and then the same contradiction. Total ordering of the pointwise lower intervals gives the proper-injective target comparison of C5.01. A locally trivial source bundle is open, so this correction applies to the strong equisingular interpretation too.

\subsection{C5.08 --- Proved and refuted: the matrix and category digression}\label{c5.08-proved-and-refuted-the-matrix-and-category-digression}

For the matrix below, \textbf{raw fibrancy} means a proper Hurewicz fibration, even when its underlying map happens to be PL or smooth. A \textbf{regular PL source} is a locally trivial PL bundle; a \textbf{regular smooth source system} has proper submersions on every source and compatible horizontal product collars, together with full projection control; C5.05 proves the necessary regularity on all vertical faces. A transversality assertion refers to the manuscript's product-collar condition, not merely to equality of inverse-image subsets. These distinctions are necessary to evaluate R7.

\begin{longtable}[]{@{}
  >{\raggedright\arraybackslash}p{(\columnwidth - 6\tabcolsep) * \real{0.2500}}
  >{\raggedright\arraybackslash}p{(\columnwidth - 6\tabcolsep) * \real{0.2500}}
  >{\raggedright\arraybackslash}p{(\columnwidth - 6\tabcolsep) * \real{0.2500}}
  >{\raggedright\arraybackslash}p{(\columnwidth - 6\tabcolsep) * \real{0.2500}}@{}}
\toprule\noalign{}
\begin{minipage}[b]{\linewidth}\raggedright
Manuscript margin
\end{minipage} & \begin{minipage}[b]{\linewidth}\raggedright
TOP
\end{minipage} & \begin{minipage}[b]{\linewidth}\raggedright
PL
\end{minipage} & \begin{minipage}[b]{\linewidth}\raggedright
Smooth
\end{minipage} \\
\midrule\noalign{}
\endhead
\bottomrule\noalign{}
\endlastfoot
/051(d): incoming flag faces are manifolds from the ambient object-manifold condition & False: C5.04 & True for PL collars: C5.05 & True for smooth collars: C5.05 \\
/051(e) and /063: mixed fibre-manifold and fibre-boundary condition follows from nonsingularity and raw fibrancy & False even for locally trivial TOP bundles: C5.03 & False: the proper PL Hurewicz example C5.06 & False: the proper smooth Hurewicz example C5.06 \\
Same mixed implication with regular PL/smooth source & A TOP bundle still does not suffice & True: C5.05 & True: C5.05 \\
/062: source transversality follows from the other bare order/face axioms & False: C5.02 & False: C5.02 & False: C5.02 \\
/062 with local product triviality of the full face diagram & True by the product charts & True by the product charts & True by the product charts \\
/067: incoming faces and their stated boundaries are manifolds from object nonsingularity & False: C5.04 & True: C5.05 & True: C5.05 \\
/070--071 for ordered spaces with open path-connected components & True: C5.01 & True: C5.01 & True: C5.01 \\
/070--071 for arbitrary open finite-dimensional connected TOP components & False even with collared strict target: C5.07 & Locally finite PL spaces are locally path-connected, so C5.01 applies & Manifolds are locally path-connected, so C5.01 applies \\
\end{longtable}

The proposed boundary formulas in the positive rows are part of C5.05, not extra assumptions about the answer: horizontal faces are the inverse images of lower faces of the base, and vertical faces are exactly those inserting an intermediate label. Their union is the boundary of the arrow/flag space; only the vertical part is the intrinsic boundary of the source fibre. Empty fibres cause no exception. These arguments are local in every dimension, including zero, one, and two. The counterexamples in larger dimensions refute the unrestricted claims without purporting to be minimal-dimensional ones.

The general categorical paragraph /067--070 can also be settled for manifold object components, without assuming the category is thin. Let a topological category have Hausdorff object and arrow spaces, open manifold components of its object space, proper locally embedding target map \(b\), Hurewicz source map \(s\), and locally nowhere-dense target image of the strict arrows. Assume every invertible arrow is an identity. Then the induced relation on object components is a partial order, and every nonidentity arrow connects distinct components.

Here is a full proof which avoids the manuscript's dimension inequalities. The identity section is closed by Hausdorffness and open by local injectivity of \(b\): near an identity, both that section and a sufficiently small arrow neighbourhood lie in a target-injective chart, so the only arrow in that chart with a given nearby target is its identity. Hence strict arrows form a clopen subspace. Source lifting along a path cannot change from a strict arrow to an identity, and target continuity keeps its target in one component. It follows that the existence of an arrow from one component to another is independent of which source point in the first component is prescribed. This also proves transitivity by lifting the second arrow along a path joining the intermediate endpoints.

Suppose a strict arrow has source and target in the same connected \(n\)-manifold component \(D\), allowing boundary or corners. Move its source by path lifting into the interior of \(D\); connected manifold components are path-connected and have nonempty interior, with the dimension-zero case interpreted directly. The target remains in \(D\). Restrict to strict arrows whose source and target lie in \(D\). Over a sufficiently small open coordinate ball \(U\subset D\), a contraction of \(U\) to a point and the Hurewicz lifting property give a continuous section \(\sigma:U\to R^{<}\) of \(s\). Starting from a constant arrow lift proves this for arbitrary parameter space \(U\); the lift stays strict and its target stays in \(D\). Shrink \(U\) around a point so that \(\sigma(U)\) is contained in one target-injective neighbourhood. Then \(b\sigma:U\to D\) is injective, because \(\sigma\) is a section and \(b\) is injective on that neighbourhood. Choose a target corner chart and include its orthant in \(\mathbb R^n\). Invariance of domain makes the image of \(b\sigma\) open in \(\mathbb R^n\); it cannot meet the boundary of that orthant, and hence is open in \(D\). This contradicts the locally nowhere-dense strict target image. Thus no strict arrow has both endpoints in one component. This argument includes dimension zero directly: a nonempty image in a discrete component is open.

If there are arrows from \(D\) to \(E\) and back, path lifting aligns their middle endpoints. Write the resulting arrows as \(u:d\to e\) and \(v:e\to d'\), with \(d,d'\in D\). The composite \(vu\) must be an identity by the preceding paragraph, so \(d'=d\). The composite \(uv\), with endpoints in \(E\), must likewise be an identity. Thus \(u\) and \(v\) are inverse; the assumption on invertible arrows forces \(D=E\). This proves antisymmetry. If the category is an ordered space with totally ordered lower intervals, the extra componentwise monomorphism statement and incomparable-face disjointness follow exactly as in C5.01; no extra monomorphism axiom is needed.

The same categorical proof applies to finite-dimensional locally finite PL object spaces, even if they are polyhedra rather than manifolds, provided the arrow charts are PL and the target map is a PL local embedding. These are the category-specific assumptions used in the PL-image argument below. For an alleged strict arrow within a component of maximal dimension \(n\), use source path lifting to move its source into the interior of a maximal \(n\)-simplex; this is an open \(n\)-ball. A local source section embeds that ball into a target-injective strict arrow chart. The locally rare PL target image contains no \(n\)-simplex and has local dimension at most \(n-1\). For completeness this cannot contain an embedded \(n\)-ball: a compact sub-ball meets only finitely many target simplices; the closed preimages of these simplices cover it, so one has interior by the Baire theorem. That gives an injective map from an open \(n\)-ball into a Euclidean space of dimension at most \(n-1\), contradicting invariance of domain after inclusion into \(\mathbb R^n\). The remaining component-order argument is unchanged. The TOP counterexample C5.07 explains why arbitrary connected finite-dimensional components do not admit this argument.

\textbf{Doubtful and ill-typed readings.} The original /051, /062, /063, /067, /070, and /071 scans have been opened visually \cite{C_Grothendieck150}. In /063(d) the transcription flags the source letter as \(s\) or \(b\); the source reading is the one explicitly typed as \(p_1\) in the preceding cartesian square on /062 and used in /065--066. All nontrivial statements above use that reading. If one literally requires the target boundary embedding \(b\) to be product-collar transverse to itself, a nonempty strict arrow chart is impossible: its inverse image of that boundary is its whole domain, which cannot itself be a nonempty collared boundary subset of that domain. Under that alternative the strict relation is empty and the nonsingular object's boundary is empty, so the corresponding mixed assertions are vacuous. The uncertain indices in /067's displayed source are typed as \(s_{jk}:\Sigma_{jk}\to\Sigma_j\); every other reading with a different source label must be reindexed to this type or is not a composable map. The uncertainty in /051(d) between deducing and separately requiring manifold flag spaces does not affect the mixed counterexamples: their arrow spaces are themselves manifolds. Finally the dimension and connected-component wording on /070--071 is clear in the original; C5.07 does not depend on an uncertain reading.

\part{Verification and scope}
\section{Completed questions and retained qualifications}
All items A.1--A.4, B.1--B.6, C.1--C.5 and leads L1--L4 are
settled by the numbered statements above. A refutation includes both a
counterexample and a proved replacement. The outcome table in the
accompanying report maps every requested item to these statements.
There is no remaining mathematical proof obligation asserted as a theorem.

The quantifiers in those statements are part of their conclusions.
In particular, the ordinary bicategorical descent theorem uses degrees
zero through two; the higher shape theorem uses the full simplicial
diagram. Weak conical data, Hurewicz fibrations, locally trivial bundles,
and controlled smooth submersions are not interchangeable. Moduli
incidence categories retain automorphisms, and the coarse-space
manifold assertion has an explicit counterexample.

\section{Manuscript uncertainty and source status}
Unresolved ink is not silently resolved by a mathematical preference.
For the final subscript of 151/005 we prove completeness for the
subscript-one reading and refute the subscript-three reading. The
overwritten index and auxiliary cutoff on 151/060 do not rescue the
gluing claim under either reading treated in A1.3. Other material
reading qualifications are stated where they enter a proof.

Every bibliography entry records whether it was opened. Sources that
could not be opened are recorded in the research notes and are not
used as unchecked proof premises. In particular, the original Mazur
1961 article was not opened; the construction in C5.06 uses the
explicitly opened Zeeman, Milnor, and Ferry sources and the intervening
arguments supplied here.

\section{Certificates and attribution}
The results register has one row per numbered statement. The check-event
table distinguishes direct proof checks, independently repeated
arguments, finite computations, and source checks. Finite computations
verify only their stated finite ranges. General theorems depend on
the written proofs and cited opened results, not on extrapolation from
those ranges. The accompanying program catalog records inputs,
outputs, counts, and observed runtimes.

This research package is AI-produced by GPT-6 Astra (OpenAI, ChatGPT Pro
mode). The supplied study is identified as an input, not adopted as
authority. We make no novelty or priority claim and attribute none of
this run's writing or calculations to a human author.

\bibliographystyle{plain}
\bibliography{RESEARCH_SA10_references}
\end{document}
